% concatenated from main_template.tex
\documentclass[11pt]{article}
% arXiv version using the user-provided main_template.tex / nexais.sty template.
% The scientific manuscript body is stored byte-for-byte in manuscript_body_exact.tex.

% Encoding required by the original source text.
\usepackage[T1]{fontenc}
\usepackage[utf8]{inputenc}
\DeclareUnicodeCharacter{2011}{-}

%%% Packages from the supplied arXiv template
\usepackage{amsmath,amssymb}
\usepackage{times}
\usepackage{bm}
\usepackage{natbib}
\usepackage{algorithm}
\usepackage{mathrsfs}
\usepackage{graphicx}
\usepackage[flushleft]{threeparttable}
\usepackage[shortlabels]{enumitem}
\usepackage{titletoc}
\usepackage{tocloft}
\usepackage{nexais}
% Put each affiliation on a separate line
%\makeatletter
%\renewcommand\affiliation[2][]{%
%  \ifdefempty{\affiliationlist}{%
%    \gappto\affiliationlist{%
%      {\normalsize \nexais@marker{#1}#2}%
%    }%
%  }{%
%    \gappto\affiliationlist{%
%      \\[2pt]{\normalsize \nexais@marker{#1}#2}%
%    }%
%  }%
%}
\makeatother
%%% Additional packages required by the unchanged manuscript body
\usepackage{microtype}
\usepackage{cmap}
\usepackage{comment}
\usepackage{booktabs,float}
\usepackage{rotating}
\usepackage{tikz}
\usepackage{array}
\usepackage{tabularx}
\usepackage{mathtools}
\usepackage[most]{tcolorbox}
\usepackage{marginnote}

\usepackage[colorlinks,
linkcolor=metablue,
anchorcolor=metablue,
citecolor=metablue,
urlcolor=metablue
]{hyperref}

\usepackage{txfonts}
\usepackage{geometry}
\geometry{
 a4paper,
 top=26mm,
 bottom=26mm,
 left=26mm,
 right=26mm}
\renewcommand{\baselinestretch}{1.10}

% Original manuscript macros retained for source compatibility.
\providecommand{\dataset}{{\cal D}}
\providecommand{\fracpartial}[2]{\frac{\partial #1}{\partial #2}}
\providecommand{\red}{\color{red}}
\providecommand{\magenta}{\color{magenta}}
\providecommand{\black}{\color{black}}
\providecommand{\F}{\mathcal{F}}
\providecommand{\huiming}{\color{blue}}
\providecommand{\qiuran}{\color{red}}
\providecommand{\lihu}{\color{magenta}}
\providecommand{\e}{\varepsilon}
\providecommand{\normal}{\color{black}}
\providecommand{\sce}{\sum_{k\geq \lceil \frac{2}{\theta} \rceil}}
\providecommand{\I}{\frac{1}{\theta}}
\providecommand{\lk}{\left (}
\providecommand{\rk}{\right )}
\providecommand{\LK}{\left [}
\providecommand{\RK}{\right ]}
\providecommand{\DD}{\mathcal{D}}
\providecommand{\pp}{\mathcal{P}}
\providecommand{\gen}{\operatorname{gen}}
\providecommand{\al}{\alpha}
\providecommand{\Ia}{\frac{1}{\alpha-1}}
\providecommand{\Unif}{\mathrm{U}}
\providecommand{\X}{\mathcal{X}}
\providecommand{\Y}{\mathcal{Y}}
\providecommand{\DKL}[2]{D_{\mathrm{KL}}\left( #1 \,||\, #2 \right)}
\providecommand{\Z}{\mathcal{Z}}
\providecommand{\paren}[1]{\left(#1\right)}
\renewcommand{\norm}[1]{\left\|#1\right\|}
\renewcommand{\abs}[1]{\left|#1\right|}
\providecommand{\bracks}[1]{\left[#1\right]}
\renewcommand{\paragraph}[1]{\noindent\textit{#1}}
\providecommand{\Renyi}{R\'enyi }
\providecommand{\subw}{\operatorname{subW}}
\providecommand{\KL}{\text{KL}}
\providecommand{\iid}{\textrm{i.i.d.}}
\DeclareMathOperator{\clip}{clip}

% Original comment macros; submission comments remain disabled as in the source.
\providecommand{\cmmnt}[1]{}
\providecommand{\scolor}[1]{{\color{magenta}#1}}
\providecommand{\rcolor}[1]{{\color{red}#1}}
\providecommand{\scomment}[1]{}
\providecommand{\sq}{\scomment{Current point.}}
\providecommand{\checkthis}{\scomment{Check this.}}
\marginparwidth60pt
\renewcommand*{\marginfont}{\color{magenta}\upshape\sffamily\tiny}

\begin{document}

\title{Tail-Aware Information-Theoretic Bounds for LLM Alignment under Heavy-Tailed Rewards}

\author[1,*]{Huiming Zhang}
\author[2,*]{Binghan Li}
\author[3]{Wan Tian}
\author[4]{Qiang Sun}

\contribution[*]{Co-first authors.}

\affiliation[1]{BUAA and BAIC-FBPC}

\affiliation[2]{BUAA}

\affiliation[3]{PKU}

\affiliation[4]{UofT and MBZUAI}
%\affiliation[1]{Institute of Artificial Intelligence, Beihang University; Beijing Advanced Innovation Center for Future Blockchain and Privacy Computing}
%
%\affiliation[2]{Institute of Artificial Intelligence, Beihang University}
%
%\affiliation[3]{Advanced Institute of Information Technology, Peking University; Wangxuan Institute of Computer Technology, Peking University}
%
%\affiliation[4]{Computer Science and Statistics, University of Toronto and MBZUAI}

%\contribution[*]{}

\abstract{

Classical information-theoretic learning bounds typically rely on KL mutual
information and moment-generating-function (MGF) arguments, which are well
matched to bounded or sub-Gaussian losses but can be ineffective when losses or
rewards are heavy-tailed. We develop a tail-aware information-theoretic
framework for sub-Weibull data, where the tail parameter $\theta$ controls the
tail heaviness: $\theta=2$ corresponds to sub-Gaussian, $\theta=1$ to
sub-exponential, and $0<\theta<1$ to genuinely heavy tails. Our key technical
ingredient is a decorrelation lemma that bounds change-of-measure expectations
using a shifted-log $f_\theta$-divergence, which admits explicit comparisons to
R\'enyi divergence without MGF arguments. On the empirical-process side, we
establish sharp maximal inequalities and a Dudley-type chaining bound for
sub-Weibull processes, with logarithmic and entropy terms raised to the power
$1/\theta$. These tools yield tail-adaptive selection bounds and a multiscale
information-theoretic Dudley inequality based on shifted-log and R\'enyi mutual
information. We apply our theory to large language models (LLMs) in the context of reward hacking within reinforcement learning from human feedback (RLHF). We show that Rényi-regularized alignment provides finite reward guarantees and ensures that best-of-N policies remain well-controlled, thereby mitigating the catastrophic Goodhart effects where standard KL-regularization fails. We illustrate R\'enyi-regularized RLHF by experiments, including controlled heavy-tailed rewards and token-space reward attacks.  The code is available at \url{https://github.com/NeXAIS/SubW-Alignment}.

}
\metadata[Keywords]{
heavy-tailed distributions, sub-Weibull processes, R\'enyi divergence, information-theoretic generalization bounds, LLM, RLHF.
}

\metadata[Code]{\url{https://github.com/NeXAIS/SubW-Alignment}}

\date{\today}

\maketitle

\bibliographystyle{apalike}






\tableofcontents




%\newpage
\section{Introduction}\label{intro}

% Generalization, the ability of a model trained on finite data to perform well
% on unseen samples, is a central concern in statistical learning.  For a
% hypothesis $w\in\mathcal W$ and loss $\ell(w,z)$, the (centered) generalization
% gap compares population and empirical risks,
% \[
% \gen(w,S)=\E_\mu[\ell(w,Z)]-\frac{1}{n}\sum_{i=1}^n\ell(w,Z_i).
% \]
% Classical analyses often control
% $\sup_{w\in\mathcal W}|\gen(w,S)|$ by uniform concentration and complexity
% measures such as VC dimension, covering numbers, or Rademacher complexity.
% While powerful, uniform bounds can be overly pessimistic for modern
% \emph{data-adaptive} algorithms, which typically explore a sample-dependent
% subset of $\mathcal W$ and may generalize well even when worst-case complexity
% is large.


% This is the
% case for bandit and reinforcement learning methods
% \citep{russo2016information, lu2019information}, stochastic gradient descent
% \citep{algo1}, neural networks \citep{asadi2020chaining}, semi-supervised
% learning \citep{he2022information}, Gibbs/posterior-sampling algorithms
% \citep{aminian2021}, and preference-optimization methods \citep{mroueh2025information}. 
% A representative example is
% the minimal-norm interpolant in over-parameterized models, which can still
% generalize effectively despite the large complexity of the full hypothesis
% class \citep{zhou2021optimistic}.


%\paragraph{Background}
The generalization gap, namely the difference between population risk and
empirical risk, has traditionally been analyzed by uniform-convergence
arguments based on classical complexity measures such as VC dimension and
Rademacher or Gaussian complexity; see
\cite{koltchinskii2011oracle, Wainwright19, shalev2014understanding}.
While powerful, these tools are inherently worst-case: they control an entire
hypothesis class rather than the smaller, sample-dependent region actually
explored by a modern learning algorithm. As a result, they can be overly
pessimistic for data-adaptive procedures, which often generalize well even when
the ambient class has large worst-case complexity \citep{zhou2021optimistic}.



Information-theoretic generalization bounds replace global complexity by an
\emph{algorithm-dependent} quantity: how much information the algorithm's
output reveals about the sample \citep{xu2017,Russo20,bu2020}. In
light-tailed settings, under sub-Gaussian assumptions, KL-based
mutual information yields clean bounds in expectation and with high
probability, and these guarantees can be further sharpened through multiscale
or chaining arguments \citep{asadi2018chaining,hellstrom2025generalization}.




A major obstacle is that KL-based arguments are tightly coupled to moment
generating functions (MGFs). Many modern pipelines, however, exhibit
heavy-tailed behavior, so the relevant MGFs may fail to exist. This issue
arises in reward-model scores for RLHF \citep{kwa2024catastrophic}, gradient noise in large-scale
optimization \citep{raj2023algorithmic}, and losses induced by heavy-tailed data
\citep{lerasle2019lecture, xu2023non}. In such settings, KL mutual information
can remain small even when rare but extreme events dominate generalization or
stability; see Example~\ref{ex:weibull}.


% \begin{definition}[Sub-Weibull random variable]\label{subw}
% A random variable $X$ is called sub-Weibull with parameter $\theta>0$, denoted
% $X\sim\subw(\theta)$, if its Orlicz norm $\|X\|_{\psi_\theta}$ is finite, where
% $\psi_\theta(x)=\exp(x^\theta)-1$. Equivalently, there exists $K>0$ such that
% $\mathbb{E}\exp(|X|^\theta/K^\theta)\le 2$.
% \end{definition}


A convenient tail model that interpolates between sub-Gaussian and heavy-tailed
behavior is the \emph{sub-Weibull} family. 
A random variable $X$ is called sub-Weibull with parameter $\theta>0$, denoted
$X\sim\subw(\theta)$, if its Orlicz norm $\|X\|_{\psi_\theta}$ is finite, where
$\psi_\theta(x)=\exp(x^\theta)-1$. Equivalently, there exists $K>0$ such that
$\mathbb{E}\exp(|X|^\theta/K^\theta)\le 2$ or the tail estimate
\begin{equation}\label{tail-estimate}
\mathbb{P}(|X|\ge t)\le 2\exp\bigl(-t^\theta/K^\theta\bigr),\qquad t\ge 0,
\end{equation}
for some $K>0$. The $\theta$ controls tail thickness: $\theta=2$
corresponds to sub-Gaussian tails, $\theta=1$ to sub-exponential tails, and
$0<\theta<1$ to genuinely heavy-tailed regimes.  Our focus is on deriving
information-theoretic generalization tools that remain meaningful throughout the
sub-Weibull family, including $0<\theta<1$ where MGF-based techniques break down.




This work develops an information-theoretic generalization framework for
modern learning algorithms under heavy-tailed distributions and rewards. Our main contributions are:
\begin{enumerate}
\item \textbf{Tail-adaptive decorrelation via shifted-log divergences.}
We introduce a shifted-log $f_\theta$-divergence
$f_\theta(x)=x\log^{1/\theta}(x+A)$ and prove a decorrelation (change-of-measure)
lemma that bounds expectations under a dependent coupling through this
tail-adaptive divergence. We further show that it admits explicit upper bounds
in terms of R\'enyi divergence, leading to usable bounds in terms of R\'enyi
mutual information (Section~\ref{sec:pre}).

\item \textbf{Maximal and chaining bounds for sub-Weibull processes.}
For $0<\theta<1$, we prove a sharp maximal inequality for a finite family of
sub-Weibull random variables and a Dudley-type entropy integral for
$(\theta,C)$-sub-Weibull processes (Section~\ref{sec:ep}).

\item \textbf{Information-theoretic selection and chaining under sub-Weibull tails.}
Combining the decorrelation lemma with the new sub-Weibull maximal and chaining
bounds, we prove tail-adaptive information bounds for data-dependent selectors.
These inequalities replace KL mutual information, which can fail under heavy
tails, by shifted-log $f_\theta$-mutual information and its R\'enyi refinements,
and extend from finite selections to Dudley-type multiscale bounds
(Section~\ref{sec:gen_bounds}).

\item \textbf{Applications to RLHF under heavy-tailed and adversarial rewards.}
We show that KL-constrained RLHF can suffer catastrophic Goodhart under
heavy-tailed proxy rewards, while R\'enyi-regularized RLHF admits tail-adaptive
reward guarantees at a fixed divergence budget. We also analyze best-of-$n$
policies and bound their reward growth under sub-Weibull reward tails
(Section~\ref{sec:rlhf}). Numerically, we study both controlled heavy-tail
perturbations and token-space reward attacks; in both settings, higher-order
R\'enyi regularization dampens tail inflation and improves proxy--gold reward
alignment relative to KL (Sections~\ref{sec:rlhf-experiments}
and~\ref{sec:tompa-experiments}; Appendix~\ref{app:add_experiments}).

% \item \textbf{Applications to stochastic optimization.}
% We apply the framework to stochastic gradient Langevin dynamics (SGLD) under
% heavy-tailed stochastic-gradient noise, yielding stability-based generalization
% bounds that expose the interaction between step sizes, injected noise, and tail
% heaviness (Appendix~\ref{se:SGDL}).
\end{enumerate}



\subsection*{Related Work}

\paragraph{Maximal inequalities for heavy tails.}
Maximal inequalities and chaining are classical tools in empirical process
theory \citep{Talagrand2005,vanhandle2016probability}.
%A comprehensive
%treatment of concentration for sub-Weibull random variables is given by
%\citet{zhang2022sharper}.
While maximal inequalities for sub-Weibull processes are well-understood for
$\theta\ge 1$ \citep{kuchibhotla2022moving}, the heavy-tailed regime $0<\theta<1$ is more
delicate because MGFs can be infinite and many standard
arguments fail.  Recent works provide maximal inequalities with explicit
prefactors \citep{mendes2023concentration} and sharp Orlicz norm bounds
\citep{zhang2022sharper}. We complement these results by giving a direct
expectation bound for maxima (avoiding an additional norm-to-moment factor) and
by extending Dudley-type chaining bounds to heavy-tailed sub-Weibull processes.



\paragraph{Information-theoretic generalization and RL.}
Information-theoretic bounds date back to \citet{Russo20} and
\citet{xu2017}. Variants based on $f$-divergences and related divergences
have been developed in \citep{jiao2017dependence,hellstrom2025generalization,bu2020}
and references therein.  Most existing information-theoretic bounds rely on
KL-based mutual information (or MGF-based change-of-measure arguments) and are
therefore tailored to bounded or sub-Gaussian losses. We extend this
toolkit to sub-Weibull tails by combining a tail-adaptive divergence and a
multiscale information-theoretic chaining analysis. \cite{behnamnia2025log} proposed a log-sum-exponential estimator that avoids the pitfalls of heavy-tailed rewards. \cite{sun2026curverl} introduced a prompt reweighting method for LLM reasoning based on a distribution‑aware utility functional.




\subsection*{Overview}
Section~\ref{sec:ep} develops the sub-Weibull maximal and Dudley-type
empirical-process tools. Section~\ref{sec:pre} introduces the shifted-log
divergence, the decorrelation lemma, and the corresponding R\'enyi comparison
bounds. Section~\ref{sec:gen_bounds} combines these ingredients to obtain
tail-adaptive selection bounds and a multiscale information-theoretic Dudley
inequality. Section~\ref{sec:rlhf} applies the framework to RLHF under
heavy-tailed rewards, including R\'enyi-regularized and best-of-$n$ policies.
Section~\ref{sec:experiments} provides numerical studies, including controlled
heavy-tail perturbations and token-space reward attacks. Additional proofs and
robustness checks are deferred to the appendix. Table~\ref{tab:roadmap}
summarizes the active results.

% \subsection*{Roadmap of results}
% Table~\ref{tab:roadmap} summarizes the main results.

\begin{table}[t]
\centering
\footnotesize
\setlength{\tabcolsep}{4pt}
\renewcommand{\arraystretch}{1.05}
\begin{tabular}{@{}>{\raggedright\arraybackslash}p{0.20\linewidth}
                >{\raggedright\arraybackslash}p{0.30\linewidth}
                >{\raggedright\arraybackslash}p{0.48\linewidth}@{}}
\toprule
\textbf{Object} & \textbf{Classical} & \textbf{This paper} \\
\midrule
Finite maxima &
$\E\max_i X_i\lesssim\sqrt{\log n}$ &
$\E\max_i X_i\lesssim(\log n)^{1/\theta}$ (Lemma~\ref{pp-MaximalSW}) \\
Process suprema &
Dudley: $\int\sqrt{\log N_\varepsilon}\,d\varepsilon$ &
Sub-Weibull Dudley: $\int(\log N_\varepsilon)^{1/\theta}\,d\varepsilon$ (Theorem~\ref{Dudley}) \\
Data-dependent selection &
KL mutual information $\sqrt{I(W;S)}$ &
Shifted-log $I_{f_\theta}(W;S)$ and R\'enyi refinements (Theorem~\ref{infob}) \\
Multiscale selection &
Sub-Gaussian information chaining &
Sub-Weibull chaining with $I_{f_\theta}([W]_k;S)$ and R\'enyi refinements (Theorem~\ref{chaining}) \\
RLHF trust regions &
KL trust regions can fail under heavy tails &
R\'enyi trust regions control reward inflation (Theorem~\ref{thm:rlhf}) \\
Best-of-$n$ alignment &
Light-tailed maximal control &
R\'enyi divergence gives $(\log n)^{1/\theta}$-type reward control (Theorem~\ref{renyiregret}) \\
\bottomrule
\end{tabular}
\caption{Roadmap of results.$\theta$ is the sub-Weibull tail parameter and $N_\varepsilon=N(T,d,\varepsilon)$.}
\label{tab:roadmap}
\end{table}




\subsection*{Notation}
For a positive integer $n$, we write $[n]=\{1,\ldots,n\}$. We use $\E$ for
expectation and $\mathbb{P}$ for probability. We denote by $\DD_{\mathrm{KL}}$
the KL divergence and by $\DD_\alpha$ the R\'enyi divergence of order $\alpha>1$.
We write $\|X\|_{\psi_\theta}$ for the Orlicz norm associated with
$\psi_\theta(x)=e^{x^\theta}-1$.  We use $\lesssim$ to hide universal positive
constants. Let $(T, d)$ be a metric space and let $K \subset T$.  A subset $\mathcal{N} \subset T$ is called an \textit{$\varepsilon$-net} for $K$ if, for every $x \in K$, 
there exists $y \in \mathcal{N}$ such that $d(x, y) \leq \varepsilon$. The {covering number}, denoted by $N(K, d, \varepsilon)$,  is the minimum cardinality of such an $\varepsilon$-net. Then $e(T):=\inf\{\varepsilon>0: N(T,d,\varepsilon)=1\}$ is the minimum covering radius of $T$. 
Let $\mathcal{X}$ be the range of a random variable $X$. A partition $\mathcal{P}$ of $\mathcal{X}$ is a finite collection of disjoint sets $P_i$ such that \(\bigcup_i P_i = \mathcal{X}\).


\section{Sub-Weibull Process Theory}\label{sec:ep}

This section develops empirical-process tools for controlling maxima and
suprema of sub-Weibull random variables and processes, with an emphasis on the
heavy-tailed regime $0<\theta<1$.  In this regime, MGFs may be infinite, so many classical MGF-based arguments must be replaced
by Orlicz-norm methods.  We establish (i) a sharp maximal inequality for a
finite family of sub-Weibull variables and (ii) a Dudley-type chaining bound
for $(\theta,C)$-sub-Weibull processes.  These are \emph{uniform} (worst-case)
bounds; Section~\ref{sec:gen_bounds} will refine them for data-adaptive
algorithms using information measures.

\subsection{A sharp maximal inequality}

\begin{lemma}[Maximal inequality for heavy-tailed sub-Weibull variables]\label{pp-MaximalSW}
Let $0<\theta<1$ and $\psi_\theta(x)=e^{x^\theta}-1$. Let $\{X_i\}_{i=1}^n$ be random variables with
\(
\max_{i\in [n]}\|X_i\|_{\psi_\theta}<\infty.
\)
Then
\begin{equation}\label{eq:MaximalSW}
\mathbb{E}\Bigl[\max_{i\in [n]}|X_i|\Bigr]
\;\le\;
\psi_\theta^{-1}\bigl(\psi_\theta(x_\theta)+n\bigr)\,
\max_{i\in [n]}\|X_i\|_{\psi_\theta}
\;=\;O\bigl((\log n)^{1/\theta}\bigr),
\end{equation}
where $
\psi_\theta^{-1}\bigl(\psi_\theta(x_\theta)+n\bigr)
=\bigl[\log\bigl(1+\psi_\theta(x_\theta)+n\bigr)\bigr]^{1/\theta}$ with $x_\theta:=\Bigl(\frac{1-\theta}{\theta}\Bigr)^{1/\theta}$.
\end{lemma}


The $1/\theta$ in~\eqref{eq:MaximalSW} is sharp: for i.i.d.
$\mathrm{Weibull}(\theta)$ variables,
$\mathbb{E}[\max_{i\in [n]}X_i]\asymp (\log n)^{1/\theta}$; see Appendix~\ref{app:sec2}.
Thus, heavier tails (smaller $\theta$) necessarily worsen the growth rate of the
maximum.  A useful heuristic is that sub-Weibull variables behave like power
transforms of unbounded sub-Gaussian variables, so decreasing $\theta$ makes
rare extremes more influential.

% \noindent\textit{Comparison with norm-to-moment bounds.}
A common route is to first control the Orlicz norm of the maximum and then
convert that Orlicz bound to an expectation bound via a norm-to-moment
inequality \citep{mendes2023concentration, zhang2022sharper}.  This often
introduces an additional factor $\Gamma(1/\theta+1)$, which grows rapidly as
$\theta\downarrow0$ and can noticeably loosen prefactors in the heavy-tailed
regime.  Lemma~\ref{pp-MaximalSW} bounds the expectation directly and avoids this
extra loss while retaining the optimal $O\bigl((\log n)^{1/\theta}\bigr)$
scaling; a quantitative comparison is given in Appendix~\ref{app:max_comparison}.




\subsection{From maxima to suprema: sub-Weibull processes}

Passing from finite maxima to suprema over a general index set raises
measurability issues: $\omega\mapsto\sup_{t\in T}X_t(\omega)$ need not be
measurable.  A standard way to avoid outer expectations is to assume
\emph{separability}, which reduces the supremum to one over a countable dense
subset \citep[Definition~5.22]{vanhandle2016probability}.

\begin{definition}[Separable process]\label{def:separable}
Let $(T,d)$ be a metric space. A stochastic process $\{X_t\}_{t\in T}$ is
\emph{separable} if there exists a countable set $T_0\subseteq T$ such that
\[
X_t \in \lim_{s \to t,~s \in T_0} X_s\;\;\text{for all } t \in T\;\text{a.s.}
\]
Equivalently, for every $t\in T$ there exists $\{s_k\}\subseteq T_0$ with
$s_k\to t$ and $X_{s_k}\to X_t$ a.s..
\end{definition}

Empirical process theory typically controls oscillations of a process via tail
bounds on increments relative to a metric.  In the sub-Gaussian case, an
increment condition of the form
$\|X_t-X_s\|_{\psi_2}\lesssim d(t,s)$ leads to Dudley's entropy bound
\(\E\sup_{t\in T}X_t\lesssim\int_0^{\infty}\sqrt{\log N(T,d,\varepsilon)}\,d\varepsilon\).
In heavy-tailed settings, MGFs may fail to exist, yet increments are still often
controlled by an exponential-of-a-power tail, motivating the following
definition.

\begin{definition}[Sub-Weibull process]\label{subWp}
Let $\theta>0$, $C>0$, and let $(T,d)$ be a metric space. A mean-zero process
$\{X_t\}_{t\in T}$ is a \emph{$(\theta,C)$-sub-Weibull process} (with respect to
$d$) if
\[
\mathbb{E}\exp\!\left(\frac{|X_t-X_s|^\theta}{[Cd(t,s)]^\theta}\right)\le 2,
\qquad \forall s,t\in T~\text{s.t.}~s\ne t.
\]
Equivalently, $\|X_t-X_s\|_{\psi_\theta}\le C\,d(t,s)$ for all $s,t\in T$.
\end{definition}

For $\theta\ge 1$, such increment control implies suitable log-MGF bounds and
classical chaining machinery applies.  When $0<\theta<1$, MGF-based arguments
generally break down; nevertheless, Orlicz increment control remains sufficient
to develop a Dudley-type entropy bound once we can control maxima at each scale
(Lemma~\ref{pp-MaximalSW}).

\begin{theorem}[Heavy-tailed Dudley inequality]\label{Dudley}
Let $(T,d)$ be a finite metric space and let $\{X_t\}_{t\in T}$ be a mean-zero
$(\theta,C)$-sub-Weibull process with $0<\theta<1$. Then
\begin{equation}\label{eq:heavy_dudley}
\mathbb{E}\Bigl[\sup_{t\in T}X_t\Bigr]
\le 4C K_\theta\int_0^\infty \bigl[\log N(T,d,\varepsilon)\bigr]^{1/\theta}\,d\varepsilon,
\end{equation}
where $N(T,d,\varepsilon)$ is the covering number and
\[
K_\theta
:=\sup_{x\ge 2}\left(\frac{\log(1+\psi_\theta(x_\theta)+x)}{\log x}\right)^{1/\theta}
<\infty,
\qquad x_\theta=\left(\frac{1-\theta}{\theta}\right)^{1/\theta}.
\]
\end{theorem}

Theorem~\ref{Dudley} is stated for finite $T$ to keep discretization explicit.
For general index sets, separability allows one to pass to a countable dense
subset, yielding the following corollary (proved in Appendix~\ref{se:Dudley1}).

\begin{corollary}\label{Dudley1}
Let $\{X_t\}_{t\in T}$ be a separable, mean-zero $(\theta,C)$-sub-Weibull process
on $(T,d)$. If we assume that $\gamma<\infty$, then
\[
\mathbb{E}\Bigl[\sup_{t\in T}X_t\Bigr]
\le 4 C K_\theta \int_0^\infty \bigl[\log N(T,d,\varepsilon)\bigr]^{1/\theta}\,d\varepsilon
=: \gamma.
\]
\end{corollary}

The finiteness of $\gamma$ implies $\E[\sup_{t\in T}X_t]<\infty$.  When
$\theta=2$ (sub-Gaussian),~\eqref{eq:heavy_dudley} reduces to Dudley's classical
entropy integral; $\theta=1$ matches the usual sub-exponential analogue.  For
$0<\theta<1$, the exponent $1/\theta>1$ inflates the entropy integrand,
quantitatively capturing the degradation in supremum control induced by heavier
tails.

\subsection{Maximal and Dudley bounds are insufficient for data-adaptive algorithms}

Lemma~\ref{pp-MaximalSW} and Theorem~\ref{Dudley} control maxima and suprema via
\emph{global} complexity terms (e.g., $\log n$ or an entropy integral).  This
worst-case nature is intrinsic: replacing a data-dependent choice by a supremum
over all candidates necessarily pays for the size/entropy of the ambient class.
In modern learning problems, however, algorithms are typically data-adaptive and
probe only a localized, sample-dependent region of the hypothesis space.
Consequently, purely uniform maximal inequalities can be overly conservative and
may become vacuous when used as surrogates for generalization control.

This motivates the information-theoretic approach developed next: rather than
bounding a random, data-dependent choice by a global supremum, we control it in
terms of how much \emph{information} the algorithm extracts from the data.
Section~\ref{sec:gen_bounds} develops tail-adaptive information bounds and a
multiscale (chaining) refinement; Section~\ref{sec:experiments} provides a
numerical illustration of the gap between uniform and information-theoretic
chaining bounds.

\section{Preliminaries on \texorpdfstring{$f$}{f}-Divergence and \Renyi Divergence}\label{sec:pre}

This section collects divergence notions and change-of-measure tools used in the
heavy-tailed analysis.  Classical information-theoretic generalization proofs
typically combine (i) an MGF bound for the loss and (ii) a KL-based variational
representation (Donsker--Varadhan).  When MGFs do not exist, this route breaks
down.  We instead use a \emph{shifted-log} $f_\theta$-divergence tailored to
sub-Weibull tails and compare it to R\'enyi divergence.  The key payoff is a
\emph{decorrelation} lemma: it bounds expectations under a dependent coupling
(e.g., $W$ depends on $S$) by a divergence term plus a decoupled moment term.

\subsection{\texorpdfstring{$f$}{f}-divergence and \Renyi divergence}

\begin{definition}[$f$-divergence and \Renyi mutual information]\label{f-divergence}
Let $P$ and $Q$ be probability measures on a measurable space $(\mathcal{X},\mathcal{F})$, such that $P\ll Q$.
For a convex function $f:(0,\infty)\to\mathbb R$ with $f(1)=0$, the $f$-divergence
is defined as
\begin{equation}\label{eq:f-divergence}
\DD_f(P\|Q)
\;:=\;
\begin{cases}
\displaystyle \int_{\mathcal X} f\left(\frac{dP}{dQ}\right)dQ, & \text{if } P\ll Q,\\
\infty, & \text{otherwise.}
\end{cases}
\end{equation}
For \( \alpha > 0 \) and \( \alpha \neq 1 \), the \emph{\Renyi divergence} of order \( \alpha \) from \( P \) to \( Q \) is defined as
\[
\DD_\alpha(P \| Q) = \frac{1}{\alpha - 1} \log \int_{\mathcal{X}} \left( \frac{dP}{dQ}\right)^\alpha dQ=\frac{1}{\alpha - 1} \log \int_{\mathcal{X}} \left( \frac{dP}{dQ}\right)^{\alpha-1} dP.
\]
\end{definition}


Important special cases of $f$-divergence include KL divergence ($f(x)=x\log x$). The \Renyi divergence is a generalization of the KL divergence, defined below. 
When $\alpha=2$, $\DD_2(P\|Q)=\log\bigl(1+\chi^2(P\|Q)\bigr)$, where $\chi^2(P\|Q)
=
\int_{\mathcal X}
\left({dP}/{dQ}\right)^2dQ - 1$ is the chi-square divergence for $P\ll Q$. As
$\alpha\to\infty$ one obtains
$\DD_\infty(P\|Q)=\log\|dP/dQ\|_\infty$. 
\begin{lemma}[Basic properties of \Renyi divergence, \citet{RenyiKL}]\label{Renyi}
For any $P,Q$ on the same measurable space, $\DD_\alpha(P\|Q)$ is nondecreasing in
$\alpha$. Moreover,
\begin{enumerate}
\item[(i)] for $0<\alpha_1\le \alpha_2$, $\DD_{\alpha_1}(P\|Q)\le \DD_{\alpha_2}(P\|Q)$;
\item[(ii)] $\lim_{\alpha\to 1}\DD_{\alpha}(P\|Q)=\DD_{\mathrm{KL}}(P\|Q)$.
\end{enumerate}
\end{lemma}


%For any $P,Q$ on the same measurable space, $\DD_\alpha(P\|Q)$ is nondecreasing in
%$\alpha$. Moreover, (i) for $0<\alpha_1\le \alpha_2$, $\DD_{\alpha_1}(P\|Q)\le \DD_{\alpha_2}(P\|Q)$; (ii) $\lim_{\alpha\to 1}\DD_{\alpha}(P\|Q)=\DD_{\mathrm{KL}}(P\|Q)$; see \citet{RenyiKL}.



% \begin{definition}[\Renyi divergence]
% Let \( P \) and \( Q \) be probability measures on a measurable space \( (\mathcal{X}, \mathcal{F}) \), such that \( P \ll Q \). For \( \alpha > 0 \) and \( \alpha \neq 1 \), the \emph{\Renyi divergence} of order \( \alpha \) from \( P \) to \( Q \) is defined as
% \[
% \DD_\alpha(P \| Q) = \frac{1}{\alpha - 1} \log \int_{\mathcal{X}} \left( \frac{dP}{dQ}\right)^\alpha dQ=\frac{1}{\alpha - 1} \log \int_{\mathcal{X}} \left( \frac{dP}{dQ}\right)^{\alpha-1} dP.
% \]
% \end{definition}



%\begin{lemma}[Basic properties of \Renyi divergence, \citet{RenyiKL}]\label{Renyi}

%\end{lemma}

\begin{definition}[Mutual information]\label{renyi_Mutual_information}
Let $(W,S)$ be a pair of random variables with joint distribution $P_{W,S}$.
For $\alpha>1$, the \emph{R\'enyi mutual information} is
\begin{equation}\label{eq:Renyi_MI}
I_\alpha(W;S):=\DD_\alpha(P_{W,S}\|P_W\otimes P_S).
\end{equation}
\end{definition}

\subsection{Shifted-log divergences and \Renyi upper bounds}

To obtain tail-adaptive bounds under sub-Weibull index $\theta$, we will use the
following family of convex functions with additive shift $A$,
\begin{equation}\label{eq:f_theta}
f_\theta(x)=x\log^{\frac{1}{\theta}}(x+A), \qquad \theta>0,\ A\ge 1.
\end{equation}
The $A$ prevents singular behavior near $0$ and the power $\theta^{-1}$ matches the power in the optimal scaling $O\bigl((\log n)^{1/\theta}\bigr)$ for sub-Weibull maximum inequality (Lemma \ref{pp-MaximalSW}) in the following heavy-tailed decorrelation arguments.
Strictly speaking, $f_\theta(1)=\log^{1/\theta}(1+A)$, so this is an
unnormalized shifted-log $f$-divergence rather than a normalized $f$-divergence
generator vanishing at independence.  We keep the notation $\DD_{f_\theta}$ and
$I_{f_\theta}$ below for the corresponding shifted-log dependence functionals;
the comparison and data-processing arguments are used only up to this fixed
baseline, and none of the results relies on the functional being zero under
independence.

\subsubsection{\Renyi divergence-based upper bounds}

\begin{lemma}\label{lem:key1}
Let \(P\ll Q\) be probability measures on a measurable space \((\mathcal X,\mathcal F)\), and suppose that
\(\DD_\alpha(P\|Q)<\infty\) for some \(\alpha>1\). Fix \(\theta>0\) and define
\(f_\theta(x):=x\log^{1/\theta}(x+A)\).
Assume that \(A\) satisfies $A\geq 1$ when $\theta\geq1$ and when $\theta<1$
\[
A \ge \exp\!\left(\frac{1}{\alpha-1}\Bigl(\frac{1}{\theta}-1\Bigr)\right)\quad \text{for }1<\alpha\le 2,
\qquad\text{and}\qquad
A \ge \exp\!\left(\frac{1}{\theta}-1\right)\quad \text{for }\alpha>2.
\]
Then
\begin{equation}\label{eq:DF1}
\DD_{f_\theta}(P\|Q) \le [\DD_\alpha(P\|Q)+C_{\alpha,\theta}]^{1/\theta},
\end{equation}
where, for \(1<\alpha\le 2\),
\[
C_{\alpha,\theta}
:=
\begin{cases}
\displaystyle \frac{1}{\alpha-1}\log\bigl(1+A^{\alpha-1}\bigr), & 0<\theta<1,\\[0.9em]
1+(A-1)^{1/\theta}, & \theta\ge 1,
\end{cases}
\qquad\text{and}\qquad
C_{\alpha,\theta}=C_{2,\theta}\ \text{for }\alpha>2.
\]
\end{lemma}

Lemma~\ref{lem:key1} provides a clean {power-type} comparison:
$\DD_{f_\theta}$ is controlled by a $1/\theta$ power of R\'enyi divergence up to
an additive constant inside the power. We use this comparison throughout the
R\'enyi refinements below.


\subsection{Decorrelation inequalities}\label{sub:decorr}

We now state the main change-of-measure tool used throughout the paper.
Given a joint law $\mu$ (e.g., of $(W,S)$) and a product reference law $\nu$
(e.g., $P_W\otimes P_S$), we want to control $\E_\mu[r]$ by a divergence between
$\mu$ and $\nu$ plus a moment term under $\nu$.  In the sub-Gaussian case, this
is commonly done via the Donsker--Varadhan representation of KL divergence.
Here we use a Young-type inequality based on $f_\theta$, which does not require
MGFs.

\begin{lemma}[Decorrelation lemma]\label{de}
Let $\theta>0$ and $A\ge (2^{\lceil \frac{2}{\theta} \rceil-2} \lceil \frac{2}{\theta} \rceil !)^2 \vee 1$. Let $\mu$ and $\nu$ be probability measures on a
measurable space $(\Omega,\mathcal F)$. Let $r\ge 0$ be measurable and assume
$r\in L^1(\mu)$ and $\exp(r^\theta)\in L^1(\nu)$. Let $f_\theta(x)=x\log^{1/\theta}(x+A)$, then
\begin{equation}\label{eq:dec}
\mathbb E_\mu r
\le 2^{1/\theta}\DD_{f_\theta}(\mu\|\nu)+\mathbb{E}_\nu\exp(r^\theta).
\end{equation}
\end{lemma}

The Lemma~\ref{de} will be applied with $\mu=P_{W,S}$ and $\nu=P_W\otimes P_S$ and
with $r$ equal to the data or the reward in RLHF in Section \ref{sec:rlhf}; the first term captures the dependence
between $W$ and $S$, while the second term is a sub-Weibull moment under the
product measure.  The proof of Lemma~\ref{de} is based on a new Young-type
inequality and is deferred to Appendix~\ref{proofse3}.




% In some heavy-tailed settings, it is useful to replace the full exponential
% moment $\E_\nu\exp(r^\theta)$ by a truncated version.  The next lemma provides
% such a refinement.

% \begin{lemma}[Truncated decorrelation lemma]\label{deh}
% Let $\theta\in(0,2]$. Define
% \[
% h(y)=e^{y^\theta}-\sum_{k=0}^{\lfloor 2/\theta\rfloor}\frac{y^{k\theta}}{k!},
% \qquad y\ge 0.
% \]
% Let
% \[
% A\ge
% \max\left\{
% 1,\,
% \left(2^{\lceil 2/\theta\rceil-2}\lceil 2/\theta\rceil!\right)^2,\,
% 2e^{\lfloor 2/\theta\rfloor}
% \right\}
% \]
% and define $f_\theta(x)=x\log^{1/\theta}(x+A)$. Let $\mu$ and $\nu$
% be probability measures on $(\Omega,\mathcal F)$, and let
% $r\ge 0$ be measurable. Assume $r\in L^1(\mu)$ and $h(r)\in L^1(\nu)$. Then
% \begin{equation}\label{eq:dec1}
% \mathbb E_\mu r
% \le 2^{1/\theta}\DD_{f_\theta}(\mu\|\nu)+2\mathbb{E}_\nu h(r).
% \end{equation}
% Moreover, if $A$ satisfies the conditions of Lemma~\ref{lem:key1}, then for
% $\alpha>1$,
% \begin{equation}\label{eq:dec2}
% \mathbb E_\mu r
% \le 2^{1/\theta}(\DD_{\alpha}(\mu\|\nu)+C_{\alpha,\theta})^{1/\theta}+2\mathbb{E}_\nu h(r).
% \end{equation}
% \end{lemma}



%%%%%%%%%%%%%%%%%%%%%%%%%%%%%%%%%%%%%%%%%%%%%%%%
%%%%%%%%%Information Theoretical Generalization%%%%%%%%%%%%
%%%%%%%%%%%%%%%%%%%%%%%%%%%%%%%%%%%%%%%%%%%%%%%%
\section{\Renyi information refinements of maximal and Dudley-type bounds}\label{sec:gen_bounds}



This section turns the sub-Weibull empirical-process tools from
Section~\ref{sec:ep} into information-theoretic bounds for data-dependent
selection. The key mechanism is the shifted-log divergence from
Section~\ref{sec:pre}: it matches sub-Weibull tails through the decorrelation
lemmas and can be upper bounded by R\'enyi mutual information via
Lemma~\ref{lem:key1}.

We begin with the selection problem of bounding $\E[X_W]$ when $W$ is a
data-dependent index.  In the sub-Gaussian setting, KL mutual information gives
such a bound \citep{Russo20}, but Section~\ref{se:whyKLfails} shows that this
route can fail under heavy tails.  Theorem~\ref{infob} gives the shifted-log and R\'enyi replacements, while Theorem~\ref{chaining} lifts the
single-index result to a multiscale Dudley-type bound for data-dependent
indices in metric spaces.


%and a multiscale chaining bound that remains non-vacuous for deterministic continuous learners (Theorem~\ref{genchain}).



% \subsection{\Renyi information refinements of maximal and Dudley-type bounds}

\subsection{Why KL mutual information fails under heavy tails}\label{se:whyKLfails}

Let $S=\{X_i\}_{i=1}^n$ and let $W=W(S)$ be a (possibly randomized) data-dependent
index taking values in $[n]$.  When the $X_i$ are zero-mean i.i.d.\ sub-Gaussian
with variance proxy $\sigma^2$, \citet{Russo20} showed that
\begin{equation}\label{eq:Russo}
\mathbb{E}[X_W] \le \sqrt{2\sigma^2I(W;S)},
\end{equation}
where $I(\cdot;\cdot)$ is KL-based mutual information.  In particular,
choosing $W=\arg\max_{i\in[n]}X_i$ and using $I(W;S)\le \log n$ recovers the
classical maximal inequality
\(
\mathbb{E}[\max_{i\in[n]} X_i]\lesssim \sqrt{\log n}.
\)

In heavy-tailed regimes, KL mutual information may no longer control
$\mathbb{E}[X_W]$.  The following example constructs a randomized maximum
selector $W$ for which $I(W;S)=O(1)$ while $\mathbb{E}[X_W]$ diverges with $n$.

\begin{example}\label{ex:weibull}
Let $S=\{X_i\}_{i=1}^n$ be i.i.d.\ $\mathrm{Weibull}(\theta)$ random variables
with $\theta\in(0,1)$, i.e.\ $\mathbb{P}(X_1\ge x)=\exp(-x^\theta)$ for $x>0$.
Define a randomized selector $W$ by
\[
W=
\begin{cases}
\arg\max_{i\in[n]}X_i, & \text{with probability }\varepsilon,\\
U, & \text{with probability }1-\varepsilon,
\end{cases}
\]
where $U$ is uniform on $[n]$ and independent of $S$.
A direct computation gives
\[
\begin{aligned}
I(W;S)
= \DD_{\mathrm{KL}}(P_{W,S}\,\|\,P_W\otimes P_S) = \frac{n-1}{n}(1-\varepsilon)\log(1-\varepsilon)
   + \frac{(n-1)\varepsilon+1}{n}\log\bigl((n-1)\varepsilon+1\bigr).
\end{aligned}
\]

Choosing $\varepsilon=c/\log n$ for any fixed $c>0$ yields $I(W;S)=O(1)$ as
$n\to\infty$.  On the other hand,
$\mathbb{E}[\max_{i\in[n]}X_i]\asymp (\log n)^{1/\theta}$ (see Appendix~\ref{app:sec2}), and therefore
\[
\mathbb{E}[X_W]
\;\ge\;
\varepsilon\,\mathbb{E}[\max_{i\in[n]}X_i]
=\Omega \bigl((\log n)^{1/\theta-1}\bigr).
\]
Thus, for heavy-tailed Weibull data, the selection bias $\mathbb{E}[X_W]$ can
diverge while the KL mutual information remains bounded.
\end{example}

\subsection{Why \Renyi mutual information and shifted-log \texorpdfstring{$f$}{f}-divergences}

Example~\ref{ex:weibull} reflects a general phenomenon: when MGFs do not exist,
KL-based information can be too weak to control the contribution of rare but
extreme events.  To obtain a tail-adaptive substitute, we applied the proposed shifted-log
family \eqref{eq:f_theta} and use the associated $f_\theta$-mutual information $I_{f_\theta}(W;S)$.
It matches sub-Weibull tails and, via Lemma~\ref{lem:key1}, admits explicit upper bounds in terms of R\'enyi mutual
information.

\begin{theorem}[Tail-adaptive selection bound]\label{infob}
Let $S=\{X_i\}_{i=1}^n$ with $X_i \sim \subw(\theta)$ for $\theta>0$, and let $W=W(S)$ be any $[n]$-valued, possibly randomized, selector. Let $I_{f_\theta}(W;S)$ be the $f_\theta$-mutual information associated with $f_\theta(x)=x\log^{1/\theta}(x+A)$, where $A$ satisfies the condition in Lemma~\ref{de}. Then
\[
\mathbb{E}|X_W|
\;\le\;
\max_{i\in[n]}\|X_i\|_{\psi_\theta}\,\Bigl(2^{1/\theta}I_{f_\theta}(W;S)+2\Bigr).
\]
Moreover, if $A$ satisfies the conditions of Lemma~\ref{lem:key1}, then
\[
\mathbb{E}|X_W|
\;\le\;
\max_{i\in[n]}\|X_i\|_{\psi_\theta}\,
\Bigl(2^{1/\theta}\bigl(I_\alpha(W;S)+C_{\alpha,\theta}\bigr)^{1/\theta}+2\Bigr),
\]
where $I_\alpha(W;S)$ is the R\'enyi-$\alpha$ mutual information and
$C_{\alpha,\theta}$ depend only on $(\alpha,\theta)$.
\end{theorem}

When $\theta=2$ and $\alpha\to 1$, Theorem~\ref{infob} is similar to Russo and Zou's bound
\eqref{eq:Russo} up to universal constants.  Taking $W=\arg\max_{i\in[n]}|X_i|$
also recovers the optimal maximal growth rate
$\mathbb{E}[\max_i|X_i|]\lesssim(\log n)^{1/\theta}$; see Example~\ref{Max} in
Appendix~\ref{se: proof2}.

\subsection{A Dudley-type bound via multiscale information}

Theorem~\ref{infob} bounds a \emph{single} data-dependent selection.  To obtain a
Dudley-type bound for a data-dependent index in a general metric space, we
combine Theorem~\ref{infob} with a multiscale chaining construction, following
\citet{asadi2018chaining,hellstrom2025generalization}.

\begin{definition}[$\varepsilon$-partitions and increasing sequences]\label{def:part}
Let $(\mathcal{W},d)$ be a metric space. A partition $\mathcal{P}=\{A_1,\ldots,A_m\}$ of $\mathcal{W}$ is called an \emph{$\varepsilon$-partition} if for every $i\in[m]$ there exists $w_i\in\mathcal{W}$ such that
\[
A_i \subseteq \mathcal{B}_d(w_i,\varepsilon),
\qquad
\mathcal{B}_d(w_i,\varepsilon):=\{w\in\mathcal{W}: d(w,w_i)\le \varepsilon\}.
\]
A sequence of partitions $\{\mathcal{P}_k\}_{k=k'}^{\infty}$ is called \emph{increasing} if for all $k\ge k'$ and every $A\in\mathcal{P}_{k+1}$ there exists $B\in\mathcal{P}_k$ such that $A\subseteq B$. For each $w\in\mathcal{W}$ and $k\ge k'$, let $[w]_k$ denote the unique set $A\in\mathcal{P}_k$ containing $w$.
\end{definition}

% Examples of $\{\mathcal{P}_k\}$ can be found in Appendix~\ref{Tighterexampleappen}.

\begin{theorem}[Information-theoretic Dudley bound under sub-Weibull tails]\label{chaining}
 Let $(T,d)$ be totally bounded and let $X_t$ be separable
$(\theta,C)$-sub-Weibull process with a.s. uniformly continuous paths, and with finite
Dudley integral in Corollary~\ref{Dudley1}. Let $W=W(S)$ be a $T$-valued random
index (equivalently, $W$ is measurable with respect to $X$). Let
$\{\mathcal{P}_k\}_{k\ge 0}$ be an increasing sequence of partitions such that
$\mathcal{P}_0=\{T\}$ and $\mathcal{P}_k$ is an $e(T)2^{-k}$-partition for $k\ge 1$.
For each $A\in\mathcal{P}_k$, choose $t_A\in T$ so that
$A\subset \mathcal{B}_d\bigl(t_A,e(T)2^{-k}\bigr)$, and if $A\subset B\in \mathcal{P}_{k-1}$, then $t_A\in B$. Let $[W]_k\in\mathcal P_k$ be the cell-valued variable containing $W$. Assume the shift in $f_\theta$ satisfies the condition in Lemma~\ref{de}. Then
\[
\mathbb{E}[X_W]
\;\le\;
Ce(T)\sum_{k=1}^\infty 2^{-(k-1)}
\Bigl(2^{1/\theta}I_{f_\theta}([W]_k;S)+2\Bigr).
\]
% Moreover, the bound admits the following R\'enyi mutual-information refinements.
Assume $A$ satisfies the conditions of Lemma~\ref{lem:key1}. Then, for any
$\alpha>1$ such that $I_\alpha([W]_k;S)<\infty$ for all $k$ and $C_{\alpha,\theta}$ in Lemma~\ref{lem:key1},
\begin{equation}\label{computmeth3}
\mathbb{E}[X_W]
\;\le\;
Ce(T)\sum_{k=1}^\infty 2^{-(k-1)}
\Bigl(
2^{1/\theta}\bigl(I_\alpha([W]_k;S)+C_{\alpha,\theta}\bigr)^{1/\theta}
+2
\Bigr).
\end{equation}
\end{theorem}

When $\theta=2$, Theorem~\ref{chaining} is similar to the mutual-information chaining
bound of \citet[Theorem 4]{asadi2018chaining} up to universal constants.  When
$\theta=1$, it yields the information-theoretic chaining bound under
sub-exponential tails.  For $\theta\in(0,1)$, the exponent $1/\theta>1$ inflates
the per-scale complexity term, quantifying the deterioration caused by heavier
tails.  The proof uses only Orlicz-norm control of increments and multiscale
information, avoiding MGFs and circumventing the KL-based failure exhibited in
Example~\ref{ex:weibull}.

%\subsection{Information-theoretic generalization bounds under sub-Weibullity}

% \subsubsection{Loss function, empirical risk, and generalization error}

% Let $\mathcal{D}$ denote the data space, $\mathcal{W}$ the hypothesis class with
% a metric space $(\mathcal{W},d)$, and $\ell:\mathcal{W}\times\mathcal{D}\to\mathbb{R}_{+}$
% a loss function. Given a sample $S=\{Z_i\}_{i=1}^n$ of i.i.d.\ random variables
% with common distribution $\mu$, the goal is to choose $w\in\mathcal{W}$ that
% minimizes the population risk
% $L_\mu(w)=\mathbb{E}_\mu[\ell(w,Z)]$.
% Since $\mu$ is unknown, one instead works with the empirical risk
% $L_S(w)=\frac{1}{n}\sum_{i=1}^n \ell(w,Z_i)$. A learning algorithm is modeled as a Markov kernel $P_{W\mid S}$ mapping the
% data $S$ to a hypothesis $W\in\mathcal{W}$. When $S\sim\mu^{\otimes n}$, the
% joint law is $P_{S,W}=\mu^{\otimes n}\otimes P_{W\mid S}$.
% The \emph{generalization error} is
% \[
% \operatorname{gen}(W,S)
% := L_\mu(W)-L_S(W)=\frac{1}{n}\sum_{i=1}^n \bar{\ell}(W,Z_i),
% \]
% where the centered loss is
% $\bar{\ell}(w,Z)=\mathbb{E}_\mu[\ell(w,Z)]-\ell(w,Z)$.


% %%%%%%%%%%%%%%%%%%%%%%%%%%%%%%%%%%%%%%%%%%%%%%%%%%%%%%%%%%%%%%%%%%%%%%%
% %%%%%%%%%%%%%%%Classical information theoretic bounds%%%%%%%%%%%%%%%%%%
% %%%%%%%%%%%%%%%%%%%%%%%%%%%%%%%%%%%%%%%%%%%%%%%%%%%%%%%%%%%%%%%%%%%%%%%
% %\subsubsection{Classical information-theoretic bounds}

% Suppose that for every $w\in\mathcal{W}$, the centered loss
% $\ell(w,Z)-\mathbb{E}\ell(w,Z)$ is $\sigma$-sub-Gaussian under $Z\sim\mu$.
% Then, \cite{xu2017} controlled generalization by KL mutual information
% \[
% \mathbb{E}\bigl|\operatorname{gen}(W,S)\bigr|
% \;\le\; \sqrt{\frac{2\sigma^2}{n}\,I(S;W)}.
% \]
% If the learner is deterministic and $W=\varphi(S)$, then $I(S;W)=H(W)$ measures
% how much the output varies across samples. If $W$ is independent of $S$, then
% $I(S;W)=0$ and the bound yields zero generalization error.

% \subsubsection{From sub-Gaussianity to sub-Weibullity}

% Proposition~\ref{Xu} relies critically on sub-Gaussian tails (and hence MGFs).
% In the heavy-tailed sub-Weibull regime, KL mutual information can be too weak to
% control generalization---it may remain small even when the generalization error
% is large (cf.\ Example~\ref{ex:weibull}).  We therefore replace KL by the
% shifted-log divergence induced by $f_\theta(x)=x\log^{1/\theta}(x+A)$.

% \begin{theorem}[Generalization bounds for sub-Weibull losses]\label{gen}
% Let $S:=\{Z_i\}_{i=1}^n$ be i.i.d.\ with law $\mu$, and let $W:=W(S)$ be a (possibly
% randomized) learner measurable from data to a metric space $(\mathcal{W},d)$. Assume
% \begin{equation}\label{eq:sub-Weibull-loss}
% \sup_{w\in\mathcal{W}}\|\bar{\ell}(w,Z)\|_{\psi_\theta}<\infty.
% \end{equation}
% Assume the shift in $f_\theta$ satisfies the condition in Lemma~\ref{deh}.
% Then there exists a constant $M_\theta>0$ that depends only on $\theta$
% (see~\eqref{eq:theta} in Appendix) such that
% \[
% \mathbb{E}\bigl|\operatorname{gen}(W,S)\bigr|
% \;\le\;
% \frac{M_\theta}{\sqrt{n}}\,
% \sup_{w\in\mathcal{W}}\|\bar{\ell}(w,Z)\|_{\psi_\theta}\,
% \bigl(2^{1/\theta} I_{f_\theta}(S;W)+4\bigr),~0<\theta\le2.
% \]
% Moreover, under the conditions of Lemma~\ref{lem:key1},
% $I_{f_\theta}(S;W)$ can be further bounded in terms of R\'enyi mutual
% information $I_\alpha(S;W)$.
% \end{theorem}

% When $\theta=2$ and $\alpha\to1$, Theorem~\ref{gen} recovers
% Proposition~\ref{Xu} up to constants.  For $\theta\in(0,1)$, the dependence on
% $I_{f_\theta}(S;W)$ makes the information--tail tradeoff explicit.

% \begin{remark}[On the uniform sub-Weibull loss assumption]\label{rem:uniform-subw}
% For many models, \eqref{eq:sub-Weibull-loss} can be verified by combining
% sub-Weibull tails of the data with Lipschitz/ boundedness properties of the
% predictor; see Example~\ref{ex:squareloss} in Appendix~\ref{ap:SGDL} for a square-loss
% illustration.
% The global supremum in \eqref{eq:sub-Weibull-loss} keeps the statement simple but
% can be conservative: in practice it may suffice to control the sub-Weibull scale
% locally on the region of $\mathcal W$ explored by the algorithm (e.g., in
% expectation under a posterior or on a data-dependent neighborhood of $W$).
% We keep the uniform form to emphasize the core information--tail mechanism and
% to facilitate comparisons with KL-based results.
% \end{remark}




% % \subsubsection{High-probability bounds and PAC-Bayes perspective}

% Beyond bounds in expectation, PAC-Bayes theory provides high-probability
% generalization guarantees. Existing PAC-Bayes bounds typically rely on bounded,
% sub-Gaussian or sub-exponential assumptions
% \citep{alquier2016properties, alquier2018simpler, germain2016pac, catoni2004statistical}.
% Heavy-tailed losses have also been studied under finite moment conditions
% \citep{alquier2018simpler}, but finite $q$-moment-based bounds generally scale
% polynomially in $n$ \citep[Corollary 7.1]{zhang2020concentration} and can be much
% weaker than the logarithmic dependence $(\log n)^{1/\theta}$ available under
% sub-Weibull tails (see Example~\ref{Max} in Appendix~\ref{se: proof2}).
% The following theorem extends \citet[Theorem~6]{chu2023unified} to sub-Weibull
% losses.

% \begin{theorem}[High-probability sub-Weibull bound]\label{gen1}
% Write $\log_+(u):=\max\{\log u,0\}$. Let $S=\{Z_i\}_{i=1}^n$ be i.i.d.\ with law
% $\mu$, and let $W=W(S)$ be specified by $P_{W\mid S}$ in a metric space
% $(\mathcal{W},d)$. Assume
% \[
% \sup_{w\in\mathcal{W}}\|\bar{\ell}(w,Z)\|_{\psi_\theta}<\infty
% \quad\text{for some }0<\theta\le2.
% \]
% Then there exists $E_\theta>0$ such that for any $\delta\in(0,1)$,
% with $P_S$-probability at least $1-\delta$,
% \[
% \mathbb{E}_{P_{W\mid S}}\bigl|\operatorname{gen}(W,S)\bigr|
% \le
% \frac{M_\theta}{\sqrt{n}}\,
% \sup_{w\in\mathcal{W}}\|\bar{\ell}(w,Z)\|_{\psi_\theta}
% \Bigl(
% \mathbb{E}_{P_{W\mid S}}\!\left\{
% \bigl[\log_+ (dP_{W\mid S}/dP_W)\bigr]^{1/\theta}
% \right\}
% +E_\theta+(\log(3/\delta))^{1/\theta}
% \Bigr),
% \]
% where $M_\theta>0$ is given in \eqref{eq:theta} in Appendix.
% \end{theorem}

% \noindent\textit{From conditional complexity to mutual information.}
% For any $A\ge1$,
% $\log_+\!\left(\frac{dP_{W\mid S}}{dP_W}\right)
% \le \log\!\left(\frac{dP_{W\mid S}}{dP_W}+A\right)$,
% so
% \[
% \mathbb{E}_{P_{W\mid S}}\!\left\{
% \bigl[\log_+ (dP_{W\mid S}/dP_W)\bigr]^{1/\theta}
% \right\}
% \le
% \mathbb{E}_{P_{W\mid S}}\!\left\{
% \Bigl[\log\!\left(\frac{dP_{W\mid S}}{dP_W}+A\right)\Bigr]^{1/\theta}
% \right\}
% =\DD_{f_\theta}(P_{W\mid S}\|P_W),
% \]
% which can be
% upper bounded by R\'enyi divergence using Lemma~\ref{lem:key1}.




% \subsubsection{From a single hypothesis to metric chaining}

% Theorems~\ref{gen} and~\ref{gen1} are \emph{single-scale} bounds: they measure
% data dependence through $I_{f_\theta}(W;S)$ (or $I_\alpha(W;S)$) for the final
% output $W$.  This can be vacuous when $W$ is a continuous deterministic function
% of $S$, because then $P_{W\mid S}$ is singular with respect to $P_W$ and
% $I_\alpha(W;S)=\infty$.  To address this, we use a multiscale discretization of
% $\mathcal W$ in the spirit of \citet{asadi2018chaining}, adapted here to
% sub-Weibull tails.

% Assume the centered loss $\{\bar\ell(w,Z)\}_{w\in\mathcal W}$ has
% sub-Weibull increments with respect to $d$:
% \[
% \|\bar\ell(u,Z)-\bar\ell(v,Z)\|_{\psi_\theta}\le d(u,v),\qquad u,v\in\mathcal W.
% \]
% Combining this increment condition with the information-theoretic chaining
% result (Theorem~\ref{chaining}) yields the following multiscale generalization
% bound.

% \begin{theorem}[Chaining generalization bound]\label{genchain}
% Let $S=\{Z_i\}_{i=1}^n$ be i.i.d.\ with law $\mu$, and let $W=W(S)$ be measurable.
% Assume $\{\bar\ell(w,Z)\}_{w\in\mathcal W}$ is a separable $(\theta,1)$-sub-Weibull
% process on $(\mathcal W,d)$ with $\theta\in(0,2]$, and assume the Dudley integral $\gamma$ for process $\ell(w,Z)$ in Corollary~\ref{Dudley1} is finite.
% Let $\{\mathcal P_k\}_{k\ge0}$ be an increasing sequence of partitions of
% $\mathcal W$, where $\mathcal P_k$ is an $e(\mathcal W)2^{-k}$-partition.
% For each $A\in\mathcal P_k$, choose $t_A\in\mathcal W$ such that
% \[
% A\subset \mathcal B_d \bigl(t_A,e(\mathcal W)2^{-k}\bigr),
% \]
% and if $A\subset B\in\mathcal P_{k-1}$, then $t_A\in B$.
% Let $[W]_k\in\mathcal P_k$ be the cell-valued variable containing $W$. Then there
% exists $M_\theta>0$ (see \eqref{eq:theta} in Appendix) such that, whenever the
% shift in $f_\theta$ satisfies the condition in Lemma~\ref{deh},
% \begin{equation}\label{eq:cgb}
% \mathbb E \bigl[\operatorname{gen}(W,S)\bigr]
% \le
% \frac{e(\mathcal W)M_\theta}{\sqrt n}
% \sum_{k=1}^\infty 2^{-(k-1)}
% \bigl(2^{1/\theta}I_{f_\theta}([W]_k;S)+4\bigr).
% \end{equation}
% \end{theorem}

% When $\theta=2$, \eqref{eq:cgb} recovers (up to constants) the chaining mutual-information
% bound of \citet{asadi2018chaining}.  When $\theta=1$, it yields the sub-exponential
% analogue; see also \citet{zhou2023stochastic} for finite-MGF settings.
% For $0<\theta<1$, the exponent $1/\theta>1$ enlarges each scale contribution,
% quantifying the impact of heavier tails.

% \begin{example}[Quadratic mean estimation; deterministic ERM]\label{ex:mean-chaining}
% Let $\mathcal W=[-1,1]$ and consider $\ell(w,Y)=(w-Y)^2$, where $Y$ is symmetric
% Weibull with $\mathbb P(|Y|\ge t)=\tfrac12 e^{-t^\theta}$ for $\theta\in(0,1)$.
% Given $S=\{Y_i\}_{i=1}^n$, define the ERM
% \[
% W
% =\arg\min_{w\in[-1,1]}L_S(w)
% =\clip\!\left(\frac1n\sum_{i=1}^n Y_i,\;[-1,1]\right).
% \]
% This $W$ is deterministic and continuous in $S$, so $P_{W\mid S}$ is singular
% w.r.t.\ $P_W$ and $I_\alpha(W;S)=\infty$ for every $\alpha>1$; hence
% Theorem~\ref{gen} is vacuous.

% Now discretize $[-1,1]$ by dyadic partitions
% \[
% \mathcal P_k
% =
% \Bigl\{[m2^{-(k-1)},(m+1)2^{-(k-1)})\cap[-1,1]: m\in\mathbb Z\Bigr\},\quad k\ge1,
% \]
% and let $W_k=[W]_k$ be the cell containing $W$.  Although $I_\alpha(W;S)=\infty$,
% $W_k$ is discrete and its R\'enyi mutual information grows at most linearly in
% $k$ under mild regularity (e.g., a bounded-below density for the sample mean on
% $[-1,1]$; see the proof in Appendix~\ref{Tighterexampleappen}).
% Since $\sum_{k\ge1}k\,2^{-k}<\infty$, the multiscale sum in
% Theorem~\ref{genchain} is finite and yields a finite bound.
% \end{example}

\section{Applications to RLHF with Heavy-Tailed Rewards}\label{sec:rlhf}

This section studies our heavy-tailed information-theoretic bounds: safety alignment of large language models (LLMs) under heavy-tailed reward functions. Recent works on statistical principles can directly design  trustworthy alignment methods \citep{su2026large}, and we focus more on robust estimation principles.

% and (ii) generalization analysis for stochastic optimization under heavy-tailed data (Appendix~\ref{se:SGDL}).

We study reinforcement learning from human feedback (RLHF) for aligning LLMs.
 Let $\pi(y\mid x)$ denote a policy (a Markov kernel from prompts to responses),
where $x\in\mathcal X$ is a prompt and $y\in\mathcal Y$ is a generated response.
Prompts are drawn from a distribution $\rho_X$ on $\mathcal X$, and
$\pi(\cdot\mid x)$ specifies the conditional distribution of responses.
In this sense, RLHF can be viewed as a contextual bandit problem.

Given a reference policy $\pi_0(\cdot\mid x)$ (typically a pretrained model),
the goal is to construct a new policy $\pi$ that achieves higher expected reward
$r(x,y)$.  In practice, $r$ is learned from preference data and can be
misspecified; we refer to the true human-aligned objective as the \emph{gold
reward} and the learned reward model as the \emph{proxy reward}.

\subsection{KL-constrained and best-of-\texorpdfstring{$n$}{n} alignment}

A standard formulation of RLHF is the KL-constrained optimization
\citep{ouyang2022training}:
\begin{equation}\label{opt}
\max_{\pi}\;
\mathbb{E}_{x\sim\rho_X,\;y\sim\pi(\cdot\mid x)}[r(x,y)]
\quad\text{s.t.}\quad
\DD_{\mathrm{KL}}\bigl(\pi(\cdot\mid x)\,\|\,\pi_0(\cdot\mid x)\bigr)\le\epsilon
\quad\forall x,
\end{equation}
where $\epsilon>0$ is an information budget.
By Donsker--Varadhan representation, the optimizer is an exponentially tilted distribution \citep{yang2024asymptotics} with a Lagrange multiplier $\lambda>0$
\begin{equation}\label{KLsol}
\pi^{*}_{\mathrm{KL}}(y\mid x)
=\frac{\pi_0(y\mid x)\exp(r(x,y)/\lambda)}
{\mathbb{E}_{Y\sim\pi_0(\cdot\mid x)}\exp(r(x,Y)/\lambda)}.
\end{equation}

An alternative inference-time strategy is the \emph{best-of-$n$} policy
\citep{stiennon2020learning,nakano2021webgpt}.
Given a prompt $x$, one samples $n$ i.i.d.\ responses
$Y_1,\dots,Y_n\sim\pi_0(\cdot\mid x)$ and selects
$Y_n^*\in \arg\max_{i\in [n]} r(x,Y_i)$. The resulting policy
$\pi_n(\cdot\mid x)$ is the distribution of $Y_n^*$.
Under bounded rewards and regularity assumptions, best-of-$n$ policies are
asymptotically close to the KL-constrained optimizer
\citep{mroueh2025information,yang2024asymptotics}.

\subsection{Catastrophic Goodhart and \Renyi-regularized RLHF}

Under light-tailed assumptions, one can derive reward guarantees using KL-based
arguments and data processing \citep{mroueh2025information}.
When the proxy reward is heavy-tailed, however, KL regularization can fail
catastrophically: the proxy reward can diverge even when the KL constraint is
small, a phenomenon known as \emph{catastrophic Goodhart}
\citep{kwa2024catastrophic}.
Formally, for any heavy-tailed distribution $Q$ and any $\epsilon>0$, there exist
distributions $P$ with \emph{arbitrarily large mean} such that
$\DD_{\mathrm{KL}}(P\|Q)\le\epsilon$.
Intuitively, exponential tilting \eqref{KLsol} can place disproportionate mass
on extreme-reward events, while the KL penalty---an average discrepancy under the
candidate policy---need not scale proportionally with such rare tail shifts.

To mitigate this, we consider RLHF with R\'enyi-$\alpha$ regularization:
\begin{equation}\label{optRENYI}
\max_{\pi}\;
\mathbb{E}_{x\sim\rho_X,\;y\sim\pi(\cdot\mid x)}[r(x,y)]
\quad\text{s.t.}\quad
\DD_{\alpha}\bigl(\pi(\cdot\mid x)\,\|\,\pi_0(\cdot\mid x)\bigr)\le\epsilon
\quad\forall x,
\end{equation}
for some $\alpha>1$.
Compared to KL, R\'enyi penalizes high density ratios more
aggressively through a power $\alpha$, which induces a qualitatively different
structure in the optimizer.

\begin{lemma}[Solution to R\'enyi-constrained optimization]\label{renyisol}
Let $\alpha>1$ and $\epsilon>0$. Fix $x\in\mathcal X$ and write $P_0:=\pi_0(\cdot\mid x), r_x(y):=r(x,y)$. Assume that policies in \eqref{optRENYI} are optimized over probability measures $\ll P_0$, and assume
$r_x\in L^{\frac{\alpha}{\alpha-1}}(P_0)$.
Let
\[r_x^\star:=\operatorname*{ess\,sup}_{y\sim P_0} r_x(y),
\qquad
M_x:=\{y\in\mathcal Y:r_x(y)=r_x^\star\},
\qquad
m_x:=P_0(M_x),
\]
with the convention $-\log 0=+\infty$.
Then the optimizer of \eqref{optRENYI} satisfies
\begin{enumerate}
\item If $\epsilon\ge -\log m_x$, then the R\'enyi ball is large enough to concentrate all probability mass on
the maximal-reward set. One has an optimal policy
$
\pi_M(dy\mid x)
=
\frac{\mathbf 1_{M_x}(y)}{m_x}\,\pi_0(dy\mid x)$.

\item If $\epsilon<-\log m_x$, then the R\'enyi constraint is active at the optimizer, i.e.
$
\mathcal D_\alpha\bigl(\pi^*(\cdot\mid x)\|\pi_0(\cdot\mid x)\bigr)
=
\epsilon$. Moreover, there exists a threshold $t\in\mathbb R$ such that
\[
\frac{d\pi^*(\cdot\mid x)}{d\pi_0(\cdot\mid x)}(y)
=
\frac{(r(x,y)-t)_+^{\frac1{\alpha-1}}}
{\mathbb E_{Y\sim\pi_0(\cdot\mid x)}
[(r(x,Y)-t)_+^{\frac1{\alpha-1}}]}.
\]
Equivalently, when densities exist,
$\pi^*(y\mid x)
=
\frac{
\pi_0(y\mid x)(r(x,y)-t)_+^{\frac1{\alpha-1}}
}{
\mathbb E_{Y\sim\pi_0(\cdot\mid x)}
[(r(x,Y)-t)_+^{\frac1{\alpha-1}}]
}$. The threshold $t$ is chosen so that the R\'enyi constraint is met with equality.
\end{enumerate}
\end{lemma}

Note that the first case is degenerate and rarely occurs in continuous problems. For instance, if $r(x,\cdot)$ is continuous and has a unique maximizer under a non-atomic reference policy $\pi_0(\cdot\mid x)$, then $M_x$ is a singleton and hence $\pi_0(M_x\mid x)=0$, so only the second case applies for any finite $\epsilon$.

Lemma~\ref{renyisol} highlights the key structural difference from KL: R\'enyi
regularization yields a \emph{truncated power-law} reweighting of the reference
policy, whereas KL yields exponential tilting.  The truncation limits the
influence of extreme-reward tail events, which is precisely the regime where KL
can be vulnerable.
The special case $\alpha=2$ corresponds to $\chi^2$-preference optimization, for
which $\pi^*(y\mid x)\propto\pi_0(y\mid x)(r(x,y)-t)_+$; see
\cite{huang2025correcting}.

\begin{theorem}[Reward guarantees for RLHF with R\'enyi regularization]\label{thm:rlhf}
Fix $x\in\mathcal X$ and let $\pi_0(\cdot\mid x)$ be the reference policy.
Abbreviate $\pi(\cdot\mid x)$ and $\pi_0(\cdot\mid x)$ by $\pi$ and $\pi_0$.
Define the centered reward
\[
\bar r(x,y)=r(x,y)-\mathbb{E}_{Y\sim\pi_0}[r(x,Y)].
\]
Assume $\|\bar r(x,Y)\|_{\psi_\theta}\le C$ for some $\theta>0$ and
$Y\sim\pi_0(\cdot\mid x)$.
Then for any policy $\pi(\cdot\mid x)$,
\[
\mathbb{E}_{\pi} r(x,Y) - \mathbb{E}_{\pi_0} r(x,Y)
\;\le\;
C\Bigl(2^{1/\theta}[\DD_\alpha(\pi\|\pi_0)+C_{\alpha,\theta}]^{1/\theta}
+2\Bigr),
\]
where $C_{\alpha,\theta}$ is the constant in Lemma~\ref{lem:key1}.
In particular, for the optimizer $\pi^*$ of~\eqref{optRENYI},
\[
\mathbb{E}_{x\sim\rho_X}\mathbb{E}_{Y\sim\pi^*(\cdot\mid x)} r(x,Y)
\;\le\;
\mathbb{E}_{x\sim\rho_X}\mathbb{E}_{Y\sim\pi_0(\cdot\mid x)} r(x,Y)
+ C\Bigl(2^{1/\theta} [\epsilon+C_{\alpha,\theta}]^{1/\theta}+2\Bigr).
\]
\end{theorem}

Theorem~\ref{thm:rlhf} shows that, 
R\'enyi-constrained RLHF rules out unbounded reward inflation at any fixed (divergence) budget under sub-Weibull reward tails, thereby preventing catastrophic Goodhart in the reward model.

\subsection{Best-of-\texorpdfstring{$n$}{n} policies under heavy-tailed rewards}

We now analyze whether catastrophic Goodhart arises for best-of-$n$ policies.
We adopt the following structural assumption
\citep{beirami2024theoretical,mroueh2025information}.

\begin{assumption}[Reward structure]\label{ass:reward-structure}
For each context $x$, define
\[
r(x)=r(x,Y),\quad Y\sim\pi_0(\cdot\mid x),
\qquad
r_n(x)=\max_{i\in [n]} r(x,Y_i),
\quad Y_i\overset{\mathrm{i.i.d.}}{\sim}\pi_0(\cdot\mid x).
\]
There exists a stochastic map $H_x$ such that
$H_x(r(x))\overset{d}{=}\pi_0(\cdot\mid x)$ and
$H_x(r_n(x))\overset{d}{=}\pi_n(\cdot\mid x)$.
\end{assumption}

Under Assumption~\ref{ass:reward-structure} and the data processing inequality
for R\'enyi divergence (cf.\ \citealp{Yihong25}, p.~120),
\[
\DD_\alpha(\pi_n(\cdot\mid x)\|\pi_0(\cdot\mid x))
\;\le\;
\DD_\alpha(r_n(x)\|r(x)).
\]

\begin{theorem}[Best-of-$n$ under heavy-tailed rewards]\label{renyiregret}
Assume $\|r(x,Y)\|_{\psi_\theta}\le C$ for some $\theta>0$, $\alpha>1$ and
$Y\sim\pi_0(\cdot\mid x)$, and assume Assumption~\ref{ass:reward-structure}.
If the trust-region budget $n\le e^\epsilon$ for $\epsilon>0$, then
\[
\DD_\alpha(\pi_n(\cdot\mid x)\|\pi_0(\cdot\mid x))
\le \frac{1}{\alpha-1}\log\!\left(\frac{n^\alpha}{\alpha(n-1)+1}\right)
\le  \epsilon,
\]
and
\[
\E_{\pi_n}r-\E_{\pi_0}r
\le
C \left( 2^{1/\theta} \left[ \DD_\alpha(\pi_n \| \pi_0) +C_{\alpha,\theta}\right]^{1/\theta} + 2 \right)
\le
C \left( 2^{1/\theta} \left[ \e +C_{\alpha,\theta}\right]^{1/\theta} + 2\right),
\]
where $C_{\alpha,\theta}$ depends only on $(\alpha,\theta)$ in
Lemma~\ref{lem:key1}.
\end{theorem}

\begin{remark}
For large $n$, $\DD_\alpha(\pi_n\|\pi_0)\asymp\log n$, yielding a reward bound of
order $\log^{1/\theta}n$.  This matches the scaling suggested by maximal
inequalities and shows that, unlike KL-regularized RLHF, best-of-$n$ policies do
not exhibit severe catastrophic Goodhart behavior under sub-Weibull rewards.
\end{remark}

\section{Numerical Studies}\label{sec:experiments}

In this section, we connect the theory to computation in three steps. We first examine the tightness of the proposed generalization bounds in a controlled example where single-scale information bounds fail. We then study
\Renyi\!-regularized RLHF under controlled heavy-tailed reward perturbations, using an end-to-end alignment pipeline based on Group Relative Policy Optimization (GRPO; \citealt{shao2024deepseekmath}) and state-of-the-art reward models. Finally, we move from synthetic tail amplification to an adversarial token-space reward-hacking setting, where heavy-tailed reward responses arise from reward-model vulnerabilities. %The code is available at \url{https://github.com/NeXAIS/SubW-Alignment}.


\subsection{Illustrating Tighter Generalization Bounds}\label{sec:tightness}


We revisit the two-dimensional process example of
\citet{asadi2018chaining} to illustrate a key phenomenon in the heavy-tailed,
data-adaptive setting: a single-scale mutual-information bound can be vacuous
even when a multiscale \emph{chained} information bound remains finite and
informative. Table~\ref{tab:bounds} summarizes the comparison.


\begin{table}[h]
\centering
% \caption{$\E[X_W]$ and its upper bounds ($\theta=0.5$).}8
\caption{\small $-\E[X_W]$ and upper bounds for $\theta=0.5$.
MI: Theorem~\ref{infob}; CM: Theorem~\ref{Dudley}; 
R\'enyi CMI: Theorem~\ref{chaining}.}\label{tab:bounds}
\renewcommand{\arraystretch}{1}
\begin{tabular}{c|ccccccc}
\hline
\(\varepsilon\) & $1/20$ & $1/30$ & $1/40$ & $1/50$ & $1/100$ & $1/200$ & $1/400$ \\
\hline
MI & $\infty$ & $\infty$ & $\infty$ & $\infty$ & $\infty$ & $\infty$ & $\infty$ \\
CM & 832.01 & 832.01 & 832.01 & 832.01 & 832.01 & 832.01 & 832.01 \\
R\'enyi CMI & 71.93 & 71.46 & 71.30 & 71.22 & 71.12 & 71.10 & 71.09 \\
$-\E[X_W]$ & 0.180 & 0.120 & 0.090 & 0.072 & 0.036 & 0.018 & 0.009 \\
\hline
\end{tabular}
\end{table}


% \textit{Setup.}
Let $S=(Z_1,Z_2)$ with i.i.d. Weibull$(\theta)$ coordinates, i.e.,
$\mathbb{P}(Z_i\ge t)=\exp(-t^\theta)$ for $t\ge 0$.
Consider the canonical Weibull process on the unit circle,
\[
X_\phi := Z_1\cos\phi + Z_2\sin\phi,\qquad \phi\in[0,2\pi),
\]
with loss $l(\phi,S):=-X_\phi$.
One can verify that $\{-X_\phi\}_{\phi\in[0,2\pi)}$ is a $(\theta,C)$-sub-Weibull process (Definition~\ref{subWp}) with respect to the arc-length metric $d(\phi_1,\phi_2)=|\phi_1-\phi_2|$.

%see Appendix~\ref{Tighterexampleappen}.

We define a randomized data-dependent selector by taking a measurable minimizer
$\phi^\star(S)\in\arg\min_\phi l(\phi,S)$ and adding a small perturbation with $\xi\perp\!\!\!\perp S$,
\[
W=\phi^\star(S)\oplus \xi \pmod{2\pi},\qquad
\mathbb P(\xi=0)=\varepsilon,\qquad
\xi\mid(\xi\neq0)\sim\mathrm{Unif}(-\pi,\pi).
\]
% \textit{Bounds compared.}
We compare four ways to upper bound $\E[X_W]$:
\begin{enumerate}
\item[(i)] the single-scale mutual-information (MI) bound (Theorem~\ref{infob});
\item[(ii)] the classical chaining/Dudley entropy bound (CM, Theorem~\ref{Dudley});
\item[(iii)] the chained mutual-information bound (R\'enyi CMI, Theorem~\ref{chaining});
\item[(iv)] the exact value of $-\E[X_W]$ for this example. 
\end{enumerate}
As shown in
Table~\ref{tab:bounds}, the single-scale MI bound is infinite. The chained information bound remains finite and is substantially
tighter than the classical entropy-based chaining bound.

% \textit{Why MI is infinite but chaining is finite.}
For any $\varepsilon>0$, the conditional law $P_{W\mid S=s}$ has an atom at
$\phi^\star(s)$, whereas the marginal law $P_W$ is non-atomic. Hence
$P_{W\mid S=s}\not\ll P_W$ for $P_S$-a.e.\ $s$, implying
$I(W;S)=I_\alpha(W;S)=\infty$ and rendering the single-scale MI bound
vacuous. In contrast, the quantized variables $[W]_k$ used in the chaining
construction take values in a \emph{finite} partition whose cells all have
strictly positive marginal probability, thanks to the uniform component of
$\xi$. Therefore each $I_\alpha([W]_k;S)$ is finite, and the chained bound
remains informative.


% Full derivations and calculations are given in Appendix~\ref{Tighterexampleappen}.



\subsection{RLHF under Heavy-Tailed Rewards}\label{sec:rlhf-experiments}

To illustrate, the mechanism by which R\'enyi-regularized RLHF mitigates catastrophic Goodhart effects, we empirically study R\'enyi-regularized RLHF under \emph{controlled} heavy-tailed reward perturbations. Since commonly used open-source reward models often appear approximately light-tailed on in-distribution prompts, we emphasize that the goal of these experiments is not to claim that the underlying reward model is intrinsically heavy-tailed; rather, we construct heavy-tailed training signals in a principled way to isolate the interaction between tail behavior and the choice of divergence regularizer. Starting from the constrained formulation~\eqref{optRENYI}, we adopt the
Lagrangian optimization:
\begin{equation}\label{optRENYI_penalty}
\pi_{\beta}^{*}
=
\arg\max_{\pi}\;
\mathbb{E}_{x \sim \mathcal{D}}[
\mathbb{E}_{y \sim \pi(\cdot\mid x)}[ r(x,y)]-
\beta\, D_\alpha(\pi(\cdot\mid x) \,\|\, \pi_{0}(\cdot\mid x))],
\quad \beta \ge 0,
\end{equation}
where $\pi_0(\cdot\mid x)$ is the reference policy, $\alpha\ge 1$ is the \Renyi
order (with $\alpha=1$ reducing to KL), and $\beta$ controls the trade-off between
reward maximization and conservatism.

\paragraph{Pipeline.}
We implement an end-to-end alignment pipeline based on Group Relative Policy
Optimization (GRPO; \citealt{shao2024deepseekmath}) with three stages:
\begin{itemize}
\item[(i)] \textbf{Supervised fine-tuning (SFT).}
We start from Qwen2.5-1.5B~\citep{qwen2.5} and perform supervised fine-tuning on \texttt{UltraChat-200k} to obtain an initial policy
that reliably follows the dialogue format and instructions. All runs are initialized from the same SFT checkpoint.

\item[(ii)] \textbf{Reward modeling.}
We train a lightweight proxy reward model (Qwen2.5-0.5B) on the
\texttt{hh-rlhf} training split to provide scalable training-time feedback.
The proxy model takes concatenated \texttt{prompt+completion} pairs as input
and is used \emph{only} to produce online rewards during GRPO. For
evaluation, we additionally use a stronger reward model,
Skywork-Reward-V2-Qwen3-1.7B~\citep{liu2025skywork}, as a surrogate evaluator
for the (unobserved) gold reward; its output is never used during training. For simplicity, we refer to its scores as the \emph{gold reward}. This surrogate evaluator is trained on a much larger dataset and achieves higher
correlation with human judgments, but it is too computationally expensive to
run routinely during RL optimization. Throughout this section, we refer to the lightweight
training-time model's scores as the \emph{proxy reward}.




\item[(iii)] \textbf{RL optimization and evaluation.}
We optimize the SFT policy via GRPO using the proxy reward model as the training-time critic.
We use a frozen reference model initialized from the same checkpoint as the policy to compute divergence terms and to
recalibrate reward statistics. This two-reward-model protocol is standard in the RLHF literature
(e.g., \citealt{kwa2024catastrophic}) and provides a practical balance between computational efficiency and evaluation fidelity.
\end{itemize}





\paragraph{Experimental settings.}
We fix $\beta=1$, sweep $\alpha\in\{1,2,\ldots,10\}$, train on the \texttt{hh-rlhf} training split for up to $1500$ GRPO steps, and evaluate every $10$ steps on 128 held-out prompts. Full training and evaluation hyperparameters (generation settings, batch sizes, optimizer details, truncation lengths, and reward recalibration) are provided in Appendix~\ref{app:rlhf_details}.\\

\paragraph{Inducing heavier tails as controlled stress tests.}
Most empirical RLHF analyses implicitly treat reward-model scores as approximately Gaussian, or more generally light-tailed~\citep{yan2024reward,yang2024bayesian}.
Our theory predicts that the choice of divergence regularizer matters most when rare, high-magnitude reward events are plausible, possibly due to distribution shift, reward-model miscalibration, or occasional outlier completions. Since widely used open-source reward models often appear approximately light-tailed on in-distribution prompts, we do not claim that the base proxy reward is intrinsically heavy-tailed; instead, we \emph{stress-test} this regime by constructing heavy-tailed training signals in controlled ways while keeping the rest of the pipeline fixed. Our construction follows \cite{bakhshizadeh2023sharp} and applies an order-preserving power transform to the recalibrated proxy reward before optimization:
\[
\tilde r(x,y)=\operatorname{sign}\!\big(r(x,y)\big)\,\big|r(x,y)\big|^{\kappa},
\qquad \kappa \in \{1,6,8,10\}.
\]
Here, $\kappa=1$ corresponds to the original reward. The $t\mapsto \operatorname{sign}(t)|t|^{\kappa}$ is strictly increasing on $\mathbb{R}$ for $\kappa>0$ and therefore preserves the preference ordering induced by the original proxy score. Moreover, if the original proxy score is sub-Gaussian, then the transformed reward $\tilde r$ is sub-Weibull with parameter \(2/\kappa\).  


% To confirm that the power transform does amplify the heavy-tailedness in
% the training signal, Figure~\ref{gold_proxy_trans} compares the empirical
% distributions of the original and transformed proxy rewards for the
% representative choice $\kappa=8$. Panels 1 and 3 show the
% distribution of the original reward $r(x,y)$ together with fitted normal
% density curves. In both cases, the untransformed reward is relatively
% concentrated and visually close to a light-tailed profile. Panels 2 and 4 show
% the corresponding transformed reward $\tilde r(x,y)$, overlaid with kernel
% density estimates.The transformed distribution becomes much more
% dispersed, with visibly heavier tails and more mass in the extremes. This
% supports that the power transform makes rare,
% large-magnitude reward realizations substantially more influential during
% training. \\

To confirm that the power transform amplifies the tail behavior of the training signal, Figure~\ref{gold_proxy_trans} compares the empirical distributions and QQ plots of the original and transformed proxy rewards for the representative choice \(\kappa=8\). The upper panels show that the original reward \(r(x,y)\) is relatively concentrated and visually close to a light-tailed profile, whereas the transformed reward \(\tilde r(x,y)\), overlaid with kernel density estimates, becomes much more dispersed with substantially more mass in the extremes. The lower panels provide the corresponding QQ plots against the normal reference distribution. The pronounced deviations in the tails after transformation further indicate heavier-than-normal tail behavior. These results support that the power transform makes rare, large-magnitude reward realizations substantially more influential during training.



%%%%%%%%%%%%%%%%%%%%%%%%%%%%%%%%%%%%%%%%%%%%%%%%%%%%%%%%%%%%%%%%%%%%%%%%%%%%%%%%%
%%%%%%%%%%%histogram of proxy reward before and after transformation%%%%%%%%%%%%%
%%%%%%%%%%%%%%%%%%%%%%%%%%%%%%%%%%%%%%%%%%%%%%%%%%%%%%%%%%%%%%%%%%%%%%%%%%%%%%%%%
\begin{figure}[t]
\centering
% \includegraphics[scale=0.295]{figure/proxy_rm_beta1_histograms.pdf}
% \includegraphics[scale=0.25]{figure/proxy_rm_beta1_qqplots.pdf}
\includegraphics[scale=0.3]{Aligned_figures/figure1.pdf}
\caption{\small Empirical distributions and QQ plots of proxy reward model scores before and after the power transformation for \(\kappa=8\). The upper panels show the empirical distributions of the original reward \(r(x,y)\) and the transformed reward \(\tilde r(x,y)=\operatorname{sign}(r(x,y))|r(x,y)|^{8}\) under \(\beta=1\) with \(\alpha=1\) (KL) and \(\alpha=2\) (R\'enyi). The original rewards are overlaid with fitted normal density curves, whereas the transformed rewards are overlaid with kernel density estimates. The lower panels report the corresponding QQ plots against the normal reference distribution. The transformation
preserves reward ordering while substantially amplifying the heavy-tailedness
of the training signal.}
\label{gold_proxy_trans}
\end{figure}

\noindent\textit{Robustness across heavy-tail constructions.}
The order-preserving power transform above provides a convenient knob (via $\kappa$) to dial tail-heaviness without changing preference rankings, but it is only one way to generate rare, high-impact reward outliers.
To verify that our conclusions are not an artifact of this particular transformation, we also study an alternative heavy-tailed construction based on adding centered Weibull noise to the reward signal prior to optimization; see Appendix~\ref{app:add_experiments}.
Across both constructions (and across sweeps of $\kappa$ and noise parameters), we observe the same qualitative pattern: KL-style control exhibits unstable reward--divergence behavior and a non-monotone proxy--gold relationship once rewards become heavy-tailed, whereas higher-order \Renyi regularization yields smoother reward--divergence curves and a more monotone proxy--gold coupling.\\






\paragraph{Reward--divergence behavior.} Figures~\ref{divergence_vs_reward11}--\ref{divergence_vs_reward1210} show the empirical relationship between reward and divergence under $\beta=1$.
Under KL regularization ($\alpha=1$), Figure~\ref{divergence_vs_reward11} shows a clear separation between the proxy reward (training-time signal) and the \emph{gold reward} (evaluation signal) once the proxy reward is heavy-tailed ($\kappa>1$). As the policy moves farther from the reference model, the proxy reward generally continues to increase, whereas the gold reward exhibits a pronounced \emph{rise-then-collapse} pattern: moderate divergence is initially beneficial, but additional divergence can eventually substantially degrade proxy-gold performance. This is a Goodhart-style signature of proxy over-optimization, in which continued improvement on the training-time reward no longer translates into improvement under the stronger evaluator.


This pattern is consistent with the mechanism in Section~\ref{se:whyKLfails}. Occasional extreme proxy rewards can dominate the policy-gradient estimate and induce abrupt, localized changes in the policy. Because the KL penalty averages discrepancies over the full action distribution, such tail-driven deviations need not incur a proportionate KL cost. The optimization can therefore drift toward policies that exploit proxy-specific artifacts while remaining relatively close to the reference policy in KL, which explains the collapse in gold reward in Figure~\ref{divergence_vs_reward11}.

Figure~\ref{divergence_vs_reward1210} shows that this behavior changes under R\'enyi regularization. For $\alpha>1$, higher-order R\'enyi divergences penalize large likelihood ratios $\pi(\cdot\mid x)/\pi_0(\cdot\mid x)$ more strongly, even when they arise on small subsets of actions. This makes concentrated, outlier-driven departures from the reference policy more expensive and thereby dampens the instability caused by heavy-tailed proxy rewards. Empirically, across $\alpha\in\{2,\ldots,10\}$, the gold reward is substantially more stable and typically increases before leveling off, rather than exhibiting the sharp collapse seen under KL. Taken together, Figures~\ref{divergence_vs_reward11} and~\ref{divergence_vs_reward1210} suggest that R\'enyi regularization provides a more reliable trust-region signal in the heavy-tailed regime.

\begin{figure}[t]
\centering 
\includegraphics[scale=0.39]{Aligned_figures/figure2.pdf}
% \begin{minipage}[t]{0.49\textwidth}
% \centering
% \includegraphics[scale=0.3]{figure_Power_Transformation_LLM/runs_e1/divergence_vs_reward_alpha1.pdf}
% \end{minipage}
% \begin{minipage}[t]{0.49\textwidth}
% \centering
% \includegraphics[scale=0.3]{figure_Power_Transformation_LLM/runs_e6/divergence_vs_reward_alpha1.pdf}
% \end{minipage}

% \begin{minipage}[t]{0.49\textwidth}
% \centering
% \includegraphics[scale=0.3]{figure_Power_Transformation_LLM/runs_e8/divergence_vs_reward_alpha1.pdf}
% \end{minipage}
% \begin{minipage}[t]{0.49\textwidth}
% \centering
% \includegraphics[scale=0.3]{figure_Power_Transformation_LLM/runs_e10/divergence_vs_reward_alpha1.pdf}
% \end{minipage}
\caption{ \small
Reward--divergence relationship for $\alpha=1$ (KL) with $\beta=1$ under different power exponents $\kappa$.
The four subplots correspond to $\kappa=1,6,8,10$, ordered from left to right and top to bottom.
As $\kappa$ increases, a clearer separation emerges between proxy and gold rewards: the proxy reward tends to keep increasing with divergence, whereas the gold reward exhibits a rise-then-collapse pattern, indicating over-optimization under heavy-tailed proxy rewards.
% In each subplot, the horizontal axis is the divergence from the reference policy and the vertical axis is the transformed reward; the dashed lines show the gold reward and the solid lines show the proxy reward. For visualization, both rewards are displayed after the same order-preserving, sign-preserving logarithmic transform.
}
\label{divergence_vs_reward11}
\end{figure}







\begin{figure}[t]
\centering 
\includegraphics[scale=0.24]{Aligned_figures/figure3.pdf}

% \begin{minipage}[t]{0.47\textwidth}
% \centering
% \includegraphics[scale=0.25]{figure_Power_Transformation_LLM/runs_e1/divergence_vs_reward_alpha2_to_10.pdf}
% \end{minipage}
% \begin{minipage}[t]{0.52\textwidth}
% \centering
% \includegraphics[scale=0.26]{figure_Power_Transformation_LLM/runs_e6/divergence_vs_reward_alpha2_to_10.pdf}
% \end{minipage}
% \begin{minipage}[t]{0.47\textwidth}
% \centering
% \includegraphics[scale=0.26]{figure_Power_Transformation_LLM/runs_e8/divergence_vs_reward_alpha2_to_10.pdf}
% \end{minipage}
% \begin{minipage}[t]{0.52\textwidth}
% \centering
% \includegraphics[scale=0.26]{figure_Power_Transformation_LLM/runs_e10/divergence_vs_reward_alpha2_to_10.pdf}
% \end{minipage}
\caption{\small
Reward--divergence relationship under R\'enyi regularization with $\beta=1$, for orders $\alpha\in\{2,\ldots,10\}$ and different power exponents $\kappa$.
The four subplots correspond to $\kappa=1,6,8,10$, ordered from left to right and top to bottom.
Across $\alpha\in\{2,\ldots,10\}$, the gold reward is substantially more stable and typically increases before leveling off, rather than exhibiting the sharp rise-then-collapse pattern observed under KL, suggesting that R\'enyi regularization provides a more reliable trust-region signal in the heavy-tailed regime.
}
\label{divergence_vs_reward1210}
\end{figure}








\paragraph{Proxy--proxy-gold coupling.}
Figure~\ref{proxy_vs_gold} examines how improvements in the proxy reward translate into improvements in the \emph{gold reward}. Under KL regularization ($\alpha=1$), this relationship becomes clearly non-monotone once the proxy reward is heavy-tailed ($\kappa>1$): the gold reward initially increases with the proxy reward, but beyond a certain point it begins to decline. Thus, continued optimization of the proxy reward can coincide with worse performance under the stronger evaluator, which manifests the over-optimization.

By contrast, under R\'enyi regularization ($\alpha\ge 2$), the proxy--gold curves remain monotone increasing and noticeably smoother across the observed range. It indicates that gains in the proxy reward are more consistently aligned with gains in the gold reward.
Overall, these results suggest that R\'enyi regularization yields a more reliable coupling between proxy and proxy-gold objectives, and therefore a more reliable trust-region signal than KL in the heavy-tailed regime.\\



\begin{figure}[t!]
\centering 
\includegraphics[scale=0.35]{Aligned_figures/figure4.pdf}

% \begin{minipage}[t]{0.47\textwidth}
% \centering
% \includegraphics[scale=0.33285]{figure_Power_Transformation_LLM/runs_e1/proxy_vs_gold_reward.pdf}
% \end{minipage}
% \begin{minipage}[t]{0.52\textwidth}
% \centering
% \includegraphics[scale=0.33285]{figure_Power_Transformation_LLM/runs_e6/proxy_vs_gold_reward.pdf}
% \end{minipage}
% \begin{minipage}[t]{0.47\textwidth}
% \centering
% \includegraphics[scale=0.32685]{figure_Power_Transformation_LLM/runs_e8/proxy_vs_gold_reward.pdf}
% \end{minipage}
% \begin{minipage}[t]{0.52\textwidth}
% \centering
% \includegraphics[scale=0.32685]{figure_Power_Transformation_LLM/runs_e10/proxy_vs_gold_reward.pdf}
% \end{minipage}
\caption{\small
Proxy--proxy-gold reward relationship under $\beta=1$, for R\'enyi orders $\alpha\in\{1,\ldots,10\}$ and power exponents $\kappa\in\{1,6,8,10\}$.
The four subplots correspond to $\kappa=1,6,8,10$, ordered from left to right and top to bottom.
As $\kappa$ increases, the KL case ($\alpha=1$) becomes clearly non-monotone, with the gold reward eventually declining as the proxy reward continues to improve, whereas the curves for $\alpha\ge 2$ remain smoother and largely monotone increasing, indicating a more reliable coupling between proxy and gold objectives.
}
\label{proxy_vs_gold}
\end{figure}


% \caption{
% Proxy--proxy-gold reward relationship under $\beta=1$, for R\'enyi orders $\alpha\in\{1,\ldots,10\}$ and power exponents $\kappa\in\{1,6,8,10\}$.
% The four subplots correspond to $\kappa=1,6,8,10$, ordered from left to right and top to bottom.
% In each subplot, the horizontal axis shows the transformed proxy reward and the vertical axis shows the transformed gold reward.
% Different colors correspond to different R\'enyi orders $\alpha=1,\ldots,10$.
% For visualization, both rewards are displayed after the same order-preserving, sign-preserving logarithmic transform.
% }






\paragraph{Additional robustness checks.}
Appendix~\ref{app:add_experiments} provides additional diagnostics and robustness checks, including tail plots for the power transformation and the full set of results for the additive Weibull-noise construction. Having established the same qualitative pattern under controlled heavy-tailed perturbations, we next examine a more adversarial source of heavy tails: token-space attacks that directly exploit reward-model vulnerabilities.


\subsection{RLHF under Token-Space Reward Attacks}\label{sec:tompa-experiments}\label{app:tompa}\label{app:toma}

The controlled experiments above isolate the effect of tail behavior by modifying
the reward signal while keeping the optimization pipeline fixed. We now consider
a less stylized experiment in which heavy-tailed rewards emerge from
reward-model exploitation itself. \citet{zhang2026beyond} introduced TOMPA and
identified a striking failure mode of reward modeling in RLHF: reward-model
vulnerabilities are not confined to semantic manipulation. Since an attacker
can search directly in token space for non-semantic adversarial patterns that
induce abnormally large reward values. The resulting high-reward outputs may
degenerate into nonsensical text while still receiving favorable evaluations
from the reward model. This phenomenon is closely aligned with the theory above:
rare token sequences can trigger disproportionately large rewards, and those
extreme values can dominate policy updates even when most rollouts receive
moderate rewards. Thus, we first inspect the tail behavior of the reward
distribution and then ask whether gains in the optimized target reward remain
aligned with an external reward diagnostic.

Building on the RLHF setup in Section~\ref{sec:rlhf-experiments}, we compare KL
regularization with R\'enyi regularization at $\alpha\in\{2,4\}$, using
$\beta=1$. Following TOMPA, we use Skywork-Reward-V2-Llama-3.1-8B
\citep{liu2025skywork} as the frozen target reward model and Qwen3-1.7B
\citep{qwen2025technical} as the attack policy. We train on 5{,}000 prompts
sampled from WildChat \citep{zhao2024wildchat,deng2024wildvis} using GRPO with
eight rollouts per prompt. For reward computation, only the generated response
token IDs are passed directly into the target reward model; the prompt is
excluded, and token IDs exceeding the target vocabulary are clipped to its
largest valid index. We update only rank-32 LoRA parameters of the policy, while
the policy backbone and reward model remain frozen. The maximum prompt and
response lengths are 2{,}048 and 512 tokens, respectively, and training uses a
learning rate of $10^{-6}$. As an online diagnostic, we periodically decode the
mapped response IDs with the target-model tokenizer and score the resulting
answer-only text using the frozen OpenAssistant reward model
\texttt{reward-model-deberta-v3-large-v2} \citep{openassistant2023rewarddeberta}.
We treat this auxiliary diagnostic score as the \emph{gold reward} in this
subsection. It is recorded only for monitoring and does not affect policy
optimization.

\begin{figure}[t]
\centering
\includegraphics[scale=0.25]{figure_token_space_attacks/kl_reward_emgf.pdf}
% \includegraphics[scale=0.15]{figure_token_space_attacks/kl_reward_qq.pdf}

\includegraphics[scale=0.25]{figure_token_space_attacks/renyi_a2_reward_emgf.pdf}
% \includegraphics[scale=0.149]{figure_token_space_attacks/renyi_a2_reward_qq.pdf}

\includegraphics[scale=0.25]{figure_token_space_attacks/renyi_a4_reward_emgf.pdf}
% \includegraphics[scale=0.149]{figure_token_space_attacks/renyi_a4_reward_qq.pdf}
\caption{\small Tail behavior of reward distributions under TOMPA token-space attacks with KL and R\'enyi regularization. The top row shows KL regularization, while the middle and bottom rows correspond to R\'enyi regularization with $\alpha=2$ and $\alpha=4$, respectively. Left panels plot the empirical MGF $M(t)$ against the standard normal benchmark. KL regularization exhibits stronger right-tail inflation and upper-tail deviations, while R\'enyi regularization attenuates these tail effects, especially when $\alpha=4$.}
\label{fig:tompa_tail_diagnostics}
\end{figure}

Figure~\ref{fig:tompa_tail_diagnostics} shows that KL regularization produces
more pronounced heavy-tailed reward behavior under token-space attack: the
empirical MGF grows faster than the normal benchmark. In contrast, R\'enyi regularization reduces the
upper-tail deviation, with a stronger attenuation at $\alpha=4$. This suggests
that R\'enyi regularization can suppress rare but abnormally high reward values
induced by token-space exploitation
\begin{figure}[t]
\centering
\includegraphics[width=0.98\textwidth]{figure_token_space_attacks/reward_vs_divergence.pdf}
\vspace{0.8em}
\includegraphics[width=0.98\textwidth]{figure_token_space_attacks/target_vs_gold_reward.pdf}
\caption{\small Reward--divergence diagnostics under TOMPA token-space attacks with KL and R\'enyi regularization. The upper panel reports the relationship between reward and divergence for $\alpha=1$ (KL), $\alpha=2,4$ (R\'enyi), where points denote observed rewards and curves denote smoothed trends. The lower panel compares the directly optimized target reward with the auxiliary gold reward, obtained by decoding the mapped response IDs and scoring the resulting answer-only text with an external frozen reward model. Under KL regularization, the target reward increases as divergence grows, whereas the gold reward first increases and then decreases. Under R\'enyi regularization, both rewards tend to increase together and gradually stabilize.}
\label{fig:reward_divergence_diagnostics}
\end{figure}
Figure~\ref{fig:reward_divergence_diagnostics} further explains how this tail
behavior translates into reward hacking. Under KL regularization, the target
reward continues to rise as divergence increases, but the gold reward first
increases and then declines. The lower panel gives a direct view of this
decoupling: in the high-target-reward region, larger optimized rewards no longer
correspond to better decoded outputs under the external diagnostic. The policy
therefore appears to exploit token-space vulnerabilities of the target reward
model rather than produce responses with improved reward behavior.

By contrast, R\'enyi regularization yields a more stable relationship between
the optimized target reward and the gold reward. For $\alpha=2$ and $\alpha=4$,
both rewards tend to increase together as divergence grows, and the gold reward
approaches a plateau rather than declining in the high-target-reward region.
Combining Figures~\ref{fig:tompa_tail_diagnostics}
and~\ref{fig:reward_divergence_diagnostics}, we find that KL regularization
leaves two related vulnerabilities in this adversarial setting: heavier-tailed
reward responses and a mismatch between the optimized target reward and the
external gold-reward diagnostic. R\'enyi regularization mitigates both effects,
supporting the broader conclusion that tail-aware divergence control is useful
not only for controlled heavy-tailed perturbations, but also for reward-model
exploitation that creates heavy tails endogenously.


\section{Conclusions and Discussions}\label{sec:conclusion}

This paper develops a tail-aware information--theoretic framework for
data-dependent selection and adaptive suprema under sub-Weibull tails. The
starting point is the observation that KL mutual-information bounds,
while powerful in light-tailed settings, can be ineffective when MGFs do not
exist and rare extreme events dominate selection effects. To address this, we
introduce a decorrelation lemma tailored to sub-Weibull tails: it controls
expectations under a changed measure through a shifted-log $f_\theta$-divergence
and relates this quantity to \Renyi divergence. This route avoids MGF arguments
and works directly with sub-Weibull Orlicz norms.

On the empirical-process side, we extend maximal and chaining tools to genuinely
heavy-tailed regimes. We establish a sharp maximal inequality for heavy-tailed
sub-Weibull random variables and a Dudley-type entropy bound for sub-Weibull
processes. The resulting rates make the role of the tail index explicit: maxima
scale as $(\log n)^{1/\theta}$, while metric entropy enters the Dudley integral
with exponent $1/\theta$.
%These results complement concentration inequalities
%for sums by controlling maxima and suprema, the objects most directly involved
%in data-adaptive selection and empirical-process arguments.

Combining the decorrelation and empirical-process ingredients, we obtain
tail-adaptive information bounds for data-dependent selectors. The single-scale
bound replaces KL mutual information with shifted-log $f_\theta$-mutual
information and its \Renyi refinements. The multiscale result then lifts this
control to a Dudley-type information inequality through the discretized variables
$[W]_k$. This chaining formulation is especially useful when a single-scale
information term is infinite or vacuous, but the information revealed at each
resolution remains controlled.

We apply the same tail-aware viewpoint to RLHF under heavy-tailed and
adversarial rewards. The theory shows why KL-constrained RLHF can suffer
catastrophic Goodhart under heavy-tailed proxy rewards, and why
\Renyi-regularized RLHF controls reward inflation at a fixed divergence budget.
We also analyze best-of-$n$ policies and show that their reward growth is
controlled through \Renyi divergence and the sub-Weibull maximal rate. The
numerical studies support this picture in two settings: controlled heavy-tail
perturbations of reward signals and TOMPA-style token-space reward attacks. In
both cases, higher-order \Renyi regularization attenuates tail inflation and
improves the coupling between proxy and gold reward diagnostics relative to KL.

Several limitations and directions remain. First, the sharpness of the
\Renyi-based constants and exponents should be better understood, especially in
structured or high-dimensional settings where localization may improve the
bounds. %Second, the analysis is stated under sub-Weibull Orlicz control and
%mostly i.i.d.\ sampling assumptions; extending the framework to dependent data,
%reinforcement-learning trajectories, and adaptive multi-epoch training dynamics
%is an important next step. 
Second, while this work focuses on shifted-log
$f$-divergences and \Renyi divergences with $\alpha>1$, other information
measures, such as $\alpha<1$ \Renyi divergences, Wasserstein distances, or hybrid
integral probability metrics, may provide complementary forms of tail control. Third, adaptively estimating \Renyi order $\alpha$ is an important next step. Similarly, \cite{ferrari2010maximum} choose $\alpha$ by minimizing an estimated asymptotic mean squared error of the maximum $L\alpha$-likelihood estimator.


From a practical standpoint, the results suggest moving beyond KL penalties when
rare high-reward events can dominate optimization. For RLHF, \Renyi-type
regularization offers a way to penalize concentrated tail exploitation without
simply forbidding policy movement. The token-space attack experiment further
suggests that such control is relevant not only for synthetic heavy-tailed
perturbations, but also when reward-model vulnerabilities create heavy tails
endogenously. We hope these tools contribute to a broader tail-adaptive
information--theoretic theory for modern learning and alignment systems. 

\section*{Acknowledgements}
\noindent H.Z. is supported in part by Beijing Advanced Innovation Center for Future Blockchain and Privacy Computing and the Fundamental Research Funds for the Central Universities (No. KG16447001). W.T. is supported in part by the Funds of the Natural science Foundation of Hangzhou (No.~2025SZRJJ1388) and the Beijing Outstanding Young Scientist Program (No.~JWZQ20240101027).  Q.S. is supported in part by MBZUAI, the Natural Sciences and Engineering Research Council of Canada (Grant RGPIN-2018-06484), computing resources provided by the Digital Research Alliance of Canada.
%\vskip 0.2in
%\begin{footnotesize}
%\newpage

\bibliography{ref}
%\end{footnotesize}

%\begin{appendix}
\newpage
\appendix
%\section*{Appendix A.}
%\label{app:theorem}
%\addcontentsline{toc}{section}
% Note: in this sample, the section number is hard-coded in. Following
% proper LaTeX conventions, it should properly be coded as a reference:



\section*{Appendix}

\startcontents[sections] % 开始为``章节''层面积累目录条目
\printcontents[sections]{}{1}{} % 在这里打印出刚积累的附录目录
% 第二个参数{}是前缀，第三个参数1设置目录深度（1=到section）


%%%%%%%%%%%%%%%%%%%%%%%%%%%%%%%%%%%%%%%%%%%%%%%%
%%%%%%%%%%%%%%%%%%Preliminaries%%%%%%%%%%%%%%%%%%
%%%%%%%%%%%%%%%%%%%%%%%%%%%%%%%%%%%%%%%%%%%%%%%%

\section{Preliminaries}\label{app:prior}

\subsection{Preliminaries on sub-Weibull random variables}\label{app:prelim}



%\subsection{Orlicz norm and sub-Weibull random variables}

We review the definition of Orlicz norm and an equivalent definition for sub-Weibull random variables. 
%\paragraph{Orlicz norm.}

\begin{definition}[Orlicz norm, \cite{Aad2023}]
For a random variable $X$, the Orlicz norm associated with a Young function $\psi$ is defined as
\[
\|X\|_{\psi} := \inf \left\{ K > 0 : \mathbb{E}\,\psi\!\left(\frac{|X|}{K}\right) \le 1 \right\},
\]
where $\psi : [0,\infty) \to [0,\infty)$ is a Young function, i.e., a convex function satisfying $\psi(0)=0$ and $\lim_{x \to \infty} \psi(x) = \infty$. 
\end{definition}


For example, the sub-Weibull Orlicz norm is defined as
\[
\|X\|_{\psi_\theta}:=\inf\left\{K>0:\ \mathbb{E}\exp\!\left(\left|\frac{X}{K}\right|^\theta\right)\le 2\right\}.
\]
The $\|\cdot\|_{\psi_\theta}$ is a norm when $\theta\ge 1$. For $0<\theta<1$, it is only a quasi-norm: the triangle inequality fails, but a weak triangle inequality holds~\cite[Lemma A.3]{gotze2021concentration}. 



%%%%%%%%%%%%%%%%%%%%%%%%%%%%%%%%%%%%%% 
%%%%%%Sub-Weibull distributions%%%%%%%
%%%%%%%%%%%%%%%%%%%%%%%%%%%%%%%%%%%%% 
As mentioned in the introduction, sub-Weibull random variables can also be equivalently characterized via their tail probabilities. Below we give a definition equivalent to the one based on the Orlicz norm.

\begin{definition}[Sub-Weibull random variables]\label{def:subweibull}
A random variable $X$ is called \emph{sub-Weibull} with tail parameter $\theta>0$, denoted $X \sim \subw(\theta)$, if there exist positive constants $a,\, b>0$ such that
\[
\mathbb{P}(|X| \ge x) \le a \exp \bigl\{-(x/b)^{\theta}\bigr\}, \qquad x>0.
\]
\end{definition}

In other words, sub-Weibull distributions have tails no heavier than those of a Weibull random variable with the same parameter $\theta$. When $\theta < 1$, $X$ is heavy-tailed, whereas $\theta \ge 1$ corresponds to light-tailed distributions, including sub-Gaussian ($\theta=2$) and sub-exponential ($\theta=1$) cases.  
The next section reviews key prior results on maximal inequalities for light-tailed sub-Weibull random variables (sub-Weibull with $\theta\geq 1$, \cite{Aad2023}), and information-theoretic generalization bounds for sub-Gaussian random variables \citep{Russo20,chu2023unified}.

\subsection{Maximum inequalities for light-tailed sub-Weibull random variables}\label{se:ap-max}

The following maximal inequality for light-tailed sub-Weibull random variables with $\theta\geq 1$ is due to \cite{Aad2023}. 

\begin{proposition}[Maximal inequality for light-tailed sub-Weibull random variables]\label{prp-MaximalSW}
For $\theta \ge 1$, let $\{X_{i}\}_{i=1}^n$ be random variables with $\max_{1 \le i \le n}\|X_i\|_{\psi_\theta} < \infty$ and $\psi_\theta(x)=e^{x^\theta}-1$,
\begin{equation}\label{eq:MaximalSW1}
\E \Bigl( \max_{1 \le i \le n} |X_i| \Bigr)
\;\le\; \psi^{-1}_{\theta}(n) \max_{1 \le i \le n}\|X_i\|_{\psi_{\theta}}
= (\log (1+n))^{1 / \theta}\, \max_{1 \le i \le n}\|X_i\|_{\psi_{\theta}}.
\end{equation}
\end{proposition}



The proof of Proposition~\ref{prp-MaximalSW} relies on Jensen's inequality, which requires the convexity of $\psi_\theta$. This property holds only when $\theta \ge 1$ and fails for $\theta<1$. Note that the sub-Gaussian case corresponds to $\theta = 2$. 

More generally, \citet{Russo20} extended maximal inequalities in the sub-Gaussian setting by relating them to mutual information. Under the assumption that $\{X_{i}\}_{i=1}^n$ are i.i.d. $\sigma^2$-sub-Gaussian, they established:
\begin{equation}\label{eq:RZ}
    \E \bigl[ X_W \bigr] \leq \sqrt{2\sigma^2\, I(W;X_1, \ldots, X_n)},
\end{equation}
where the algorithm output is a random index $W = W(X_1, \ldots, X_n)$. When $W$ is taken to be the maximum selector in $[n]$, the mutual information bound yields the classical sub-Gaussian maximal inequality
\begin{equation}\label{eq:RZ1}
\mathbb{E}\bigl[\max_{ i \in [n]} X_i\bigr] \leq \sqrt{2 \log n},
\end{equation}
by using the fact that $I(W;X_1, \ldots, X_n)\le \log n$. Hence, \eqref{eq:RZ} implies \eqref{eq:RZ1} as a special case, so \eqref{eq:RZ} is more general than \eqref{eq:RZ1}.

When the index set is infinite or even uncountable, it is natural to generalize maximal inequalities to upper bounds on the supremum of empirical or stochastic processes. Dudley's inequality (Theorem 5.24 in \cite{vanhandle2016probability}) states that
\[
\mathbb{E}\left[\sup _{t \in T} X_t\right] 
\;\leq\; 6 \sum_{k \in \mathbb{Z}} 2^{-k} \sqrt{\log N\left(T, d, 2^{-k}\right)}
\]
if $\left\{X_t\right\}_{t \in T}$ is a separable 1-sub-Gaussian process of zero mean on the bounded metric space $(T, d)$, i.e.,
\[
\E \exp\bigl(\lambda(X_t-X_s)\bigr) \leq \exp\bigl(\lambda^2 d(t, s)^2 / 2\bigr)
\quad \text{for all}~t, s \in T,\; \lambda \geq 0.
\]



\subsection{Information-theoretic generalization bounds}


Building upon \citet{Russo20}, \citet{asadi2018chaining} introduced a chaining method to derive information-theoretic generalization bounds, which can be viewed as a refinement of Dudley's inequality through mutual information.


\begin{proposition}[\citet{asadi2018chaining}]\label{Chaining MI Random Process}
Assume that $\{X_t\}_{t\in T}$ is a separable sub-Gaussian process of zero mean on a bounded metric space $(T,d)$, and let $k_1(T)$ be an integer such that $2^{-\left(k_1(T)-1\right)} \geq \operatorname{diam}(T)$. Let $\{\mathcal{P}_k\}_{k=k_1(T)}^{\infty}$ be an increasing sequence of partitions of $T$, where for each $k\geq k_1(T)$, $\mathcal{P}_k$ is a $2^{-k}$-partition\footnote{Recall the definition of partition in Definition \ref{def:part}.} of $(T,d)$. Then:
\begin{enumerate}
\item
$\E X_W \leq 3\sqrt{2}\sum_{k=k_1(T)}^{\infty}2^{-k}\sqrt{I([W]_k;X_T)}.$
\item For any arbitrary $t_0\in T$,
$
\E\bigl\lvert X_W-X_{t_0}\bigr\rvert \leq 3\sqrt{2}\sum_{k=k_1(T)}^{\infty}2^{-k}\sqrt{I([W]_k;X_T)+\log 2}.
$
\end{enumerate}
\end{proposition}

When $W$ is taken to be the maximum selector, Proposition~\ref{Chaining MI Random Process} yields an infinite-series bound that refines Dudley's inequality, since $\mathbb{E}[X_W] \leq \mathbb{E}\bigl[\sup _{t \in T} X_t\bigr]$.

However, standard MGF-based techniques for deriving maximal inequalities and mutual information bounds are ineffective for heavy-tailed sub-Weibull distributions. This motivates the development of new tools to obtain sharp information-theoretic generalization bounds in the heavy-tailed regime. 


%%%%%%%%%%%%%%%%%%%%%%%%%%%%%%%%%%%%%
%%%%%%%%%%%%%%Proofs%%%%%%%%%%%%%%%%%
%%%%%%%%%%%%%%%%%%%%%%%%%%%%%%%%%%%%%

\section{Proofs for Section \ref{sec:ep}}\label{se: proof2}





\subsection{Proof of Lemma \ref{pp-MaximalSW}}
\begin{proof}
The function $\psi_\theta$ is increasing on $[0,\infty)$ and convex on $[x_\theta,\infty)$.
Indeed,
\[
\psi_\theta'(x)=\theta x^{\theta-1}e^{x^\theta},
\qquad
\psi_\theta''(x)=\theta x^{\theta-2}e^{x^\theta}\bigl(\theta x^\theta-(1-\theta)\bigr),
\]
so $\psi_\theta''(x)\ge 0$ for all $x\ge x_\theta=:\paren{\frac{1-\theta}{\theta}}^{1/\theta}$. 
Let $0<\theta<1$ and define $
g_\theta(x):=\paren{\psi_\theta(x)-\psi_\theta(x_\theta)}_+$.
Then $g_\theta$ is convex and nondecreasing on $[0,\infty)$. 

Let $K:=\max_{1\le i\le n}\norm{X_i}_{\psi_\theta}$ and $M:=\max_{1\le i\le n}\frac{\abs{X_i}}{K}$.
Then, the elementary bound
$\max_{1\le i\le n} a_i \le \sum_{i=1}^n a_i~(a_i\ge 0)$ gives
\[
\E \psi_\theta(M)\le \sum_{i=1}^n \E \psi_\theta\!\paren{\frac{\abs{X_i}}{K}}\le n.
\]
By Jensen's inequality,
\[
\paren{\psi_\theta(\E M)-\psi_\theta(x_\theta)}_+= g_\theta(\E M)\le \E g_\theta(M)\le \E\psi_\theta(M)\le n.
\]
Hence,
$\psi_\theta(\E M)\le \psi_\theta(x_\theta)+n$. Since $\psi_\theta$ is increasing,
\[
\E M=\E \paren{\max_{1\le i\le n}\frac{\abs{X_i}}{K}}\le \psi_\theta^{-1}\!\paren{\psi_\theta(x_\theta)+n},
\]
and multiplying by $K:=\max_{1\le i\le n}\norm{X_i}_{\psi_\theta}$ yields \eqref{eq:MaximalSW}. The big-$O$ rate is from $\psi_\theta^{-1}$.
\end{proof}


\subsection{Optimality and comparison with classical bounds}\label{app:sec2}

The following result shows that the bound in Lemma \ref{pp-MaximalSW} is optimal in order: the rate $\log^{1/\theta} n$ cannot be improved in general.   

\begin{proposition}\label{prop:lowerbound}
Let $\{X_i\}_{i=1}^n$ be i.i.d. Weibull variables $\mathbb{P}(X\ge x)=\exp(-(x/b)^\theta)$. Then, 
\[
\E ( \max_{i\in [n]} X_i ) \gtrsim \log^{1 / \theta}(n).
\]
\end{proposition}
\begin{proof}
The upper bound follows from the proof of Lemma \ref{pp-MaximalSW}. For the lower bound:
\[
\begin{aligned}
\E ( \max_{i\in [n]} X_i ) & =\int_0^{\infty} \mathbb{P}(\max_{i\in [n]} X_i>x) \, d x= \int_0^\infty \bigl[1-(1-\exp(-(x/b)^\theta))^n\bigr] d x\\
&\geq \int_0^{b\log(n)^{1 / \theta}} \left\{1-(1-\exp(-(x/b)^\theta))^n\right\} d x\\
&\geq \int_0^{b\log(n)^{1 / \theta}} \left\{1-(1-\frac{1}{n})^n\right\}d x\geq b\log^{1 / \theta}(n) \left(1-\frac{1}{e}\right),
\end{aligned}
\]
where the third line uses monotonicity of function $t\mapsto t^n$ and the last line uses $(1-\frac{1}{n})^n\leq \frac{1}{e}$.
\end{proof}



%%%%%%%%%%%%%%%%%%%%%%%%%%%%%%%%%%%%%%%%%%%%%%%%%%%%%%%%%%%%%%%%%%%%%%%%%%%%%%%%%
%%%%%%%%%%%%%%%%%%%%%%%%Comparison  with Classical Bounds%%%%%%%%%%%%%%%%%%%%%%%%
%%%%%%%%%%%%%%%%%%%%%%%%%%%%%%%%%%%%%%%%%%%%%%%%%%%%%%%%%%%%%%%%%%%%%%%%%%%%%%%%%
\subsection{Comparison with norm-to-moment bounds}\label{app:max_comparison}


A standard alternative to bounding $\mathbb{E}\max_{i\in [n]} |X_i|$ is to
first control the Orlicz norm of the maximum and then convert that norm bound
into a moment bound. For instance, \citet[(S2.2)]{mendes2023concentration}
showed that
\[
\left\|\max_{i\in [n]}|X_i|\right\|_{\psi_\theta}
\le c_1 \psi_\theta^{-1}(2n)\,\max_{i\in [n]}\|X_i\|_{\psi_\theta},
\qquad\text{where}~~c_1=\frac{2}{\log(1.5)}.
\]
using \cite[Lemma~2.2.2]{Aad2023}. Combining this with the first-moment bound
for the sub-Weibull norm \cite[Corollary~3]{zhang2022sharper},
$
\mathbb{E}|X|\le 2\|X\|_{\psi_\theta}\Gamma\!\left(\frac{1}{\theta}+1\right),
$
yields the bound
\[
\mathbb{E}\Bigl[\max_{i\in [n]} |X_i|\Bigr]
\le 2c_1\,\psi_\theta^{-1}(2n)\,\Gamma\!\left(\frac{1}{\theta}+1\right)
\max_{i\in [n]}\|X_i\|_{\psi_\theta}.
\]
In contrast, Lemma~\ref{pp-MaximalSW} controls the expectation directly and
avoids the additional multiplicative factor
$\Gamma(1/\theta+1)$ introduced by the norm-to-moment conversion. The ratio
between the two prefactors is
\begin{equation}\label{eq:sharplow}
    R_{n,\theta}=\frac{4/ \log (1.5) \psi_\theta^{-1}(2 n)\Gamma\!\left(\frac{1}{\theta}+1\right)}{\psi_\theta^{-1}\bigl(\psi_\theta(x_\theta)+n\bigr)
}\ge \frac{4}{\log(1.5)} \Gamma\left(\frac{1}{\theta} + 1\right) \approx 9.865 \, \Gamma\left(\frac{1}{\theta} + 1\right).
\end{equation}
Since $\Gamma(1/\theta+1)$ grows rapidly as $\theta\downarrow 0$, the
norm-to-moment route can be substantially looser in the heavy-tailed regime,
whereas Lemma~\ref{pp-MaximalSW} retains the sharp logarithmic scaling.






To verify \eqref{eq:sharplow}, note that
$\psi_\theta(x)=e^{x^\theta}-1$ has inverse
$\psi_\theta^{-1}(y)=\bigl[\log(y+1)\bigr]^{1/\theta}$. Thus,
\begin{equation*}
    R_{n,\theta}= \frac{\frac{4}{\log(1.5)} \psi_\theta^{-1}(2n) \Gamma\left(\frac{1}{\theta}+1\right)}{\psi_\theta^{-1}\bigl(\psi_\theta(x_\theta)+n\bigr)}.
\end{equation*}
We substitute $\psi_\theta(x_\theta) = e^{x_\theta^\theta} - 1=e^{\frac{1-\theta}{\theta}} - 1$, where $x_\theta^\theta = \frac{1-\theta}{\theta}$.
Substituting the inverse functions gives
\begin{equation*}
    R_{n,\theta} = \frac{4}{\log(1.5)} \Gamma\left(\frac{1}{\theta}+1\right) \left[ \frac{\log(2n+1)}{\log(n + e^{\frac{1-\theta}{\theta}})} \right]^{1/\theta}.
\end{equation*}
As $n\to\infty$,
\begin{equation*}
    \lim_{n \to \infty} \frac{\log(2n+1)}{\log(n + e^{\frac{1-\theta}{\theta}})} = \lim_{n \to \infty} \frac{\log n + \log(2 + 1/n)}{\log n + \log(1 + e^{\frac{1-\theta}{\theta}}/n)} = 1.
\end{equation*}
Therefore, the ratio converges to the constant
$\frac{4}{\log(1.5)} \Gamma\left(\frac{1}{\theta}+1\right)$. Moreover, for any
$n \ge e^{\frac{1-\theta}{\theta}}$, we have
$\frac{\log(2n+1)}{\log(n + e^{\frac{1-\theta}{\theta}})} \geq 1$. Hence,
\begin{equation*}
    R_{n,\theta} \ge \frac{4}{\log(1.5)} \Gamma\left(\frac{1}{\theta}+1\right).
\end{equation*}


The next example shows that taking $W=\arg\max_{i\in[n]}|X_i|$ in Theorem \ref{infob} recovers the maximal growth rate $\mathbb{E}[\max_i|X_i|]\lesssim(\log n)^{1/\theta}$.



\begin{example}\label{Max}
Let $S= \{X_1, \ldots, X_n\}$ be a vector of $n$ i.i.d. continuous random
variables with a common density. Let
$
W = \argmax_{i \in [n]} |X_i|$. Then
\[
I(W;S)=\lim_{\alpha \to 1} I_\alpha (W;S) =\log n,
\]
and for every $\alpha \in (0, \infty) \setminus \{1\}$, the R\'enyi-$\alpha$
mutual information is
\[
    I_\alpha(W; S) = \log n.
\]
\end{example}
\begin{proof}
By continuity, ties in $|X_1|,\ldots,|X_n|$ occur with probability $0$.
Hence the maximizer
$W=\arg\max_{i\in [n]} |X_i|$
is almost surely unique and takes values in $\{1,\ldots,n\}$. By the
exchangeability of $\{X_1,\ldots,X_n\}$, we have for every $i$,
\[
\mathbb{P}(W=i)=\frac{1}{n}.
\]
Moreover, conditional on $S=s=(x_1,\ldots,x_n)$, the value of $W$ is deterministic:
$\mathbb{P}(W=i\mid S=s)=\mathbf{1}\{i=\arg\max_{j\in [n]} |x_j|\}$, which
is $\{0,1\}$-valued and
$\sum_{i=1}^n \mathbb{P}(W=i\mid S)=1$.

\paragraph{Kullback--Leibler mutual information:}
Write $P_{W\mid S}$ for the conditional law of $W$ given $S$ and $P_W$ for the marginal law. Then
\[
\begin{aligned}
I(W;S)&=\E\Bigl[ \DD_{\mathrm{KL}}\bigl(P_{W\mid S}\,\|\,P_W\bigr)\Bigr]
      =\E\biggl[\sum_{i=1}^n \mathbb{P}(W=i\mid S)\log\frac{\mathbb{P}(W=i\mid S)}{\mathbb{P}(W=i)}\biggr]\\
      &=\E\biggl[\sum_{i=1}^n \mathbb{P}(W=i\mid S)\log\frac{1}{1/n}\biggr]=\log n.
\end{aligned}
\]
\paragraph{R\'enyi mutual information.}
For $\alpha\neq 1$, 
\[
I_\alpha(W;S)=\DD_\alpha(P_{W,S}\,\|\,P_W\otimes P_S)
=\frac{1}{\alpha-1}\log \int \Bigl(\frac{dP_{W,S}}{d(P_W\otimes P_S)}\Bigr)^\alpha \, d(P_W\otimes P_S).
\]
Because $W$ is discrete and $S$ has density, the Radon--Nikodym derivative exists and equals
\[
\frac{dP_{W,S}}{d(P_W\otimes P_S)}(i,s)=\frac{\mathbb{P}(W=i\mid S=s)}{\mathbb{P}(W=i)}
= n\,\mathbb{P}(W=i\mid S=s).
\]
Hence, using $\mathbb{P}(W=i\mid S=s)\in\{0,1\}$ and $\sum_{i=1}^n \mathbb{P}(W=i\mid S=s)=1$,
\begin{align}\label{eq:p-moment}
\int \Bigl(\frac{dP_{W,S}}{d(P_W\otimes P_S)}\Bigr)^\alpha d(P_W\otimes P_S)
&=\E\Bigl[\sum_{i=1}^n \mathbb{P}(W=i)\bigl(n\mathbb{P}(W=i\mid S)\bigr)^\alpha\Bigr]\nonumber\\
&=\E\Bigl[\frac{1}{n}\sum_{i=1}^n n^\alpha \mathbb{P}(W=i\mid S)\Bigr]=\E\bigl[n^{\alpha-1}\bigr]
= n^{\alpha-1}.
\end{align}
Therefore, $I_\alpha(W;S)=\frac{1}{\alpha-1}\log\bigl(n^{\alpha-1}\bigr)=\log n$ for $\alpha\neq 1$. Taking $\alpha\to 1$ yields $I(W;S)=\lim_{\alpha\to 1} I_\alpha(W;S)=\log n$.
\end{proof}









\subsection{Proof of Theorem \ref{Dudley}}
\begin{proof}
The proof proceeds in four steps.\\
\medskip
\noindent\textbf{Step 1: Discretization via $\varepsilon$-nets.}
If $\abs{T}=1$, the result is trivial. Assume $\abs{T}\ge 2$. For $k\ge 0$, set
$\varepsilon_k:=2^{-k}e(T)$. Since $T$ is finite, there exists
\[
\Delta:=\min\{d(s,t):s,t\in T,\ s\neq t\}>0.
\]
Choose $K$ so large that $\varepsilon_K<\Delta/2$. Then every $\varepsilon_K$-ball contains at most one point of $T$, so every $\varepsilon_K$-net is all of $T$ and $T_K=T$.

Fix $k\in\{1,\dots,K\}$ and $u\in T_k$. Since $T_{k-1}$ is an
$\varepsilon_{k-1}$-net of $T$ and $u\in T_k\subseteq T$, there exists
$p_k(u)\in T_{k-1}$ such that
\[
d\bigl(u,p_k(u)\bigr)\le \varepsilon_{k-1}.
\]
This proves the existence of the parent map $p_k$.

\medskip
\noindent\textbf{Step 2: Apply the maximal inequality.}
Next, for each $k$, the set of possible $k$th increments is contained in
$
\bigl\{X_u-X_{p_k(u)}:u\in T_k\bigr\}$. For each $k\ge 1$ and $u\in T_k$, set $Y_{k,u}:=X_u-X_{p_k(u)}$. Then
\[
\norm{Y_{k,u}}_{\psi_\theta}\le C\,d\paren{u,p_k(u)}\le 2C\varepsilon_k.
\]
Moreover,
$
\sup_{t\in T}\abs{X_{\pi_k(t)}-X_{\pi_{k-1}(t)}}
\le
\max_{u\in T_k}\abs{Y_{k,u}}$. Applying Lemma~\ref{pp-MaximalSW} at scale $k$ gives
\[
\E\sup_{t\in T}\abs{X_{\pi_k(t)}-X_{\pi_{k-1}(t)}}
\le
2C\varepsilon_k\,
\psi_\theta^{-1}\!\paren{\psi_\theta(x_\theta)+N(T,d,\varepsilon_k)}.
\]
For $k\ge 1$ we have $N(T,d,\varepsilon_k)\ge 2$, so by the definition of $K_\theta$,
\[
\psi_\theta^{-1}\!\paren{\psi_\theta(x_\theta)+N(T,d,\varepsilon_k)}
\le
K_\theta \bracks{\log N(T,d,\varepsilon_k)}^{1/\theta}.
\]
Hence
\[
\E \sup_{t\in T} X_t
\le
2CK_\theta
\sum_{k=1}^K
\varepsilon_k \bracks{\log N(T,d,\varepsilon_k)}^{1/\theta}.
\]
\medskip
\noindent\textbf{Step 3: Converting to an integral.}
Note that $\varepsilon_k=2(\varepsilon_k-\varepsilon_{k+1})$ and that $\varepsilon\mapsto N(T,d,\varepsilon)$ is nonincreasing. Therefore
\[
\varepsilon_k \bracks{\log N(T,d,\varepsilon_k)}^{1/\theta}
\le
2\int_{\varepsilon_{k+1}}^{\varepsilon_k}
\bracks{\log N(T,d,\varepsilon)}^{1/\theta}\,d\varepsilon.
\]
Summing over $k$ yields the desired bound
\[
\E \sup_{t\in T} X_t
\le
4CK_\theta
\int_0^{e(T)}
\bracks{\log N(T,d,\varepsilon)}^{1/\theta}\,d\varepsilon
\le
4CK_\theta
\int_0^\infty
\bracks{\log N(T,d,\varepsilon)}^{1/\theta}\,d\varepsilon.
\]
\end{proof}



\subsection{Proof of Corollary \ref{Dudley1}}\label{se:Dudley1}
\begin{proof}
By separability, there exists a countable dense set $T_0=\{t_1,t_2,\dots\}\subseteq T$ such that
\[
\sup_{t\in T} X_t = \sup_{t\in T_0} X_t
\qquad\text{a.s.}
\]
For $m\ge 1$, let
$T^{(m)}:=\{t_1,\dots,t_m\},~
M_m:=\sup_{t\in T^{(m)}}X_t$. Then $M_m\uparrow \sup_{t\in T_0}X_t$ a.s. Since $T^{(m)}\subseteq T$,
\[
N(T^{(m)},d,\varepsilon)\le N(T,d,\varepsilon),
\qquad \varepsilon>0.
\]
Applying Theorem~\ref{Dudley} to the finite set $T^{(m)}$,
\[
\E M_m
\le
4CK_\theta
\int_0^\infty
\bracks{\log N\paren{T^{(m)},d,\varepsilon}}^{1/\theta}\,d\varepsilon
\le
4CK_\theta
\int_0^\infty
\bracks{\log N(T,d,\varepsilon)}^{1/\theta}\,d\varepsilon.
\]
By the monotone convergence theorem,
\[
\E\sup_{t\in T} X_t
=
\E\sup_{t\in T_0} X_t
=
\lim_{m\to\infty}\E M_m
\le
4CK_\theta
\int_0^\infty
\bracks{\log N(T,d,\varepsilon)}^{1/\theta}\,d\varepsilon.
\]
\end{proof}



%%%%%%%%%%%%%%%%%%%%%%%%%%%%%%%%%
%%%%%%Proofs for Section 3%%%%%%%
%%%%%%%%%%%%%%%%%%%%%%%%%%%%%%%%%
\section{Proofs for Section \ref{sec:pre}}\label{proofse3}

This section collects proofs for the main results in Section \ref{sec:pre}.
\subsection{Proof of Lemma \ref{lem:key1}}
\begin{proof} 
We first establish the result for the case $1<\alpha\le 2$, and then treat the case $\alpha>2$.

\paragraph{Case $1<\alpha\le 2$.}
When $\theta\ge 1$ and $A=1$, the claim follows directly by combining
\cite[Proposition~1]{chu2023unified} with the monotonicity of \Renyi divergence stated in Lemma~\ref{Renyi}.

So we only need to bound the gap between $A>1$ and $A=1$ case:
\[
b_1=\E_Q \frac{dP}{dQ} \log^{1/\theta} \bigl(\frac{d P}{d Q} +A \bigr )-\E_Q \frac{dP}{dQ} \log^{1/\theta} \bigl(\frac{d P}{d Q} +1 \bigr ).
\]
Note that
\[
x\log^{1/\theta}(x+A)-x\log^{1/\theta}(x+1)\leq x \log^{1/\theta}\bigl ( 1+\frac{A-1}{x+1}\bigr ) \leq x (A-1)^{1/\theta}.
\]
Let $x=dP/dQ$. Taking expectation with respect to $Q$ yields $b_1\leq (A-1)^{1/\theta}$, which completes this case.

We therefore focus on the case $0<\theta<1$. Using the elementary inequality
$(x+y)^k \le x^k+y^k$ for $x,y\ge 0$ and $0<k\le 1$, we obtain
\[
\begin{aligned}
\DD_{f_\theta}(P\|Q)
&= \mathbb{E}_P\!\left[
\frac{1}{\alpha-1}
\log\!\left(\left|\frac{dP}{dQ}\right|+A\right)^{\alpha-1}
\right]^{1/\theta} \le
\mathbb{E}_P\!\left[
\frac{1}{\alpha-1}
\log\!\left(
\left|\frac{dP}{dQ}\right|^{\alpha-1}
+ A^{\alpha-1}
\right)
\right]^{1/\theta}.
\end{aligned}
\]

By Lemma~\ref{lemma:concave} in Appendix~\ref{app:lemmas}, the $\log^{1/\theta}(x+A^{\alpha-1})$ is concave on $[0,\infty)$ provided that
\[
A^{\alpha-1} \;\ge\; \exp\!\left(\frac{1}{\theta}-1\right).
\]
Applying Jensen's inequality then yields
\[
\DD_{f_\theta}(P\|Q)
\le
\left[
\frac{1}{\alpha-1}
\log\!\left(
\mathbb{E}_P\!\left|\frac{dP}{dQ}\right|^{\alpha-1}
+ A^{\alpha-1}
\right)
\right]^{1/\theta}.
\]

To separate the logarithmic terms, for $x\ge 1$ and $y\ge 0$,
$x+y \le x(1+y)$, 
\[
\log(x+y) \le \log x + \log(1+y).
\]
Thus it suffices to verify that
\[
\mathbb{E}_P\!\left|\frac{dP}{dQ}\right|^{\alpha-1} \;\ge\; 1.
\]
Indeed, writing $L:=\frac{dP}{dQ}\ge 0$ (defined $Q$-a.s.), Jensen's inequality gives
\[
\mathbb{E}_P[L^{\alpha-1}]
=\mathbb{E}_Q[L^\alpha]
\;\ge\;
\bigl(\mathbb{E}_Q L\bigr)^\alpha
=1.
\]
Consequently,
\[
\DD_{f_\theta}(P\|Q)
\le
\left[
\DD_\alpha(P\|Q)
+
\frac{1}{\alpha-1}\log(1+A^{\alpha-1})
\right]^{1/\theta},
\]
which establishes the desired bound for $1<\alpha\le 2$.

\paragraph{Case $\alpha>2$.}
For $\alpha>2$, we invoke the monotonicity of \Renyi divergence,
$\DD_2(P\|Q)\le \DD_\alpha(P\|Q)$, to obtain
\[
\DD_{f_\theta}(P\|Q)
\;\le\;
\inf_{\alpha>2}
\left\{
\bigl(\DD_\alpha(P\|Q)+C_{\alpha,\theta}\bigr)^{1/\theta}
\right\}
=
\bigl(\DD_2(P\|Q)+C_{2,\theta}\bigr)^{1/\theta}.
\]
\end{proof}

\subsection{Proof of Lemma \ref{de}}
\begin{proof}
Since
$\E_\mu r = \E_\nu \frac{d \mu}{d \nu} r$. Now set $x=\frac{d \mu}{d \nu}$, $y=r$ in Lemma \ref{young}, taking expectation gives
\[
\E_\mu r = \E_\nu \frac{d \mu}{d \nu} r\leq 2^{\frac{1}{\theta}} \E_\nu \{\frac{d \mu}{d \nu} [\log(\frac{d \mu}{d \nu} + A)]^{\frac{1}{\theta}}\} + \E_\nu\exp(y^\theta),
\]
and the definition of $\DD_{f_\theta}(\mu \| \nu)$ shows
\[
\E_\mu r \leq 2^\I\DD_{f_\theta}(\mu \| \nu) + \E_\nu \exp(r^\theta).
\]
Finally, Lemma~\ref{lem:key1} gives the R\'enyi upper bound
$\E_\mu r \leq 2^\I[\DD_\al(\mu \| \nu)+ C_{\alpha,\theta}]^\I+ \E_\nu \exp(r^\theta)$.
\end{proof}


% Next, we prove the truncated version of the above lemma.


% \subsection{Proof of Lemma \ref{deh}}
% \begin{proof}
% Suppose that $\mu$ and $\nu$ are probability measures on $(\Omega,\mathcal{F})$ such that $\mu \ll \nu$.
% Let $r : \Omega \to [0,\infty)$ be a non-negative $\mathcal{F}$-measurable function such that $r \in L^1(\mu)$ and $h(r) \in L^1(\nu)$. Let
% \[
% \E_\mu r = \E_\nu \frac{d \mu}{d \nu} r.
% \]
% Now we set $x=\frac{d \mu}{d \nu}$, $y=r$ in Lemma \ref{young1}, giving
% \[
% \E_\mu r \leq 2^\I \DD_{f_\theta}(\mu \| \nu) + 2\E_\nu h(r).
% \]
% Finally, applying Lemma~\ref{lem:key1}, we upper bound $\DD_{f_\theta}(\mu \| \nu)$ by R\'enyi divergence:
% \[
% \E_\mu r \leq 2^\I[\DD_\al(\mu \| \nu) + C_{\alpha,\theta}]^\I+2\E_\nu h(r).
% \]
% Here $h(y)=\exp(y^\theta)-\sum_{k=0}^{\lfloor \frac{2}{\theta} \rfloor} {y^{k\theta}}/{k!}$ is a truncated version of $\exp(y^\theta)$.
% \end{proof}








%%%%%%%%%%%%%%%%%%%%%%%%%%%%%%%%%%%%%%
%%%%%%%%%%%%%%%Section 4%%%%%%%%%%%%%%
%%%%%%%%%%%%%%%%%%%%%%%%%%%%%%%%%%%%%%
\section{Proofs for Section \ref{sec:gen_bounds}}\label{se:PCT}
\subsection{Mutual information for randomized maximum selector}
Let $S=\{X_1,\ldots,X_n\}$ with i.i.d.\ Weibull\((\theta)\) coordinates
(\(\theta<1\)), so $\mathbb{P}(X>x)=e^{-x^\theta}$. The selector \(W\in[n]\) is
defined by
\[
\mathbb{P}(W=y\mid S=x)=:p(y\mid x)=
\begin{cases}
\epsilon+\dfrac{1-\epsilon}{n}, & \text{if } y=\arg\max_{i\in[n]} x_i,\\[6pt]
\dfrac{1-\epsilon}{n}, & \text{otherwise,}
\end{cases}
\]
where the argmax is a.s.\ unique by the continuity of $\{X_i\}$. 

By symmetry, the marginal of \(W\) is uniform:
\[
p(y)=\mathbb{P}(W=y)=\frac{1}{n},\qquad y\in[n].
\]
Write
\[
p_m:=\epsilon+\frac{1-\epsilon}{n}=\frac{1+\epsilon(n-1)}{n},
\qquad
p_o:=\frac{1-\epsilon}{n}.
\]
Then, for any realization \(x\), exactly one index is the maximizer and
\[
\sum_{y=1}^n p(y\mid x)\log\frac{p(y\mid x)}{p(y)}
=\sum_{y=1}^n p(y\mid x)\log\bigl(n\,p(y\mid x)\bigr)
= p_m\log(n p_m) + (n-1)\,p_o\log(n p_o),
\]
which is constant in \(x\). Therefore,
\begin{align*}
I(W;S)
&=\int p(x)\sum_{y=1}^n p(y\mid x)\log\frac{p(y\mid x)}{p(y)}dx= p_m\log(n p_m) + (n-1)\,p_o\log(n p_o)\\
&=\frac{1}{n}\Bigl[(1+\epsilon(n-1))\,\log\bigl(1+\epsilon(n-1)\bigr)
+ (n-1)\,(1-\epsilon)\,\log(1-\epsilon)\Bigr].
\end{align*}

\subsection{Proof of Theorem \ref{infob}}
\begin{proof}
Without loss of generality, we assume $\max_{i\in[n]}\|X_i\|_{\psi_\theta}=1$. By changing measure,
\[
\mathbb{E}_{P_{W,S}}(|X_W|) = \mathbb{E}_{P_W \otimes P_S} \left(|X_W| \frac{d P_{W,S}}{d P_W \otimes P_S} \right).
\]
Let $\mu = P_{W,S}$, $\nu=P_{W} \otimes P_{S}$ and $r=|X_W|$. Use Lemma \ref{de} (decorrelation lemma)
\[
\E_{P_{W,S}}(|X_W|) \leq 2^\I I_{f_\theta}(W,S) + \E_{P_W \otimes P_S} \exp(|X_W|^\theta).
\]
By sub-Weibullity,
$\E_{P_W \otimes P_S} \exp |X_W|^\theta = \E_{w \sim P_W} \E_{P_S} \exp |X_w|^\theta \leq 2$. Then, if $A$ satisfies the conditions of Lemma~\ref{lem:key1}, then
\[
\mathbb{E}|X_W|
\;\le\;
\max_{i\in[n]}\|X_i\|_{\psi_\theta}\,
\Bigl(2^{1/\theta}\bigl(I_\alpha(W;S)+C_{\alpha,\theta}\bigr)^{1/\theta}+2\Bigr).
\]
\end{proof}

\subsection{Proof of Theorem \ref{chaining}}
\begin{proof}
Without loss of generality, assume $C=1$. \\
\textbf{Step 1: The partitions get finer.}
Assume that we are given increasing sequence partitions $\{\mathcal P_k\}_{k\ge 0}$ such that $\mathcal P_0=\{T\}$, and $\mathcal P_k$ is a $\epsilon_k$-partition for $k\ge1$ with $\epsilon_k := e(T)2^{-k}$. Let $[W]_k\in\mathcal P_k$ be the cell-valued random variable that contains $W$ with 
$$
T=[W]_0 \supset[W]_1 \supset[W]_2 \supset[W]_3 \supset \cdots.
$$
For each $k \ge 0$, let $W_k:=t_{[W]_k}$ be a designated representative point of $[W]_k \in \mathcal{P}_k$. 

For each cell $A \in \mathcal{P}_k$, choose a representative $t_A \in T$ such that 
$$A \subset \mathcal{B}_d\left(t_A, \varepsilon_k\right).$$
Moreover, choose these representatives hierarchically so that whenever $A \in \mathcal{P}_k$ is contained in its parent cell $B \in \mathcal{P}_{k-1}$, one has $t_A \in B$. Hence, for every such parent-child pair, it gives a parent-child distance 
\begin{equation}\label{eq: parent-child}
    d\left(t_A, t_B\right) \leq \varepsilon_{k-1} .
\end{equation}


\begin{figure}[t]
\centering 
\includegraphics[scale=0.078]{partition_cell_chaining_diagram_W_labels_hd.png}
\caption{Partition-cell chaining argument.}
\end{figure}


\textbf{Step 2: Telescope along the representatives for finite \( |T| \).}
For finite $|T|$, there exists a large $m$ such that each cell in $\mathcal{P}_m$ contains at most one point of $T$. Hence $W_{m}=W$. Pick an arbitrary anchor point $t_0 \in [W]_0=T$. Then
\[
X_W = X_{t_0} + \sum_{k=1}^{m}(X_{W_k}-X_{W_{k-1}}).
\]
For $k \geq 1$, by the hierarchical choice of representatives and $[W]_k \subset [W]_{k-1}$, \eqref{eq: parent-child} implies
\[
    d(W_{k}, W_{{k-1}})\le \epsilon_{k-1} = e(T)2^{-k+1}.
\]
Let $\delta_k = e(T)2^{-(k-1)}$. Because $\E[X_{t_0}] = 0$, the linearity of expectation shows,
\[
    \E[X_W] = \sum_{k=1}^m \delta_k \, \E \left[ \frac{X_{W_{k}} - X_{W_{{k-1}}}}{\delta_k} \right]\le \sum_{k=1}^m \delta_k \, \E \left| \frac{X_{W_{k}} - X_{W_{{k-1}}}}{\delta_k} \right|.
\]
To bound the expectation of the $k$-th term, set
$U_k:=(W_k,W_{k-1})$ and
$Z_k=|X_{W_k}-X_{W_{k-1}}|/\delta_k$. Applying the decorrelation lemma
(Lemma~\ref{de}) with $\mu=P_{U_k,S}$ and
$\nu=P_{U_k}\otimes P_S$ gives
\[
  \E_{P_{U_k,S}}[Z_k]
  \le
  2^{1/\theta} I_{f_\theta}(U_k; S)
  + \E_{P_{U_k}\otimes P_S} \left[ \exp(Z_k^\theta) \right].
\]
The pair $U_k$ is a deterministic function of the cell $[W]_k$, because
$W_k=t_{[W]_k}$ and $W_{k-1}$ is the representative of the parent cell.
Hence, by data processing for $f$-divergences,
$I_{f_\theta}(U_k;S)\le I_{f_\theta}([W]_k;S)$. Under the product measure
$P_{U_k}\otimes P_S$, condition on $U_k=(u,v)$. The indices $u,v$ are then
fixed and satisfy $d(u,v)\le\delta_k$, so the $(\theta,1)$-sub-Weibull
increment assumption gives
\[
    \E_{P_S}\exp\left(\frac{|X_u-X_v|^\theta}{\delta_k^\theta}\right)\le 2.
\]
Integrating over $P_{U_k}$ yields
$\E_{P_{U_k}\otimes P_S}[\exp(Z_k^\theta)]\le 2$. Substituting this back into
the sum gives
\[
    \E[X_W] \le \sum_{k=1}^m e(T)2^{-(k-1)} \left( 2^{1/\theta} I_{f_\theta}([W]_k; S) + 2 \right),
\]
which proves the theorem for finite $T$.

\textbf{Step 3: Extension to separable processes.}
For $m\ge 1$, define the $m$-th hierarchical approximation of $W$ by
$W_m:=t_{[W]_m}$. The representatives are chosen along the separability
approximation described in Definition~\ref{def:separable}. Since
$W\in[W]_m\subset B_d(W_m,e(T)2^{-m})$, we have
\[
d(W_m,W)\le e(T)2^{-m}\xrightarrow[m\to\infty]{}0
\quad\text{a.s.}
\]
%Although $W_m$ is random, separability is applied pathwise: for almost every
%outcome, $W_m(\omega)$ is the deterministic separability approximation of
%$W(\omega)$. 
Almost-sure sample-path continuity
 implies
\[
X_{W_m}\xrightarrow[m\to\infty]{}X_W
\quad\text{a.s.}
\]
To justify dominated convergence $\E[X_W] = \lim_{m \to \infty} \E[X_{W_m}]$, fix any $u_0 \in T$. Then
$$
\left|X_{W_m}\right| \leq\left|X_{u_0}\right|+\sup _{t \in T}\left|X_t-X_{u_0}\right| .
$$
The first term is integrable because a sub-Weibull random variable has finite first moment. For the second term,
$$
\sup _{t \in T}\left|X_t-X_{u_0}\right| \leq \sup _{t \in T}\left(X_t-X_{u_0}\right)+\sup _{t \in T}\left(X_{u_0}-X_t\right) .
$$
Both processes
$\left\{X_t-X_{u_0}\right\}_{t \in T}$ and $\left\{X_{u_0}-X_t\right\}_{t \in T}
$ are again separable $(\theta, C)$-sub-Weibull processes, so Corollary \ref{Dudley1} implies that each supremum has finite expectation. The dominated convergence theorem justifies the passage to the limit:
\[
    \E[X_W] = \lim_{m \to \infty} \E[X_{W_m}] \le \sum_{k=1}^\infty e(T)2^{-(k-1)} \left( 2^{1/\theta} I_{f_\theta}([W]_k; S) + 2 \right).
\]
	The R\'enyi refinement follows immediately by applying Lemma~\ref{lem:key1} to bound the $I_{f_\theta}$ term via R\'enyi mutual information.
\end{proof}

%To generalize to separable case, what we need is $\lim_{k\to \infty} \E [ X_{W_k} ]=\E[X_W]$. This is true when we can apply dominated convergence theorem, which need $X_{W_k}$ be uniformly bounded by a integrable function. When $\E[\sup_w X_w]<\infty$ hold, which can be bounded by $\gamma$ in corollary \ref{Dudley1} this condition is satisfied.  





% \subsection{Proof of Theorem \ref{gen}}

% The proof is based on the following lemma, adapted from the proof of Theorem~3.1 in \citet{kuchibhotla2022moving}.
% %\scomment{this sentence is unfinished.}
% \begin{lemma} \label{moment}
% Let $\{X_i\}_{i=1}^n$ be a sequence of independent zero-mean r.v.s with $\| X_i\|_{\psi_\theta} \leq 1$. Then, for $p\geq 2$ and $0<\theta \leq1 $, we have
% \begin{equation}\label{eq:lessone}
% \left\| \sum_{i=1}^{n} a_{i} X_{i} \right\|_{p} \leq L_\theta \left[ \sqrt{p} \|a\|_{2} + p^{1/\theta} \|a\|_{\infty} \right].
% \end{equation}
% where $L_\theta=:\sqrt{8} e^{3} (2\pi)^{1/4} e^{1/24} (e^{2/e} / \theta)^{1/\theta}$.
% And if $\theta \ge 1$,
% \begin{equation}\label{eq:bigone}
% \left\| \sum_{i=1}^{n} a_{i} X_{i} \right\|_{p}  \leq (4 e +2 (\log 2)^{1/\theta} )\sqrt{p} \|a\|_{2} + 4e p^{1/\theta} \|a\|_{\frac{\theta}{\theta-1}} .
% \end{equation}
% \end{lemma}

% Note that the moment inequality above only holds for $p\geq 2$, we need to introduce a truncated function of $\exp(y^\theta)$: $h(y)=\exp(y^\theta)-\sum_{k=0}^{\lfloor \frac{2}{\theta} \rfloor} \frac{y^{k\theta}}{k!}$.
% %----------------------------------------------hbound hbound
% \begin{lemma}[A truncated exponential moment bound]\label{hbound}
% Let $\{X_i\}_{i=1}^n$ be a sequence of independent zero-mean sub-Weibull r.v.s with $\|X_i\|_{\psi_\theta}\leq 1$. Then, for a positive constant
% \begin{equation}\label{eq:theta}
%     M_\theta = 
% \begin{cases}
% 2^{\frac{1}{\theta}}(\sqrt{\theta} + \theta^{\frac{1}{\theta}})e^{\frac{1}{\theta}} L_\theta, & 0 < \theta < 1 \\[10pt]
% \left[(4e + 2(\log 2)^{1/\theta})\sqrt{\theta} + 4e\theta^{\frac{1}{\theta}}\right](2e)^{\frac{1}{\theta}}, & \theta \geq 1
% \end{cases}
% \end{equation}
% and $L_\theta=:\sqrt{8} e^{3} (2\pi)^{1/4} e^{1/24} (e^{2/e} / \theta)^{1/\theta}$, which depend only on $\theta$. 
% For $0<\theta\leq2$, we have
% \begin{equation*}
% \E \left[h\left (\frac{\sum_{k=1}^n X_k t_k}{M_\theta\|t\|_2}\right )\right]\leq 2.
% \end{equation*}
% And for $\theta>2$,
% \begin{equation}
% \E \left[h\left (\frac{\sum_{k=1}^n X_k t_k}{M_\theta \|t\|_{\frac{\theta}{\theta-1}}}\right )\right]\leq 2.
% \end{equation}
% \end{lemma}

% \begin{proof}
% We prove the result by cases.

% \noindent\textbf{Case 1 ($\theta<1$).} Note that $h(y)=\exp(y^\theta)-\sum_{k=0}^{\lfloor \frac{2}{\theta} \rfloor} \frac{y^{k\theta}}{k!}$. By Lemma~\ref{moment},
% \[
% \begin{aligned}
% \E \left[h\left (\frac{\sum_k X_k t_k}{M_\theta \|t\|_2}\right )\right]
% &\leq\sce \frac{1}{k!}  \E \left[\left (\frac{\sum_k X_k t_k}{M_\theta \|t\|_2}\right )^{k\theta}\right]\leq \sce \frac{1}{k!} \left(\frac{L_\theta}{M_\theta \|t\|_2}\right)^{k \theta}  (\|t\|_2 \sqrt{k \theta} + (k\theta)^\I \|t\|_\infty)^{k\theta}\\
% &\leq \sce (\sqrt{\theta} k^{\frac{1}{2}-\I}+\theta^\I\frac{\|t\|_\infty}{\|t\|_2})^{k\theta} \left (\frac{e^\I L_\theta}{M_\theta}\right )^{k\theta} \quad (k!\geq (k/e)^k)\\
% &\leq \sum_{k=1}^{\infty} \left (\frac{(\sqrt{\theta}+\theta^\I)e^\I L_\theta}{M_\theta} \right)^{k\theta}=\sum_{k=1}^{\infty}2^{-k}=2,
% \end{aligned}
% \]
% where we choose $M_\theta=2^\I(\sqrt{\theta}+\theta^\I)e^\I L_\theta$ for $\theta<1$.

% \noindent\textbf{Case 2 ($1 \leq \theta \leq 2$)}. Let $\beta=\frac{\theta}{\theta-1}\geq 2$. By monotonicity of the $\ell_p$ norm, we have 
% \begin{equation}\label{Norm}
% \|t\|_\beta \leq \|t\|_2.
% \end{equation}
% Then,
% \[
% \begin{aligned}
% \E \left[h\left (\frac{\sum_k X_k t_k}{M_\theta \|t\|_2}\right )\right]
% &\leq\sce \frac{1}{k!}  \E \left[\left (\frac{\sum_k X_k t_k}{M_\theta\|t\|_2}\right )^{k\theta}\right]\\
% &\leq \sce \frac{1}{k!} \left(\frac{1}{M_\theta \|t\|_2}\right)^{k \theta}[(4 e +2 (\log 2)^{1/\theta})\|t\|_2 \sqrt{k \theta} + (4e )(k\theta)^{\I} \|t\|_\beta]^{k\theta}\\
% (\text{By}~k!\geq (k/e)^k~\text{and}~\eqref{Norm})
% &\leq \sce\lk (4 e +2 (\log 2)^{1/\theta} ) \sqrt{\theta} k^{\frac{1}{2}-\I}+(4e)\theta^{\I}  \rk^{k\theta} \left(\frac{e^{\I}}{M_\theta}\right)^{k\theta} \\
% &\leq \sum_k \left(\frac{\lk (4 e +2 (\log 2)^{1/\theta} )  \sqrt{\theta}+ 4 e \theta^\I \rk e^\I}{M_\theta} \right)^{k\theta}.
% \end{aligned}
% \]
% So we choose $M_\theta= [(4 e +2 (\log 2)^{1/\theta} )  \sqrt{\theta}+ 4 e \theta^\I ](2e)^\I $ for $1 \leq \theta \leq 2$.

% \noindent\textbf{Case 3~($\theta > 2$).} Let $\beta=\frac{\theta}{\theta-1}\in(1,2)$, we have 
% \begin{equation}\label{Norm1}
% \|t\|_2 \leq \|t\|_\beta
% \end{equation}
% Using the same derivation as Case 2, it gives
% \$
% \E \left[h\left (\frac{\sum_k X_k t_k}{M_\theta \|t\|_\beta}\right )\right]
% &\leq\sce \frac{1}{k!}  \E \left[\left (\frac{\sum_k X_k t_k}{M_\theta\|t\|_\beta}\right )^{k\theta}\right]\\
% &\leq \sce \frac{1}{k!} \left(\frac{1}{M_\theta \|t\|_\beta}\right)^{k \theta}  ((4 e +2 (\log 2)^{1/\theta} ) \|t\|_2 \sqrt{k \theta} + (4e )(k\theta)^\I \|t\|_\beta)^{k\theta}\\
% (k!\geq (k/e)^k \text{and}~\eqref{Norm1})~&\leq \sce \lk (4 e + 2 (\log 2)^{1/\theta} ) \sqrt{\theta} k^{\frac{1}{2}-\I} +(4e)\theta^{\I} \rk^{k\theta} \left( \frac{e^{\I}}{M_\theta} \right)^{k\theta}  \\
% &\leq \sum_k \Big (\frac{\lk (4 e +2 (\log 2)^{1/\theta} )  \sqrt{\theta}+ 4 e \theta^\I \rk e^\I}{M_\theta} \Big)^{k\theta}.
% \$
% So we choose $M_\theta= [ (4 e +2 (\log 2)^{1/\theta} )  \sqrt{\theta}+ 4 e \theta^\I ](2e)^\I $ for $\theta > 2$.
% \end{proof}



% \begin{proof}[Proof of Theorem~\ref{gen}]
% Write $S=\{Z_1,\dots,Z_n\}$ and $\bar\ell(w,Z):=\E_\mu[\ell(w,Z)]-\ell(w,Z)$,
% \[
% \gen(W,S)=\frac1n\sum_{i=1}^n \bar\ell(W,Z_i).
% \]
% Let
% $C:=\sup_{w\in\mathcal W}\|\bar\ell(w,Z)\|_{\psi_\theta}\in(0,\infty)$, and define the normalized variables
% \[
% X_i:=\frac{\bar\ell(w,Z_i)}{C},\qquad i=1,\dots,n.
% \]
% Random variables $X_1,\dots,X_n$ are independent, centered, and satisfy
% $\|X_i\|_{\psi_\theta}\le 1$.

% \noindent
% \textbf{Step 1: a truncated exponential moment bound under $P_W\otimes P_S$.}
% Fix $t=(t_1,\dots,t_n)$ with $t_i\equiv 1/n$.
% Then $\|t\|_2 = n^{-1/2}$.
% Applying Lemma~\ref{hbound} to $\{X_i\}_{i=1}^n$ with this choice of $t$, we obtain
% (for some constant $M_\theta>0$)
% \[
% \E_{P_W\otimes P_S}\Biggl[
% h\Bigl(\frac{\sum_{i=1}^n X_i t_i}{M_\theta\|t\|_2}\Bigr)
% \Biggr]\le 2.
% \]
% Since $\sum_{i=1}^n X_i t_i = \frac1n\sum_{i=1}^n X_i=\gen(W,S)/C$ and
% $M_\theta\|t\|_2=M_\theta/\sqrt n$, the last display is
% \[
% \E_{P_W\otimes P_S}\Biggl[
% h\Bigl(\frac{\sqrt n\,|\gen(W,S)|}{M_\theta\,C}\Bigr)
% \Biggr]\le 2.
% \]
% \textbf{Step 2: change of measure via the truncated decorrelation lemma.}
% Let $\mu:=P_{W,S}$ and $\nu:=P_W\otimes P_S$.
% Define the nonnegative random variable
% \[
% r:=\frac{\sqrt n\,|\gen(W,S)|}{M_\theta\,C}.
% \]
% Then Step~1 gives $\E_\nu[h(r)]\le 2$, hence $h(r)\in L^1(\nu)$.
% Applying Lemma~\ref{deh} with this $(\mu,\nu,r)$ yields
% \[
% \E_\mu[r]
% \le 2^{1/\theta} D_{f_\theta}(\mu\|\nu) + 2\,\E_\nu[h(r)].
% \]
% Using $\E_\nu[h(r)]\le 2$ from Step~1, we obtain the upper bound
% \[
% \E_{P_{W,S}}[r]\le 2^{1/\theta} D_{f_\theta}(P_{W,S}\|P_W\otimes P_S) + 4.
% \]
% \textbf{Step 3: conclude and identify the information term.}
% Multiplying the last inequality by $M_\theta C/\sqrt n$ gives
% \[
% \E\bigl|\gen(W,S)\bigr|
% \le \frac{M_\theta C}{\sqrt n}\Bigl(
% 2^{1/\theta} D_{f_\theta}(P_{W,S}\|P_W\otimes P_S) + 4
% \Bigr).
% \]
% The definition of the $f_\theta$-mutual information $I_{f_\theta}(S;W)=D_{f_\theta}(P_{W,S}\|P_W\otimes P_S)$ gives
% \[
% \E\bigl|\gen(W,S)\bigr|
% \le \frac{M_\theta C}{\sqrt n}\,
% \Bigl(2^{1/\theta} I_{f_\theta}(S;W)+4\Bigr).
% \]
% Finally, the refinement using R\'enyi mutual information $I_\alpha(S; W)$ follows immediately by applying Lemma~\ref{lem:key1} to upper bound $I_{f_\theta}(S; W)$.
% \end{proof}



% \subsection{Proof of Theorem \ref{gen1}}
% First, we need the following lemma.
% \begin{lemma}[A change-of-measure inequality with truncated sub-Weibull tail]\label{lem:pb_trunc}
% Let $(\Omega,\mathcal F)$ be a measurable space and let $\mu,\nu$ be probability
% measures on $(\Omega,\mathcal F)$ such that $\mu\ll \nu$. Fix $0<\theta\le 2$, and let
% $h(x)=\exp(x^\theta)-\sum_{k=0}^{\lfloor 2/\theta\rfloor}\frac{x^{k\theta}}{k!}$  with $x\ge 0$. Let $g:\Omega\to[0,\infty)$ be measurable such that $h(g)\in L^1(\nu)$ and
% $\bigl(\log\frac{d\mu}{d\nu}\bigr)_+^{1/\theta}\in L^1(\mu)$. Then
% \begin{equation}\label{eq:pb_trunc_main}
% \E_\mu[g]\le
% 4^{1/\theta}\Biggl\{
% \Bigl(\log(\E_\nu[h(g)]+1)\Bigr)^{1/\theta}
% +\E_\mu\left[\left(\log\frac{d\mu}{d\nu}\right)_+^{1/\theta}\right]
% +E_\theta
% \Biggr\},
% \end{equation}
% where the constant $E_\theta>0$ satisfying
% $4^{1/\theta}E_\theta:=\sup_{x\ge 1}\frac{x\,\exp(x^\theta/2)}{h(x)+1}+1.$
% \end{lemma}
% \begin{proof}
% Set $\phi:=\frac{d\mu}{d\nu}$ and denote $
% G:=\E_\nu[h(g)]+1\in(1,\infty)$. Define the measurable set
% \[
% T
% :=
% \Bigl\{\omega\in\Omega:\;
% \phi(\omega)\ge \frac{\exp(g(\omega)^\theta/2)}{G}\Bigr\}.
% \]
% We decompose $E_\mu[g]=\int g\,d\mu=\int_T g\,d\mu+\int_{T^c} g\,d\mu$, and bound the two terms separately.

% \textbf{Step 1: bound over $T^c$.}
% Write $T^c=(T^c\cap\{g<1\})\cup(T^c\cap\{g\ge 1\})$. Since $g\ge 0$,
% \[
% \int_{T^c\cap\{g<1\}} g\,d\mu\le 1.
% \]
% On $T^c$ we have $\phi<\exp(g^\theta/2)/G$, hence on $T^c\cap\{g\ge 1\}$,
% \[
% g\,d\mu = g\phi\,d\nu
% \le \frac{g\exp(g^\theta/2)}{G}\,d\nu.
% \]
% Therefore,
% \begin{equation}\label{eq:Tc_bound}
% \int_{T^c} g\,d\mu
% \le
% 1+\int_{\{g\ge 1\}} \frac{g\exp(g^\theta/2)}{G}\,d\nu
% \le
% 1+\frac{1}{G}\sup_{x\ge 1}\frac{x\exp(x^\theta/2)}{h(x)+1}\int (h(g)+1)\,d\nu.
% \end{equation}
% Since $\int (h(g)+1)\,d\nu=G$, the right-hand side of \eqref{eq:Tc_bound} becomes
% \[
% \int_{T^c} g\,d\mu
% \le
% \sup_{x\ge 1}\frac{x\exp(x^\theta/2)}{h(x)+1}+1
% =
% 4^{1/\theta}E_\theta.
% \]

% \textbf{Step 2: bound over $T$.}
% On $T$ we have $\phi \ge \exp(g^\theta/2)/G$, hence
% $g^\theta
% \le
% 2\log(G\phi)=2\bigl(\log G+\log\phi\bigr)
% \le
% 2\bigl(\log G+(\log \phi)_+\bigr)$, so that
% \begin{equation}\label{eq:g_control_on_T}
% g
% \le
% 2^{1/\theta}\bigl(\log G+(\log \phi)_+\bigr)^{1/\theta}
% \le
% 4^{1/\theta}\bigl((\log G)^{1/\theta}+(\log \phi)_+^{1/\theta}\bigr),
% \end{equation}
% where the last inequality uses $(a+b)^{1/\theta}\le 2^{1/\theta}(a^{1/\theta}+b^{1/\theta})$ for
% $0<\theta\le 2$ by $C_r$-inequality. Integrating \eqref{eq:g_control_on_T} w.r.t.\ $\mu$ over $T$ gives
% \[
% \int_T g\,d\mu
% \le
% 4^{1/\theta}(\log G)^{1/\theta}
% +
% 4^{1/\theta}\E_\mu(\log \phi)_+^{1/\theta}.
% \]
% Combining the bounds for $\int_T g\,d\mu$ and $\int_{T^c} g\,d\mu$, we obtain \eqref{eq:pb_trunc_main}.
% \end{proof}

% \begin{proof}[Proof of Theorem~\ref{gen1}]
% Let $C:=\sup_{w\in\mathcal W}\|\bar\ell(w,Z)\|_{\psi_\theta}<\infty$ for $0<\theta\le 2$,
% and write the generalization error as
% \[
% \operatorname{gen}(w,S)
% =\mathbb E[\ell(w,Z)]-\frac1n\sum_{i=1}^n \ell(w,Z_i)
% =-\frac1n\sum_{i=1}^n \bar\ell(w,Z_i).
% \]
% Fix $w\in\mathcal W$, $\bar\ell(w,Z_1),\dots,\bar\ell(w,Z_n)$ are independent,
% mean-zero with $\|\bar\ell(w,Z_i)\|_{\psi_\theta}\le C$.

% \medskip
% \noindent\textbf{Step 1: truncated MGF bound and a change-of-measure inequality.}
% Let $h(y)=\exp(y^\theta)-\sum_{k=0}^{\lfloor 2/\theta\rfloor} y^{k\theta}/k!$.
% Apply Lemma~\ref{hbound} to
% $X_i:=\bar\ell(w,Z_i)/C$ and $t_i:=1/\sqrt n$.
% Since $\|X_i\|_{\psi_\theta}\le 1$ and $\|t\|_2=1$, there exists a constant
% $K_\theta>0$ (depending only on $\theta$) such that
% \begin{equation}\label{eq:prod-h-bound}
% \mathbb E_{P_S}\, h\!\left(\frac{\sqrt n\,|\operatorname{gen}(w,S)|}{K_\theta C}\right)\le 2.
% \end{equation}
% Integrating \eqref{eq:prod-h-bound} w.r.t.\ $w\sim P_W$ yields
% \begin{equation}\label{eq:prod-h-bound-W}
% \mathbb E_{P_S\otimes P_W}\, h\!\left(\frac{\sqrt n\,|\operatorname{gen}(W,S)|}{K_\theta C}\right)\le 2.
% \end{equation}
% Define the random variable
% $g(w):=\frac{\sqrt n\,|\operatorname{gen}(w,S)|}{K_\theta C}$, viewed as a function of $w$ for the realized sample $S$.
% For each fixed $S$, apply Lemma~\ref{lem:pb_trunc} with $\mu=P_{W\mid S}$ and $\nu=P_W$ to obtain
% \begin{equation}\label{eq:lemma41-applied}
% \mathbb E_{P_{W\mid S}} g
% \le 4^{1/\theta}\Bigl(
% \bigl(\log(\mathbb E_{P_W} h(g)+1)\bigr)^{1/\theta}
% +\mathbb E_{P_{W\mid S}}\!\left[
% \left(\log\frac{dP_{W\mid S}}{dP_W}\right)_+^{1/\theta}
% \right]
% +E_\theta
% \Bigr),
% \end{equation}
% where $E_\theta>0$ depends only on $\theta$. Multiplying \eqref{eq:lemma41-applied} by $K_\theta C/\sqrt n$ gives
% \begin{equation}\label{eq:cond-bound-before-hp}
% \mathbb E_{P_{W\mid S}}|\operatorname{gen}(W,S)|
% \le
% \frac{4^{1/\theta}K_\theta C}{\sqrt n}\Bigl(
% \mathbb E_{P_{W\mid S}}\!\left[
% \left(\log\frac{dP_{W\mid S}}{dP_W}\right)_+^{1/\theta}
% \right]
% +E_\theta
% +\bigl(\log(\mathbb E_{P_W} h(g)+1)\bigr)^{1/\theta}
% \Bigr).
% \end{equation}

% \medskip
% \noindent\textbf{Step 2: high-probability control of $\mathbb E_{P_W} h(g)$.}
% By \eqref{eq:prod-h-bound-W},
% \[
% \mathbb E_{P_S}\,\mathbb E_{P_W}\,[h(g)+1]
% =
% \mathbb E_{P_S\otimes P_W}[h(g)+1]
% \le 3.
% \]
% Hence, by Markov's inequality, for any $\delta\in(0,1)$,
% \[
% P_S\!\left(\mathbb E_{P_W}[h(g)+1]>\frac{3}{\delta}\right)
% \le \frac{\delta}{3}\,\mathbb E_{P_S}\mathbb E_{P_W}[h(g)+1]
% \le \delta.
% \]
% Therefore, on an event $\mathcal E:=\{\mathbb E_{P_W}[h(g)+1]\le 3/\delta\}$ with $P_S(\mathcal E)\ge 1-\delta$, we have
% $\mathbb E_{P_W}[h(g)+1]\le 3/\delta$ and thus
% \[
% \log(\mathbb E_{P_W} h(g)+1)
% \le \log(\mathbb E_{P_W}[h(g)+1])
% \le \log(3/\delta).
% \]
% Substituting this into \eqref{eq:cond-bound-before-hp} proves that on $\mathcal E$,
% \[
% \mathbb E_{P_{W\mid S}}|\operatorname{gen}(W,S)|
% \le
% \frac{4^{1/\theta}K_\theta C}{\sqrt n}\Bigl(
% \mathbb E_{P_{W\mid S}}\!\left[
% \left(\log\frac{dP_{W\mid S}}{dP_W}\right)_+^{1/\theta}
% \right]
% +E_\theta
% +(\log(3/\delta))^{1/\theta}
% \Bigr)
% \]
% with probability at least $1-\delta$.
% \end{proof}


% \subsection{Proof of Theorem \ref{genchain}}
% \begin{proof}
% Similar to the proof of Theorem \ref{chaining}, for each $k \ge 1$, let $\epsilon_k = e(\mathcal{W})2^{-k}$. Let $W_k:=t_{[W]_k}$ be a representative element of the cell $[W]_k \in \mathcal{P}_k$. 

% First, consider the case where $|\mathcal{W}| < \infty$. The sequence of partitions is increasing and finite; hence, there exists an integer $m$ sufficiently large such that $W_m= W$ almost surely. We decompose $\gen(W, S)$ as a telescoping sum. 


% Choose a deterministic point $w_0\in \mathcal{W}$ and set $W_0:=w_0$. Since $\bar\ell(w_0,Z)$ is centered,
% \[
% \E[\gen(w_0,S)] = \E\!\left[\frac1n\sum_{i=1}^n \bar\ell(w_0,Z_i)\right] = 0.
% \]
% Therefore,
% \[
% \gen(W,S)-\gen(w_0,S)
% =
% \sum_{k=1}^m \bigl(\gen(W_k,S)-\gen(W_{k-1},S)\bigr),
% \]
% and hence
% \begin{equation}
% \label{eq:finite22-1}
% \E[\gen(W,S)]
% \le
% \sum_{k=1}^m \E\left|\gen(W_k,S)-\gen(W_{k-1},S)\right|.
% \end{equation}
% Therefore, the linearity implies
% \[
% \mathbb E[\gen(W_m,S)]
% =
% \sum_{k=1}^m \mathbb E\Bigl[\gen(W_k,S)-\gen(W_{k-1},S)\Bigr]= \sum_{k=1}^m \E \left[ \frac{1}{n} \sum_{i=1}^n \left( \bar{\ell}(W_k, Z_i) - \bar{\ell}(W_{k-1}, Z_i) \right) \right].
% \]
% By the sub-Weibull process assumption, $\|\bar{\ell}(W_k, Z) - \bar{\ell}(W_{k-1}, Z)\|_{\psi_\theta} \le d(W_k,W_{k-1})$. By construction of the partitions, $d(W_k, W_{k-1}) \le \epsilon_{k-1} =e(\mathcal{W})2^{-(k-1)}$.
% Let $\delta_k = e(\mathcal{W})2^{-(k-1)}$. We normalize the increments by $\delta_k$:
% \[
%     \E[\gen(W_m, S)] = \sum_{k=1}^m \frac{\delta_k}{\sqrt{n}} \, \E \left[ \frac{\sqrt{n} \sum_{i=1}^n (\bar{\ell}(W_k, Z_i) - \bar{\ell}(W_{k-1}, Z_i))}{n \delta_k} \right].
% \]
% Let $Y_k = \frac{\sqrt{n} \sum_{i=1}^n (\bar{\ell}(W_k, Z_i) - \bar{\ell}(W_{k-1}, Z_i))}{n \delta_k}$ and set
% $U_k:=(W_k,W_{k-1})$. We apply the truncated decorrelation lemma
% (Lemma~\ref{deh}) for $r_k=|Y_k|/M_\theta$ with
% $\mu=P_{U_k,S}$ and $\nu=P_{U_k}\otimes P_S$, then
% \begin{equation}\label{eq:thm22}
%       \E_{P_{U_k,S}}[r_k] \le 2^{1/\theta} I_{f_\theta}(U_k; S) + 2 \, \E_{P_{U_k} \otimes P_S} [h(r_k)],
% \end{equation}
% where $h(y)$ is the truncated exponential. Since $U_k$ is a deterministic
% function of $[W]_k$, data processing gives
% $I_{f_\theta}(U_k;S)\le I_{f_\theta}([W]_k;S)$.
% Under the product measure $P_{U_k} \otimes P_S$, condition on
% $U_k=(u,v)$. The indices $u,v$ are fixed and satisfy $d(u,v)\le\delta_k$, and
% the term $Y_k$ is a normalized sum of independent, mean-zero sub-Weibull random
% variables.

% Thus the truncated moment bound (Lemma \ref{hbound}) with appropriate scaling
% $M_\theta$ for $0<\theta\leq2$ applies conditionally on $U_k=(u,v)$ with the
% choice $t_i\equiv 1/n$, whose Euclidean norm is $\|t\|_2 = n^{-1/2}$. Let
% $h(y)=e^{y^\theta}-\sum_{j=0}^{\lfloor 2/\theta\rfloor}y^{j\theta}/j!$. We have
% \[
%     \E_{P_S} \left[ h\left( \frac{|Y_k|}{M_\theta}\right) \mid U_k=(u,v)\right] \le 2.
% \]
% Integrating over $P_{U_k}$ gives
% $\E_{P_{U_k} \otimes P_S} [h(r_k)] \le 2$ and substituting it into
% \eqref{eq:thm22}, we have
% \[
%     \E [ | \gen(W_k, S) - \gen(W_{k-1}, S) | ] \le \frac{M_\theta \delta_k}{\sqrt{n}} \left( 2^{1/\theta} I_{f_\theta}([W]_k; S) + 4 \right).
% \]
% Summing over $k$:
% \[
%     \E [\gen(W_m, S)] \le \frac{M_\theta e(\mathcal{W})}{\sqrt{n}} \sum_{k=1}^m 2^{-(k-1)} \left( 2^{1/\theta} I_{f_\theta}([W]_k; S) + 4 \right).
% \]
% This proves the result for finite $\mathcal{W}$.

% For a general separable metric space $\mathcal{W}$, define
% $W_m:=t_{[W]_m}$ with representatives chosen along the separability
% approximation of the loss process. Since
% $W\in[W]_m\subset B_d(W_m,e(\mathcal{W})2^{-m})$, we have
% $d(W_m,W)\to 0$ a.s. Although $W_m$ is random, the separability property is
% applied pathwise with $w=W(\omega)$, and therefore
% \[
% \gen(W_m,S)\xrightarrow[m\to\infty]{} \gen(W,S)
% \quad\text{a.s.}
% \]


% The condition that the Dudley integral of the process is finite implies that the supremum of the process is integrable:
% \[
% |\gen(W_m,S)|\le \sup_{w\in\mathcal{W}}|\gen(w,S)|\in L^1.
% \]
% By the dominated convergence theorem, we can pass to the limit $m \to \infty$ in the expectation,
% \[
% \mathbb{E}\bigl[\operatorname{gen}(W,S)\bigr]= \lim_{m \to \infty} \mathbb{E}\bigl[\operatorname{gen}(W_m,S)\bigr]\le
% \frac{e(\mathcal{W})M_\theta}{\sqrt{n}}
% \sum_{k=1}^\infty 2^{-(k-1)}
% \bigl(2^{1/\theta} I_{f_\theta}([W]_k;S)+4\bigr)
% \]
% This yields the infinite-series bound \eqref{eq:cgb}.
% \end{proof}



%%%%%%%%%%%%%%%%%%%%%%%%%%%%%%%%%%%%%
%%%%%%%%%%%%%Section E%%%%%%%%%%%%%%%
%%%%%%%%%%%%%%%%%%%%%%%%%%%%%%%%%%%%%
\section{Proofs in Section \ref{sec:rlhf}}\label{se:API}
\subsection{Proof of Lemma \ref{renyisol}}
\begin{proof}
Fix $x\in\mathcal X$ and write $P_0:=\pi_0(\cdot\mid x),~
r(y):=r(x,y)$. For any feasible policy $\pi(\cdot\mid x)\ll P_0$, let
$
q(y):=\frac{d\pi(\cdot\mid x)}{dP_0}(y).
$
The optimization problem is equivalent to
\[
\sup_{q\ge0}
\int_{\mathcal Y} r(y)q(y)\,P_0(dy)
\]
subject to $\int_{\mathcal Y}q\,dP_0=1,~\int_{\mathcal Y}q^\alpha\,dP_0\le \rho$, where
$\rho:=\exp((\alpha-1)\epsilon)$.

Since $r\in L^{\alpha/(\alpha-1)}(P_0)$ and $q\in L^\alpha(P_0)$ for every
feasible $q$, the objective is finite by H\"older's inequality. The feasible
set is convex, and the objective is linear. Moreover, since $\epsilon>0$,
we have $\rho>1$, and $q\equiv1$ satisfies
\[
\int q\,dP_0=1,
\qquad
\int q^\alpha\,dP_0=1<\rho.
\]
Thus Slater's condition holds, so the KKT conditions are necessary for
optimality. Let
\[
r_x^\star:=\operatorname*{ess\,sup}_{y\sim P_0}r(y),
\qquad
M_x:=\{y:r(y)=r_x^\star\},
\qquad
m_x:=P_0(M_x).
\]

We first characterize when a policy supported entirely on $M_x$ is feasible.
Suppose $q$ is supported on $M_x$. Then
\[
\int_{M_x}q\,dP_0=1.
\]
If $m_x=0$, no such density exists. If $m_x>0$, Jensen's inequality under the
conditional probability measure $P_0(\cdot\mid M_x)$ gives
\[
\int q^\alpha\,dP_0
=
\int_{M_x}q^\alpha\,dP_0
=
m_x\int_{M_x}q^\alpha\,dP_0(\cdot\mid M_x)
\ge
m_x
\left(
\int_{M_x}q\,dP_0(\cdot\mid M_x)
\right)^\alpha.
\]
Since $\int_{M_x}q\,dP_0(\cdot\mid M_x)
=
\frac1{m_x}\int_{M_x}q\,dP_0
=
\frac1{m_x}$, we get
\[
\int q^\alpha\,dP_0
\ge
m_x^{1-\alpha}.
\]
Equality is attained by $
q_M(y)=\frac{\mathbf 1_{M_x}(y)}{m_x}$. Therefore, the minimal R\'enyi divergence among policies supported on $M_x$ is $\frac1{\alpha-1}\log m_x^{1-\alpha}
=
-\log m_x$. Hence, if
\[
\epsilon\ge -\log m_x,
\]
then $q_M$ is feasible. Since $q_M$ puts all its mass on the set where
$r(y)=r_x^\star$, it achieves the largest possible objective value
\[
\int r(y)q_M(y)\,P_0(dy)=r_x^\star.
\]
No policy can achieve objective value larger than $r_x^\star$ by the definition
of essential supremum. Thus $q_M$ is optimal. This proves the first case.

It remains to consider the case
\[
\epsilon<-\log m_x.
\]
In this case, no feasible policy can be supported entirely on $M_x$.

Let $t\in\mathbb R$ be the multiplier for
the normalization constraint, and let $\lambda\ge0$ be the multiplier for the
R\'enyi-power constraint. The Lagrangian is
\[
\mathcal L(q,t,\lambda)
=
\int r(y)q(y)\,dP_0
-
t\left(\int q\,dP_0-1\right)
-
\lambda\left(\int q^\alpha\,dP_0-\rho\right).
\]
Equivalently,
\[
\mathcal L(q,t,\lambda)
=
t+\lambda\rho
+
\int
\left[
(r(y)-t)q(y)-\lambda q(y)^\alpha
\right]P_0(dy).
\]
The non-negativity constraint $q\ge0$ is treated as a domain constraint.

The complementary slackness condition is
\[
\lambda
\left(
\int q^*(y)^\alpha\,P_0(dy)-\rho
\right)
=
0.
\]

We claim that in the present case one must have
\[
\lambda>0.
\]
Suppose instead that $\lambda=0$. Then the pointwise part of the Lagrangian is
\[
(r(y)-t)q(y),
\qquad q(y)\ge0.
\]
For the supremum over $q(y)\ge0$ to be finite, it is necessary that
\[
r(y)-t\le0
\qquad P_0\text{-a.s.}
\]
Thus, $t\ge r_x^\star$. On the other hand, since $q^*$ is feasible,
\[
\int q^*\,dP_0=1,
\]
so $q^*>0$ on a set of positive $P_0$-measure. On such a set, the pointwise
linear function $(r(y)-t)q$ can be maximized at a positive value of $q$ only if
\[
r(y)-t=0.
\]
Therefore $q^*$ must be supported on
\[
M_x=\{y:r(y)=r_x^\star\}.
\]
But in the case $\epsilon<-\log m_x$, no feasible policy can be supported
entirely on $M_x$. This is a contradiction. Hence, $\lambda>0$.

By complementary slackness, $\lambda>0$ implies
\[
\int q^*(y)^\alpha\,P_0(dy)=\rho.
\]
Equivalently, $\mathcal D_\alpha(\pi^*(\cdot\mid x)\|\pi_0(\cdot\mid x))
=
\epsilon$. Thus the optimum is attained on the boundary of the R\'enyi ball.

It remains to derive the form of $q^*$. For fixed $t$ and $\lambda>0$, the
Lagrangian is separable in $y$. For each $y$, one has to solve
\[
\max_{q\ge0}
\left\{
(r(y)-t)q-\lambda q^\alpha
\right\}.
\]
Define
$
\phi_y(q):=(r(y)-t)q-\lambda q^\alpha$ for $q\ge0$. Since $\alpha>1$ and $\lambda>0$, $\phi_y$ is strictly concave.

If $r(y)\le t$, then for every $q>0$,
\[
\phi_y'(q)
=
r(y)-t-\alpha\lambda q^{\alpha-1}
<0.
\]
Thus the maximizer is $
q^*(y)=0$. If $r(y)>t$, then
\[
\phi_y'(0+)=r(y)-t>0,
\]
so the maximizer is an interior point and satisfies
$
r(y)-t-\alpha\lambda q(y)^{\alpha-1}=0$. Therefore, $q^*(y)
=
\left(
\frac{r(y)-t}{\alpha\lambda}
\right)^{\frac1{\alpha-1}}$. Combining the two cases above yields
\[
q^*(y)
=
\left(
\frac{r(y)-t}{\alpha\lambda}
\right)_+^{\frac1{\alpha-1}}
\qquad P_0\text{-a.s.}
\]

Let $\beta:=\frac1{\alpha-1}$. Then,
$
q^*(y)
=
(\alpha\lambda)^{-\beta}(r(y)-t)_+^\beta$. The multiplicative constant is determined by the constraint $\int q^*\,dP_0=1$. Hence
\[
q^*(y)
=
\frac{(r(y)-t)_+^\beta}
{\int_{\mathcal Y}(r(y')-t)_+^\beta\,P_0(dy')}.
\]
Returning to the original notation,
\[
\frac{d\pi^*(\cdot\mid x)}{d\pi_0(\cdot\mid x)}(y)
=
\frac{(r(x,y)-t)_+^{\frac1{\alpha-1}}}
{\mathbb E_{Y\sim\pi_0(\cdot\mid x)}
[(r(x,Y)-t)_+^{\frac1{\alpha-1}}]}.
\]
The threshold $t$ is chosen so that $\mathcal D_\alpha(\pi^*(\cdot\mid x)\|\pi_0(\cdot\mid x))
=
\epsilon$. This completes the proof.
\end{proof}


\subsection{Proof of Theorem \ref{thm:rlhf}}
\begin{proof}
Fix $x\in\mathcal X$ and abbreviate $\pi(\cdot\mid x)$ and $\pi_0(\cdot\mid x)$
by $\pi$ and $\pi_0$, respectively. Let
\[
\bar r(x,y):= r(x,y)-\E_{Y\sim\pi_0}r(x,Y),
\qquad y\in\mathcal Y.
\]
Define the nonnegative random variable
$Z:=\frac{|\bar r(Y,x)|}{C}$ with $Y\sim\pi_0(\cdot\mid x)$. By the assumption $\|\bar r(x,Y)\|_{\psi_\theta}\le C$, under $Y\sim\pi_0(\cdot\mid x)$, we have
\[
\E_{\pi_0}\exp(Z^\theta)
=
\E_{\pi_0}\exp\!\left(\frac{|\bar r(Y)|^\theta}{C^\theta}\right)
\le 2.
\]

We now apply the decorrelation lemma (Lemma~\ref{de}) with $\mu=\pi$, $\nu=\pi_0$, and
the measurable function $r(\cdot):=|\bar r(\cdot,x)|/C$. This yields
\[
\E_{\pi}\!\left[\frac{|\bar r(Y)|}{C}\right]
\;\le\;
2^{1/\theta}\,\DD_{f_\theta}(\pi\|\pi_0)
\;+\;
\E_{\pi_0}\exp\!\left(\frac{|\bar r(Y)|^\theta}{C^\theta}\right)
\;\le\;
2^{1/\theta}\,\DD_{f_\theta}(\pi\|\pi_0)+2.
\]
Multiplying both sides by $C$ gives
\[
\E_{\pi}|\bar r(Y)|
\;\le\;
C\bigl(2^{1/\theta}\,\DD_{f_\theta}(\pi\|\pi_0)+2\bigr).
\]
Since $\E_\pi \bar r(Y)=\E_\pi r(x,Y)-\E_{\pi_0}r(x,Y)$, we conclude by
$|\E_\pi \bar r|\le \E_\pi|\bar r|$ that
\[
\bigl|\E_\pi r(x,Y)-\E_{\pi_0}r(x,Y)\bigr|
\;\le\;
C\bigl(2^{1/\theta}\,\DD_{f_\theta}(\pi\|\pi_0)+2\bigr).
\]
Finally, invoking Lemma~\ref{lem:key1} to upper bound
$\DD_{f_\theta}(\pi\|\pi_0)$ in terms of $\DD_\alpha(\pi\|\pi_0)$ yields the
displayed R\'enyi-based bound in the theorem.

For $\pi^*$ solving \eqref{optRENYI}, we have
$\DD_\alpha(\pi^*\|\pi_0)\le\epsilon$ for each $x$. Substituting this into the
preceding bound (and averaging over $x$ if desired) gives the stated guarantee
for $\pi^*$.
\end{proof}

\subsection{Proof of Theorem \ref{renyiregret}}

\begin{proof}
By Assumption \ref{ass:reward-structure}, we need to bound 
\[
\DD_\alpha (r_n(x)\|r(x)).
\]
This reduces to the following lemma by setting $r_n(x)\sim\pi_n=:\nu$ and $r(x)\sim\pi_0=:\mu$.
\begin{lemma}[\Renyi divergence of the maximum of i.i.d.\ draws]\label{DDD} 
Let $\mu$ and $\nu$ be any two univariate probability distributions. Assume $\alpha>1$. Let $X$ and $\{X_i\}_{i=1}^n \overset{\iid}{\sim} \mu$ be random variables, and put $R_n=\max_{i\in [n]} X_i \sim \nu$. Then
\[
D_\alpha(\nu \| \mu) \leq \frac{1}{\alpha-1}\log \left( \frac{n^\alpha}{\alpha(n-1)+1} \right)
\]
and $D_{KL}(\nu \| \mu) \leq \log n-\frac{n-1}{n}$.
%(KL case is from Theorem 1 in \cite{mroueh2025information}).
%In particular, for fixed $\alpha$, $D_\alpha(\nu \| \mu) = O(\log n)$ and grows sublinearly in $n$.
\end{lemma}

\begin{proof}
The KL case is given by Theorem 1 in \cite{mroueh2025information} and can also be proved by taking the limit $\alpha \to 1$, so here we prove the \Renyi case: $\alpha>1$.

Let $Q$ be the left-continuous quantile function of $X$. For $U$ and $\{U_i\}_{i=1}^n \overset{\iid}{\sim}\mathrm{U}[0,1]$,
\[
R\stackrel{d}=Q(U),
\qquad
R_n\stackrel{d}=Q(U_{(n)}),
\]
where $U_{(n)}:=\max_{1\le i\le n}U_i$.
By data processing inequality for R\'enyi divergence,
\[
D_\alpha(R_n\|R)\le D_\alpha(U_{(n)}\|U).
\]
Now $U_{(n)}$ has density $n u^{n-1}$ on $[0,1]$, so
\[
D_\alpha(R_n\|R)\le D_\alpha(U_{(n)}\|U)
=
\frac{1}{\alpha-1}\log\int_0^1 \paren{n u^{n-1}}^\alpha\,du
=
\frac{1}{\alpha-1}\log\frac{n^\alpha}{\alpha(n-1)+1}.
\]
By L'Hôpital's Rule, $\lim_{\alpha \to 1}\frac{1}{\alpha-1}\log\frac{n^\alpha}{\alpha(n-1)+1}=\log n-\frac{n-1}{n}$ and we have $D_{KL}(\nu \| \mu) \leq \log n-\frac{n-1}{n}$.
%Let the support of $\mu$ be $\{t_k\}$. Then
%\[
%\exp \big( (\alpha-1)D_\alpha(\nu\|\mu) \big)
%= \int \left( \frac{d(F^n(t_k))}{dF(t_k)}\right)^\alpha d F(x),
%\]
%where $F$ is the CDF of $\mu$. Expanding over the support points $\{t_k\}$ gives
%\[
%\frac{d(F^n(t_k))}{dF(t_k)}
%= \frac{F^n(t_k) - F^n(t_{k-1})}{F(t_k) - F(t_{k-1})}.
%\]
%We use the following inequality: for $a>b$,
%\begin{equation}\label{eqab}
%(a-b) \left( \frac{a^n - b^n}{a-b} \right)^\alpha
%\leq \int_b^a n^\alpha t^{\alpha(n-1)} \, dt.
%\end{equation}
%Let $g(t) = n t^{n-1}$ on $[b,a]$. Since $\alpha > 1$, the function
%$\varphi(x) = x^\alpha$ is convex on $[0,\infty)$. By Jensen's inequality,
%\[
%\varphi\!\left(\frac{1}{a-b}\int_b^a g(t)\,dt\right)
%\leq \frac{1}{a-b}\int_b^a \varphi(g(t))\,dt.
%\]
%Multiplying both sides by $(a-b)$ yields
%\[
%(a-b)\left(\frac{1}{a-b}\int_b^a g(t)\,dt\right)^\alpha
%\leq \int_b^a g^\alpha(t)dt.
%\]
%Note that
%\[
%\frac{1}{a-b}\int_b^a g(t)\,dt
%= \frac{1}{a-b}\int_b^a n t^{n-1}\,dt
%= \frac{a^n - b^n}{a-b},
%\]
%and $g(t)^\alpha = n^\alpha t^{\alpha(n-1)}$. Substituting these expressions yields
%\eqref{eqab}. Therefore,
%\[
%\begin{aligned}
%\exp \big( (\alpha-1)D_\alpha(\nu \| \mu) \big)
%&= \sum_k \left( \frac{F^n(t_k) - F^n(t_{k-1})}{F(t_k)-F(t_{k-1})} \right)^\alpha \big(F(t_k)-F(t_{k-1})\big) \\
%&\leq \sum_k \int_{t_{k-1}}^{t_k} n^\alpha t^{\alpha(n-1)} \, dt 
%= \int_0^1 n^\alpha t^{\alpha(n-1)} \, dt
%= \frac{n^\alpha}{\alpha(n-1)+1}.
%\end{aligned}
%\]
\end{proof}

Follows from Assumption~\ref{ass:reward-structure} and another application of data processing:
\[
D_\alpha(\pi_n\|\pi_0)\le D_\alpha(R_n\|R).
\]
Now, in this setting, since $\frac{1}{\alpha-1}\log \left( \frac{n^\alpha}{\alpha(n-1)+1} \right)\leq \log n$ for $\alpha>1$ and $n\ge 1$, we have 
\[
\DD_\alpha(\pi_n \|\pi_0) \leq \frac{1}{\alpha-1}\log \left( \frac{n^\alpha}{\alpha(n-1)+1} \right)\leq \log n
\]
So if $n\leq \exp(\varepsilon)$, then $\DD_\alpha(\pi_n\|\pi_0)\leq \varepsilon$.

Applying the decorrelation lemma (Lemma \ref{de}) to $\bar r(x,Y):= |r(x,Y)-\E_{Y\sim\pi_0}r(x,Y)|/C,~~y\in\mathcal Y$, similar to the proof of Theorem \ref{thm:rlhf}, we have
\[
\mathbb{E}_{\pi_n(\cdot|x)} r - \mathbb{E}_{\pi_0(\cdot|x)} r \le |\mathbb{E}_{\pi_n(\cdot|x)} r - \mathbb{E}_{\pi_0(\cdot|x)} r|
\leq C \left( 2^{1/\theta} \left[ \DD_\alpha(\pi_n \| \pi_0) +C_{\alpha,\theta}\right]^{1/\theta} + 2 \right).
\]
from $\DD_\alpha(\pi_n\|\pi_0)\leq \varepsilon$. This completes the proof.
\end{proof}





% \subsection{Proof of Lemma \ref{SGLD1}}
% We first state a lemma for the $\alpha$-R\'enyi divergence between two $d$-dimensional Gaussian distributions.

% \begin{lemma}
% Let $p=\mathcal{N}(\mu_1,\sigma^2 I_d)$ and $q=\mathcal{N}(\mu_2,\sigma^2 I_d)$ be $d$-dimensional Gaussian distributions with identical isotropic covariance $\sigma^2 I_d$. Then for any $\alpha>0$, $\alpha\neq 1$, the $\alpha$ R\'enyi divergence is
% \[
% \DD_\alpha(p\|q)=\frac{\alpha}{2\sigma^2}\,\|\mu_1-\mu_2\|^2.
% \]
% \end{lemma}

% \begin{proof}
% Recall the definition of the R\'enyi divergence,
% \[
% \DD_\alpha(p\|q)
% =\frac{1}{\alpha-1}\log\!\int p(x)^\alpha q(x)^{1-\alpha}\,dx.
% \]
% The densities of $p$ and $q$ are
% $
% p(x)=\frac{1}{(2\pi\sigma^2)^{d/2}}
% \exp\!\Bigl(-\frac{\|x-\mu_1\|^2}{2\sigma^2}\Bigr)$ and $q(x)=\frac{1}{(2\pi\sigma^2)^{d/2}}
% \exp\!\Bigl(-\frac{\|x-\mu_2\|^2}{2\sigma^2}\Bigr)$. Therefore,
% \[
% p(x)^\alpha q(x)^{1-\alpha}
% =(2\pi\sigma^2)^{-d/2}
% \exp\!\Bigl(
% -\frac{1}{2\sigma^2}\bigl(
% \alpha\|x-\mu_1\|^2+(1-\alpha)\|x-\mu_2\|^2
% \bigr)
% \Bigr).
% \]

% We now complete the square in the exponent. Let $m=\alpha\mu_1+(1-\alpha)\mu_2$. A direct expansion shows that
% \[
% \alpha\|x-\mu_1\|^2+(1-\alpha)\|x-\mu_2\|^2
% =\|x-m\|^2+\alpha(1-\alpha)\|\mu_1-\mu_2\|^2.
% \]
% Thus,
% \[
% p(x)^\alpha q(x)^{1-\alpha}
% =(2\pi\sigma^2)^{-d/2}
% \exp\!\Bigl(-\frac{\|x-m\|^2}{2\sigma^2}\Bigr)
% \exp\!\Bigl(-\frac{\alpha(1-\alpha)\|\mu_1-\mu_2\|^2}{2\sigma^2}\Bigr).
% \]

% Integrating over $\mathbb{R}^d$, the first exponential integrates to $(2\pi\sigma^2)^{d/2}$, since it is the normalization constant of a Gaussian distribution with mean $m$ and covariance $\sigma^2 I_d$. Hence,
% \[
% \int p(x)^\alpha q(x)^{1-\alpha}\,dx
% =\exp\!\Bigl(-\frac{\alpha(1-\alpha)\|\mu_1-\mu_2\|^2}{2\sigma^2}\Bigr).
% \]

% Taking the logarithm and dividing by $\alpha-1$ yields
% \[
% \DD_\alpha(p\|q)
% =\frac{1}{\alpha-1}\Bigl(
% -\frac{\alpha(1-\alpha)\|\mu_1-\mu_2\|^2}{2\sigma^2}
% \Bigr)
% =\frac{\alpha}{2\sigma^2}\|\mu_1-\mu_2\|^2.
% \]
% \end{proof}



% Let \(K_\theta:=2^{(1/\theta-1)_+}\). Denote the joint distribution of $(X, Y)$ by \(P_{X,Y}\). By independence, the distributions of \(X + \varepsilon\) and \(Y + \varepsilon\) can be written as
% \[
% P_{X+\varepsilon} = \int_{x,y} \mathcal{N}(x, \sigma^2 I) \, \mathrm{d}P_{X,Y}(x, y)
% \quad \text{and} \quad
% P_{Y+\varepsilon} = \int_{x,y} \mathcal{N}(y, \sigma^2 I) \, \mathrm{d}P_{X,Y}(x, y),
% \]
% respectively, where \(\mathcal{N}(x, \sigma^2 I)\) is the Gaussian distribution with mean \(x\) and covariance \(\sigma^2 I\). Since $f_\theta(x)=x\log^{\frac{1}{\theta}}(x+A)$ is convex (Lemma~\ref{lemma:convex}), $\DD_{f_\theta}(P \| Q)$ is jointly convex. Using this observation, we can write
% \[
% \begin{aligned}
% &~~~~ \mathcal{D}_{f_\theta}\left(P_{X+\varepsilon} \| P_{Y+\varepsilon}\right)
% = \mathcal{D}_{f_\theta}\left( \int_{x,y} \mathcal{N}(x, \sigma^2 I) \, \mathrm{d}P_{X,Y}(x, y) \bigg\| \int_{x,y} \mathcal{N}(y, \sigma^2 I) \, \mathrm{d}P_{X,Y}(x, y) \right) \\
% &\leq \int_{x,y} \mathcal{D}_{f_\theta}\left( \mathcal{N}(x, \sigma^2 I) \| \mathcal{N}(y, \sigma^2 I) \right) \, \mathrm{d}P_{X,Y}(x, y) = \mathbb{E}_{X,Y} \left[ \mathcal{D}_{f_\theta}\left( \mathcal{N}(X, \sigma^2 I) \| \mathcal{N}(Y, \sigma^2 I) \right) \right] \\
% &\leq \mathbb{E}_{X,Y} \left[
% \left(\mathcal{D}_\alpha\left( \mathcal{N}(X, \sigma^2 I) \| \mathcal{N}(Y, \sigma^2 I) \right)+C_{\alpha,\theta}\right)^\I
% \right] \\
% &=
% \mathbb{E}_{X,Y}\left[
% \left(\frac{\alpha}{2\sigma^2}\|X-Y\|^2+C_{\alpha,\theta}\right)^\I
% \right]\\
% &\leq
% K_\theta\left\{
% \left(\frac{\alpha}{2\sigma^2}\right)^\I
% \E \|X - Y\|^{\frac{2}{\theta}}
% +C_{\alpha,\theta}^\I
% \right\},
% \end{aligned}
% \]
% where the second line uses Jensen's inequality and the joint convexity of \(\mathcal{D}_{f_\theta}(\cdot \| \cdot)\)  (Page 120 of \cite{Yihong25}, Theorem 7.5(b)) in its arguments. The next two lines use Lemma~\ref{lem:key1}, the Gaussian R\'enyi divergence formula above, and \((u+v)^p\le 2^{(p-1)_+}(u^p+v^p)\) with \(p=1/\theta\). 


% %\section{Asymptotic analysis and examples of results for sub-Weibull random variables}\label{asym-Weibull}
% \subsection{Proof of Theorem \ref{eq:SGD0}}
% \begin{proof}
% Let $S$ be a random dataset sampled from $P_S$, and let $W_{1:T}:=(W_1,\dots,W_T)$ be the sequence of model parameters generated by a learning algorithm conditioned on $S$. Let $\widehat W = F(W_{1:T})$ be any measurable output
% (e.g.\ $\widehat W=W_T$ or $\widehat W=T^{-1}\sum_{t=1}^T W_t$). Recall the shifted-log
% divergence $f_\theta(x)=x\log^{1/\theta}(x+A)$ and the corresponding $f_\theta$--mutual information:
% \[
% I_{f_\theta}(F(W_{1:T}),S)
% = \E_{P_S}\,\DD_{f_\theta}\!\left(P_{F(W_{1:T})\mid S}\,\Big\|\,P_{F(W_{1:T})}\right).
% \]
% By the data processing inequality for $f$--divergence (c.f. \citealp{Yihong25}, p.~120),
% \[
% I_{f_\theta}(F(W_{1:T}),S)
% \;\le\;
% \E_{P_S}\,\DD_{f_\theta}\!\left(P_{W_{1:T}\mid S}\,\Big\|\,P_{W_{1:T}}\right).
% \]
% where the conditional $f$-divergence of $P_{W_{1:T}\mid S}$ with respect to the marginal $P_{W_{1:T}}$ is defined as
% $\DD_{f_\theta}(P_{W_{1:T}\mid S}\,\|\,P_{W_{1:T}})
% =\int f_\theta\Big(\frac{dP_{W_{1:T}\mid S}}{dP_{W_{1:T}}}(w_{1:T})\Big)\, dP_{W_{1:T}}(w_{1:T})$. Define $L_t=\frac{dP_{W_t\mid W_{1:t-1},S}}{dP_{W_t\mid W_{1:t-1}}}$. Using the chain rule for the Radon--Nikodym derivative,
% \[
% \log \left ( \frac{dP_{W_{1:T}\mid S}}{dP_{W_{1:T}}}+A \right )
% =\log \left ( \prod_t L_t  +A\right )
% \leq \log \left ( (1+A) \prod_t (L_t+A)\right)=\log(1+A)+\sum_t \log(L_t+A)
% .
% \]

% If $1/\theta \ge 1$ and $x\mapsto x^{1/\theta}$ is convex, Jensen's inequality yields
% $
% (\sum_{t=1}^m a_t)^{1/\theta}
% \;\le\; m^{1/\theta-1}\sum_{t=1}^m a_t^{1/\theta}
% $
% for nonnegative $\{a_t\}_{t=1}^m$. If $0<1/\theta\le 1$, the map $x\mapsto x^{1/\theta}$ is subadditive, so
% $(\sum_{t=1}^m a_j)^{1/\theta}\le \sum_{t=1}^m a_j^{1/\theta}$. These two cases give
% \begin{equation}\label{eq:subadditive}
%     \paren{\sum_{j=1}^m a_j}^{1/\theta}
% \le
% m^{(\I-1)_+}\sum_{j=1}^m a_j^{1/\theta}.
% \end{equation}
% Applying this to $m=T+1$
% $a_t=\log(L_t+A)=\log \left ( \frac{dP_{W_t\mid W_{1:t-1},S}}{dP_{W_t\mid W_{1:t-1}}}+A\right )\geq 0$ for $t\leq T$ and $a_{T+1}=\log(1+A)$ , we obtain
% \begin{align}\label{eq:IFtheta-upper}
%  &   I_{f_\theta}(S; W_{1:T}) 
%     =\E_{P_S\otimes P_{W_{1:T}}}\left[ \frac{dP_{W_{1:T}|S}}{dP_{W_{1:T}}} \log^{1/\theta} \left ( \frac{dP_{W_{1:T}\mid S}}{dP_{W_{1:T}}}+A \right ) \right] \nonumber\\
%     &\leq \E_{P_S\otimes P_{W_{1:T}}}\left[ \frac{dP_{W_{1:T}|S}}{dP_{W_{1:T}}} \left ( \sum\log(L_t+A)+ \log(1+A) \right )^{1/\theta} \right] \nonumber\\
%     &\le (T+1)^{(\I-1)_+}\left [ \log^{1/\theta}(1+A)+\sum_{t=1}^T\E_{P_S\otimes P_{W_{1:T}}} \left[ \frac{dP_{W_{1:T}|S}}{dP_{W_{1:T}}} \log^{1/\theta}(L_t+A) \right] \right ] \nonumber\\
%     &= (T+1)^{(\I-1)_+}\left [ \log^{1/\theta}(1+A)+\sum_{t=1}^T \E_{S, W_{1:t-1}} \DD_{f_\theta}(P_{W_t | W_{1:t-1}, S} \| P_{W_t | W_{1:t-1}})\right ].
% \end{align}

% Next, we characterize conditional distributions of $W_t$ under SGLD.
% The update rule is
% \[
% W_{t+1}=W_t-\eta_t g(W_t,B_t)+\epsilon_t,
% \qquad
% \epsilon_t\sim\mathcal N(0,\sigma_t^2 I_d).
% \]

% \paragraph{One-step coupling.}
% Conditioned on $\mathcal{F}_{t-1}:=\sigma(W_{1:t-1},S,B_{t-1})$,
% $W_t\sim\mathcal N\bigl(W_{t-1}-\eta_{t-1}g(W_{t-1},B_{t-1}),\,\sigma_{t-1}^2 I_d\bigr)$,
% so we may write
% \[
% W_t=X_t+\varepsilon_{t-1},
% \qquad
% X_t:=W_{t-1}-\eta_{t-1}g(W_{t-1},B_{t-1}),
% \]
% with $\varepsilon_{t-1}\perp\!\!\!\perp X_t$ and
% $P_{X_t+\varepsilon_{t-1}\mid W_{1:t-1},S}=P_{W_t\mid W_{1:t-1},S}$.

% To couple this with $P_{W_t\mid W_{1:t-1}}$, let $S'$ be a conditionally
% independent copy of $S$ given $W_{1:t-1}$: that is,
% $S'\mid W_{1:t-1}\sim P_{S\mid W_{1:t-1}}$ and $S'\perp\!\!\!\perp S\mid W_{1:t-1}$. Note that $S$ is not an independent fresh dataset unconditionally.
% Let $J'_{t-1}$ be an independent copy of the batch index $J_{t-1}$, and define
% $B'_{t-1}:=Z'_{J'_{t-1}}\subset S'$ (the mini-batch drawn from $S'$). Set
% \[
% Y_t:=W_{t-1}-\eta_{t-1}g(W_{t-1},B'_{t-1}).
% \]
% By construction, $B'_{t-1}\mid W_{1:t-1}$ has the same law as
% $B_{t-1}\mid W_{1:t-1}$ (since $S'\mid W_{1:t-1}\stackrel{d}{=}S\mid W_{1:t-1}$
% and $J'_{t-1}\stackrel{d}{=}J_{t-1}$), so
% $P_{Y_t+\varepsilon_{t-1}\mid W_{1:t-1}}=P_{W_t\mid W_{1:t-1}}$. Since \(S'\perp S\mid W_{1:t-1}\), we also have
% \[
% P_{Y_t+\varepsilon_{t-1}\mid W_{1:t-1},S}
% =
% P_{Y_t+\varepsilon_{t-1}\mid W_{1:t-1}}
% =
% P_{W_t\mid W_{1:t-1}}.
% \]
% %Moreover,
% %$B'_{t-1}\perp\!\!\!\perp S\mid W_{1:t-1}$, so Lemma~\ref{SGLD1} applied to the
% %pair $(X_t,Y_t)$ yields a non-trivial bound on
% %$\mathcal{D}_{f_\theta}(P_{W_t\mid W_{1:t-1},S}\,\|\,P_{W_t\mid W_{1:t-1}})$.




% Applying Lemma~\ref{SGLD1} to the pair $(X_t,Y_t)$  yields
% \begin{align*}
% &D_{f_\theta}\!\left(
% P_{W_t\mid W_{1:t-1},S}\,\big\|\,P_{W_t\mid W_{1:t-1}}
% \right) \le
% K_\theta\left[
% \Bigl(\frac{\alpha}{2\sigma_{t-1}^2}\Bigr)^{1/\theta}
% \E\bigl[\|X_t-Y_t\|^{2/\theta}\,\big|\,W_{1:t-1},S\bigr]
% +C_{\alpha,\theta}^{1/\theta}
% \right]
% \\
% &=
% K_\theta\left[
% \Bigl(\frac{\alpha\eta_{t-1}^2}{2\sigma_{t-1}^2}\Bigr)^{1/\theta}
% \E\bigl[\|g(W_{t-1},B_{t-1})-g(W_{t-1},B'_{t-1})\|^{2/\theta}\,\big|\,W_{1:t-1},S\bigr]
% +C_{\alpha,\theta}^{1/\theta}
% \right].
% \end{align*}
% For $t=1$, $W_1$ is independent of $S$, so $\DD_{f_\theta}(P_{W_1|S} \| P_{W_1}) = \log^{1/\theta}(1+A)$. Then the sum starts at $t=2$. Thus by \eqref{eq:IFtheta-upper},
% {\small
% \begin{align*}
% &~~~~I_{f_\theta}(F(W_{1:T}),S)
% \le (T+1)^{(\I-1)_+}\left [ \log^{1/\theta}(1+A)+\sum_{t=1}^T \E_{S, W_{1:t-1}} \DD_{f_\theta}(P_{W_t | W_{1:t-1}, S} \| P_{W_t | W_{1:t-1}})\right ] \\
% & \le (T+1)^{(\I-1)_+}
%     \left[ 2\log^{1/\theta}(1+A) + \sum_{t=2}^T 
%         \left(
%             K_\theta\left[
%             \left(\frac{\alpha\eta_{t-1}^2}{2\sigma_{t-1}^2}\right)^\I
%                 \E\|g(W_{t-1},B_{t-1})-g(W_{t-1},B'_{t-1})\|^{\frac{2}{\theta}}  + C_{\alpha,\theta}^{1/\theta}
%             \right]
%         \right)
% \right].
% \end{align*}
%     }


% We now bound the stochastic gradient difference with $\bar g(w)=\mathbb{E}_Z g(w,Z)$. By the triangle inequality,
% \[
% \|g(W_t,B_t)-g(W_t,B'_t)\|
% \le
% \|g(W_t,B_t)-\bar g(W_t)\|
% +
% \|g(W_t,B'_t)-\bar g(W_t)\|.
% \]
% Using $C_r$-inequality: $(x+y)^k \le 2^{k-1}(x^k+y^k)$ for $k=\frac{2}{\theta}\ge 1$, we obtain
% \[
% \E\|g(W_t,B_t)-g(W_t,B'_t)\|^{\frac{2}{\theta}}
% \le 2^{\frac{2}{\theta}}G_t
% ~\text{with}~G_t:=\E\|g(W_t,B_t)-\bar g(W_t)\|^{\frac{2}{\theta}}.
% \]

% Substituting this bound into the previous display and invoking
% Theorem~\ref{gen} to convert $I_{f_\theta}(F(W_{1:T}),S)$ into a bound on the expected generalization error, we conclude that
% \[
% \E\,|\gen(\hat W,S)|
% \le
% \frac{M_\theta v_\theta}{\sqrt n}
% \left\{
% 4
% +
% 2^{1/\theta}\,
% (T+1)^{(\I-1)_+}
% \left[
% 2\log^{1/\theta}(1+A)
% +
% \sum_{t=1}^{T-1}
% K_\theta\left(
% C_{\alpha,\theta}^{1/\theta}
% +
% \left(\frac{2\alpha\eta_t^2}{\sigma_t^2}\right)^{1/\theta}G_t
% \right)
% \right]
% \right\},
% \]
% which completes the proof.
% \end{proof}




% \section{Stochastic Optimization with Heavy-tailed Losses}\label{se:SGDL}


% Beyond empirical risk minimization, most information-theoretic generalization
% bounds for stochastic optimization rely on light-tailed assumptions on the loss
% and its gradients; see, e.g., \citet[Section~8]{hellstrom2025generalization} for
% a survey. In noisy iterative methods such as stochastic gradient algorithms,
% however, both losses and stochastic gradients may exhibit heavy tails \citep{raj2023algorithmic, madden2024high, fatkhullin2025can, sun2025distributed}.

% In this section, we apply our information-theoretic framework to stochastic
% gradient descent, with a particular focus on stochastic gradient Langevin
% dynamics (SGLD; \citealp{welling2011bayesian}). Our starting point is the
% analysis-by-perturbation viewpoint of \citet{algo1}, which controls generalization
% via information accumulated along the optimization trajectory. Under
% sub-Gaussian assumptions, this approach yields sub-Gaussian generalization
% bounds. We extend this analysis to sub-Weibull losses and heavy-tailed gradient noise, obtaining analogous
% bounds whose constants depend explicitly on the tail parameter $\theta$ and on
% pathwise gradient variability.

% \paragraph{SGLD dynamics.}
% Consider SGLD in $\mathbb{R}^d$ with step sizes $\{\eta_t\}_{t\ge 1}$ and
% injected Gaussian noises $\varepsilon_t\sim\mathcal{N}(0,\sigma_t^2 I_d)$,
% independent across $t$ and independent of $S$ and $W_1$:
% \begin{equation}\label{SGLDupdate}
% W_{t+1}=W_t-\eta_t\,g(W_t,B_t)+\varepsilon_t,
% \qquad
% g(w,B_t):=\frac{1}{|B_t|}\sum_{Z\in B_t} g(w,Z),
% \end{equation}
% where $B_t=(Z_{I_{t,1}},\dots,Z_{I_{t,b}})$ is the mini-batch used at
% iteration $t$. The batch index set $J_t=\{I_{t,1},\dots,I_{t,b}\}\subset[n]$
% is sampled (with or without replacement) according to a fixed rule that is
% \emph{independent of $S$, $W_1$, $\varepsilon_{1:T}$, and $J_{1:t-1}$}. In
% particular, $J_t$ is independent of the past, so given $W_{1:t}$, the
% mini-batch index $J_t$ has its prescribed sampling distribution (and the data
% $Z_{J_t}\subset S$ is then determined by $J_t$ and $S$).

% Under these conditions, the unbiasedness convention
% \[
% \mathbb{E}[g(w,B_t)] = \bar g(w) := \mathbb{E}_{Z\sim\mu}\,g(w,Z), 
% \qquad \forall\,w\in\mathcal W,
% \]
% holds when $S\sim\mu^{\otimes n}$ (with the expectation jointly over $J_t$
% and the data). The algorithm output is any measurable function
% $\widehat W=F(W_{1:T})$, and we assume $W_1$ is independent of $S$.
% Following the convention of \cite{algo1}, we center the mini-batch
% gradient at the population gradient
% \(\bar g(w)=\mathbb E_{Z\sim\mu}g(w,Z)\), rather than at the conditional mean
% of the stochastic gradient given the algorithmic state. Thus \(G_t\) should be
% viewed as a population-centered local gradient fluctuation moment, not as a
% conditional variance. In particular, our proof does not require
% \(\mathbb E[g(W_t,B_t)\mid W_t]=\bar g(W_t)\); it only uses the moment
% assumption defining \(G_t\) and the coupling/triangle-inequality argument below.


% % Theorem~\ref{eq:SGD0} separates the expected generalization gap into a
% % \emph{distributional} term and an \emph{algorithmic/path} term.
% % The distributional term is the sub-Weibull scale $v_\theta$, which captures the
% % uniform tail-heaviness of the (centered) loss.

% % The path term aggregates stability contributions along the SGLD trajectory:
% % \[
% % \sum_{t=1}^{T-1}\Bigl(\frac{\eta_t^2}{\sigma_t^2}\Bigr)^{1/\theta} G_t,
% % \qquad
% % G_t \coloneqq \E\bigl\|g(W_t,B_t)-\bar g(W_t)\bigr\|^{2/\theta}.
% % \]
% % This expression makes the roles of the two sources of randomness explicit.
% % The heavy-tailed component comes from the \emph{stochastic-gradient noise}
% % $g(W_t,B_t)-\bar g(W_t)$, whose burstiness is quantified by the moment $G_t$
% % and the tail index $\theta$ (smaller $\theta$ corresponds to heavier tails and
% % places more weight on rare but large gradient fluctuations).
% % In contrast, the injected Gaussian noise level $\sigma_t$ is a \emph{smoothing}
% % mechanism: convolving each update with $\mathcal N(0,\sigma_t^2 I)$ reduces the
% % distinguishability between neighboring trajectories (e.g., under neighboring
% % datasets) by smoothing the iterate distribution, thereby limiting how much
% % data-dependence can accumulate along the path.

% % The tradeoff is governed by the noise-normalized drift budget above:
% % increasing $\sigma_t$ decreases the per-step contribution, while increasing
% % $\eta_t$ increases it.
% % Compared to the sub-Gaussian baseline $\theta=2$, where the dependence is
% % $(\eta_t^2/\sigma_t^2)^{1/2}$, heavier tails (smaller $\theta$) increase the
% % exponent $1/\theta$ and thus make the bound more sensitive to the effective ratio
% % $\eta_t^2/\sigma_t^2$; in particular, for $\theta<1$ this dependence becomes
% % super-linear.
% % Thus, Gaussian injection need not eliminate heavy tails in the stochastic
% % gradients; rather, it provides stabilizing smoothing whose strength must be
% % calibrated against tail-heaviness through $\theta$ and $G_t$.
% % Simulation studies for linear regression are given in Appendix~\ref{ap:SGDL}.


% % When $\theta=2$, the sub-Weibull norm coincides with the usual sub-Gaussian
% % scale and the bound recovers the result by \citet{algo1}, discussed in detail later. As $\theta$ decreases below $1$, the exponent
% % $1/\theta$ amplifies the cumulative path term, reflecting the fact that rare
% % gradient outliers can dominate stability unless the effective per-step
% % noise-to-signal ratio $\eta_t^2/\sigma_t^2$ is decreased accordingly. 



% % When $\theta=2$, the loss is sub-Gaussian and $G_t=\E\|g(W_t,B_t)-\bar g(W_t)\|$.
% % Specializing Theorem~\ref{eq:SGD0} to $\theta=2$ yields, up to constants,
% % \[
% % \E\bigl|\gen(\widehat W,S)\bigr|
% % \;\lesssim\;
% % \frac{v_2}{\sqrt n}\left[
% % 1
% % +\frac{1}{\sqrt{T}}\sum_{t=1}^{T-1}
% % \left(
% % B_{\alpha,2}
% % +\sqrt{\frac{\eta_t^2}{\sigma_t^2}}\,
% % \E\bigl\|g(W_t,B_t)-\bar g(W_t)\bigr\|
% % \right)
% % \right],
% % \]
% % where the $\sqrt{\eta_t^2/\sigma_t^2}=|\eta_t|/\sigma_t$ factor is the
% % sub-Gaussian specialization of $(\eta_t^2/\sigma_t^2)^{1/\theta}$.

% % By contrast, \citet[Proposition~3]{algo1} (for $\theta=2$) obtain a KL-based
% % path bound of the schematic form
% % \[
% % O\!\left(
% % \sqrt{\frac{1}{n}\sum_{t=1}^{T-1}\frac{\eta_t^2}{\sigma_t^2}\,
% % \E\bigl\|g(W_t,B_t)-\bar g(W_t)\bigr\|^{2}}
% % \right).
% % \]
% % Both bounds capture the same stability mechanism: generalization improves with
% % larger injected noise $\sigma_t$ and degrades with larger step sizes $\eta_t$ and
% % more variable stochastic gradients; but they differ in form for two reasons.
% % First, our result is stated for an arbitrary post-processed output
% % $\widehat W=F(W_{1:T})$, whereas \citet{algo1} focus on the final iterate.
% % Second, our heavy-tail analysis proceeds by controlling stepwise
% % $f_\theta$-divergence increments and then upper bounding them by R\'enyi
% % divergence; unlike KL, R\'enyi divergences do not admit the same additive chain
% % rule along the optimization path, which naturally leads to an additive
% % per-step accumulation before applying outer norm inequalities.
% % For $\theta=2$, one can relate the two forms using Cauchy--Schwarz and Jensen's inequalities:
% % \[
% % \frac{1}{\sqrt{T}}\sum_{t=1}^{T-1}\frac{|\eta_t|}{\sigma_t}\,\E\|X_t\|
% % \;\le\;
% % \sqrt{\sum_{t=1}^{T-1}\frac{\eta_t^2}{\sigma_t^2}\,(\E\|X_t\|)^2}
% % \;\le\;
% % \sqrt{\sum_{t=1}^{T-1}\frac{\eta_t^2}{\sigma_t^2}\,\E\|X_t\|^{2}},
% % \]
% % where $X_t=g(W_t,B_t)-\bar g(W_t)$.
% % Thus, in the sub-Gaussian regime the present bound is comparable to the KL
% % path bound, while our formulation continues to apply in genuinely heavy-tailed settings where MGF-based KL techniques break down.

% Theorem~\ref{eq:SGD0} decomposes the expected generalization gap into a
% \emph{distributional} component and an \emph{algorithmic/path} component.
% The distributional component is the sub-Weibull scale $v_\theta$, which
% uniformly quantifies the tail-heaviness of the centered loss.

% The path component aggregates stability contributions along the SGLD trajectory
% through the step sizes $\eta_t$, injected noise levels $\sigma_t$, and the
% stochastic-gradient noise:
% \[
% \sum_{t=1}^{T-1}\Bigl(\frac{\eta_t^2}{\sigma_t^2}\Bigr)^{1/\theta} G_t,
% \qquad
% G_t \coloneqq \EE\bigl\|g(W_t,B_t)-\bar g(W_t)\bigr\|^{2/\theta}.
% \]
% This term makes the roles of the two sources of randomness explicit.
% Heavy tails arise from the \emph{minibatch (stochastic-gradient) noise}
% $g(W_t,B_t)-\bar g(W_t)$, whose burstiness is captured by the moment $G_t$ and the
% tail index $\theta$: smaller $\theta$ corresponds to heavier tails and increases
% the influence of rare but large gradient fluctuations.
% By contrast, the injected Gaussian noise is \emph{algorithmic} and plays the role
% of a smoothing mechanism: adding $\mathcal N(0,\sigma_t^2 I)$ at each step blurs
% the iterate distribution, reducing the distinguishability between neighboring
% trajectories (e.g., under neighboring datasets) and thereby limiting how much
% data-dependence can accumulate along the path.

% The resulting tradeoff is governed by the noise-normalized drift budget above:
% larger $\sigma_t$ suppresses each per-step contribution, while larger $\eta_t$
% amplifies it.
% Relative to the sub-Gaussian baseline $\theta=2$, where the dependence is
% $(\eta_t^2/\sigma_t^2)^{1/2}$, heavier tails (smaller $\theta$) increase the
% exponent $1/\theta$ and make the bound more sensitive to the effective ratio
% $\eta_t^2/\sigma_t^2$; in particular, for $\theta<1$ this dependence becomes
% super-linear.
% Thus, Gaussian injection need not make the stochastic gradients light-tailed;
% rather, it provides stabilizing smoothing whose strength must be calibrated
% against tail-heaviness by $\theta$ and $G_t$.

% This interpretation is consistent with the sub-Gaussian specialization.
% When $\theta=2$, the sub-Weibull norm reduces to the usual sub-Gaussian scale and
% $G_t=\E\|g(W_t,B_t)-\bar g(W_t)\|$.
% Specializing Theorem~\ref{eq:SGD0} to $\theta=2$ yields, up to constants,
% \[
% \E\bigl|\gen(\widehat W,S)\bigr|
% \;\lesssim\;
% \frac{v_2}{\sqrt n}\left[
% 1
% +\frac{1}{\sqrt{T}}\sum_{t=1}^{T-1}
% \left(
% B_{\alpha,2}
% +\frac{|\eta_t|}{\sigma_t}\,
% \E\bigl\|g(W_t,B_t)-\bar g(W_t)\bigr\|
% \right)
% \right],
% \]
% where $\eta_t/\sigma_t$ is exactly the $\theta=2$ instance of
% $(\eta_t^2/\sigma_t^2)^{1/\theta}$.

% By contrast, \citet[Proposition~3]{algo1} obtain (for $\theta=2$) a KL-based path
% bound of as the form
% \[
% O\!\left(
% \sqrt{\frac{1}{n}\sum_{t=1}^{T-1}\frac{\eta_t^2}{\sigma_t^2}\,
% \E\bigl\|g(W_t,B_t)-\bar g(W_t)\bigr\|^{2}}
% \right).
% \]
% Both bounds reflect the same stability mechanism: generalization improves with
% larger injected noise $\sigma_t$ and degrades with larger step sizes $\eta_t$ and
% more variable stochastic gradients.
% They differ in form for two main reasons.
% First, our result applies to an arbitrary post-processed output
% $\widehat W=F(W_{1:T})$, whereas \citet{algo1} focus on the final iterate.
% Second, our heavy-tail analysis controls stepwise $f_\theta$-divergence increments
% and then upper bounds them using R\'enyi divergence; unlike KL, R\'enyi
% divergences do not admit the same additive chain rule along the optimization
% path, which naturally leads to an additive per-step accumulation before applying
% outer norm inequalities.

% In the sub-Gaussian regime, the two forms can be related by standard inequalities.
% Let $X_t \coloneqq g(W_t,B_t)-\bar g(W_t)$. Then Cauchy--Schwarz and Jensen yield
% \[
% \frac{1}{\sqrt{T}}\sum_{t=1}^{T-1}\frac{|\eta_t|}{\sigma_t}\,\E\|X_t\|
% \;\le\;
% \sqrt{\sum_{t=1}^{T-1}\frac{\eta_t^2}{\sigma_t^2}\,(\E\|X_t\|)^2}
% \;\le\;
% \sqrt{\sum_{t=1}^{T-1}\frac{\eta_t^2}{\sigma_t^2}\,\E\|X_t\|^{2}}.
% \]
% Thus, when $\theta=2$ our bound is comparable to the KL path bound up to constants,
% while our formulation continues to apply in genuinely heavy-tailed settings where
% MGF-based KL techniques can become ineffective or inapplicable.
% Simulation studies for linear regression are given in Appendix~\ref{ap:SGDL}.



%%%%%%%%%%%%%%%%%%%%%%%%%%%%%%%%%%%%%%%%%%%%%%%%
%%%%%%%%%%%%%%%%%%Section F%%%%%%%%%%%%%%%%%%%%%
%%%%%%%%%%%%%%%%%%%%%%%%%%%%%%%%%%%%%%%%%%%%%%%%
% \section{Calculation Details for Examples and Simulations}\label{Tighterexampleappen}

% \subsection{Example of quadratic mean estimation}
% Let the hypothesis class be $\mathcal{W}=[-1,1]$, and consider the squared loss
% \[
% \ell(w,Y)=(w-Y)^2,
% \]
% where $Y$ is a symmetric Weibull with
% $\mathbb{P}(|Y|\ge t)=\tfrac12 e^{-t^\theta}$ for $\theta \in (0,1)$. In this example, $f(w,x)=w$ and there is no $X$. Let $S=\{Y_i\}_{i=1}^n$ and take the empirical risk minimizer
% \[
% W(S):=\arg\min_{w\in[-1,1]} L_S(w)
% \;=\; \clip\!\left(\frac1n\sum_{i=1}^n Y_i,\;[-1,1]\right),
% \]
% which is a deterministic and continuous function of $S$. Consequently,
% $P_{W\mid S}$ is a point mass and is singular with respect to the marginal $P_W$,
% so $I_\alpha(W;S)=\infty$ for every $\alpha>1$. Therefore, Theorem~\ref{gen} does not produce a meaningful bound.

% To apply the chaining bound (Theorem~\ref{genchain}), fix the dyadic partitions
% \[
% \mathcal{P}_k
% =\Bigl\{\,[m2^{-(k-1)},(m+1)2^{-(k-1)})\cap[-1,1]\;:\; m\in\mathbb{Z}\Bigr\},
% \qquad k\ge1,
% \]
% which form an increasing partition sequence of $T$. 

% Denote $\mathcal{P}_k=\cup_{i\in \mathbb{Z}}\mathcal{P}_{k,i}$ denote the $i$-th cell in $k$-partition. With the metric from the
% square-loss example (see Example \ref{ex:squareloss}), we have $M=1$ and $L=1$ with
% $$
% e(\mathcal{W})=d(0,1)= C_2(\theta) D_\theta\bigl(2+2\|Y\|_{\psi_\theta}\bigr),
% $$
% and the corresponding Dudley integral $\gamma<\infty$ can be verified directly.

% Let $W_k=[W]_k$ denote the cell in $\mathcal{P}_k$ that contains $W$. Below are the rigorous, step-by-step derivation for the three claims to obtain our desired result.

% \medskip
% \noindent\textbf{Claim 1: }
% $
% I_\alpha(W_k;S) = \frac{1}{\alpha-1}\log \int p(s)\Bigl(\frac{1}{\mathbb{P}(W_k\in\mathcal{P}_{k,i(s)})}\Bigr)^{\alpha-1}ds.
% $
% \begin{proof}
% By definition, 
% \[
% I_\alpha(W_k; S) = D_\alpha(P_{W_k, S} \| P_{W_k} \otimes P_S) = \frac{1}{\alpha-1} \log \int_{\mathcal{S}} \sum_{i} \frac{\bigl( p_{W_k, S}(i, s) \bigr)^\alpha}{\bigl( P_{W_k}(i) p(s) \bigr)^{\alpha-1}} \, ds
% \]
% where $p_{W_k, S}(i, s)$ is the joint density-mass function, $p(s)$ is the marginal density of $S$, and $P_{W_k}(i) = \mathbb{P}(W_k \in \mathcal{P}_{k,i})$ is the marginal probability mass of $W_k$.

% The joint distribution is
% $p_{W_k, S}(i, s) = \mathbb{P}(W_k\in\mathcal{P}_{k,i}\mid S = s) p(s).$ Because $W$ is a deterministic function of $S$, $W_k$ (the cell containing $W$) is also a deterministic function of $S$. Let $i(s)$ be the unique index of the cell containing $W(s)$. Then,
% \[
% \mathbb{P}(W_k\in\mathcal{P}_{k,i} \mid S = s) = I\{i = i(s)\} = \begin{cases} 1 & \text{if } i = i(s) \\ 0 & \text{otherwise} \end{cases}
% \]
% Substitute this back into the R\'enyi mutual information formula:
% \begin{align*}
%     I_\alpha(W_k; S) &= \frac{1}{\alpha-1} \log \int_{\mathcal{S}} \sum_{i} \frac{\bigl( I\{i = i(s)\} p(s) \bigr)^\alpha}{\bigl( P_{W_k}(i) p(s) \bigr)^{\alpha-1}}ds = \frac{1}{\alpha-1} \log \int_{\mathcal{S}} \frac{p(s)^\alpha}{\bigl( P_{W_k}(i(s)) p(s) \bigr)^{\alpha-1}}ds\\
%     &\le  \frac{1}{\alpha-1} \log \int_{\mathcal{S}} p(s) \left( \frac{1}{\mathbb{P}(W_k \in \mathcal{P}_{k,i(s)})} \right)^{\alpha-1} ds.
% \end{align*}
% Since $I\{i = i(s)\}^\alpha =I\{i = i(s)\}$ for $\alpha > 1$, all terms in the sum over $i$ evaluate to $0$ except the single term where $i = i(s)$. 
% \end{proof}

% \medskip
% \noindent\textbf{Claim 2: } $\mathbb{P}(W_k\in\mathcal{P}_{k,i}) \;\ge\; C_{n,\theta}\,2^{-(k-1)}$.
% \begin{proof}
% Recall the deterministic mapping: $W = \text{clip}(\bar{Y}, [-1, 1])$, where $\bar{Y} = \frac{1}{n}\sum_{j=1}^n Y_j$.
% The event $\{W_k \in \mathcal{P}_{k,i}\}$ means that the clipped mean $W$ falls into the $i$-th cell $\mathcal{P}_{k,i} \subset[-1, 1]$. 

% Based on the definition $\mathcal{P}_{k,i} =[i 2^{-(k-1)}, (i+1)2^{-(k-1)}) \cap[-1, 1]$, for any internal cell strictly within $(-1, 1)$, we have $|\mathcal{P}_{k,i}| = 2^{-(k-1)}$.

% \begin{itemize}
%     \item For an internal cell $\mathcal{P}_{k,i} \subset (-1, 1)$, $W \in \mathcal{P}_{k,i}$ if and only if $\bar{Y} \in \mathcal{P}_{k,i}$ (no clipping).
%     \item For a boundary cell containing $-1$ or $1$, $W \in \mathcal{P}_{k,i}$ if $\bar{Y} \in \mathcal{P}_{k,i}$ or $\bar{Y} > 1$.
% \end{itemize}
% Therefore, for \textit{any} cell $\mathcal{P}_{k,i}$, 
% $\mathbb{P}(W_k \in \mathcal{P}_{k,i}) = \mathbb{P}(W \in \mathcal{P}_{k,i}) \ge \mathbb{P}(\bar{Y} \in \mathcal{P}_{k,i})$. If $f_{\bar Y}$ denotes the density of $\bar Y=n^{-1}\sum_{i=1}^n Y_i$ and
% $C_{n,\theta}:=\inf_{|y|\le 1} f_{\bar Y}(y)>0$. Thus, evaluating the probability by integrating the density over the cell gives
% \[
% \mathbb{P}(W_k \in \mathcal{P}_{k,i}) \ge\mathbb{P}(\bar{Y} \in \mathcal{P}_{k,i}) = \int_{\mathcal{P}_{k,i}} p_{\bar{Y}}(y) \, dy \ge \int_{\mathcal{P}_{k,i}} C_{n,\theta} \, dy = C_{n,\theta} |\mathcal{P}_{k,i}| = C_{n,\theta} 2^{-(k-1)}.
% \]
% \end{proof}

% From Claims 1 and 2, we established that:
% \begin{align*}
% I_\alpha(W_k;S)& = \frac{1}{\alpha-1}\log \int p(s)\left(\frac{1}{\mathbb{P}(W_k\in\mathcal{P}_{k,i(s)})}\right)^{\alpha-1}ds \le \frac{1}{\alpha-1}\log \int p(s)\left(\frac{1}{C_{n,\theta} 2^{-(k-1)}}\right)^{\alpha-1}ds\\
% & = \frac{1}{\alpha-1}\log \left[ \left(\frac{1}{C_{n,\theta} 2^{-(k-1)}}\right)^{\alpha-1} \int_{\mathcal{S}} p(s) \, ds \right]= \frac{1}{\alpha-1}\log \left[ \left(\frac{1}{C_{n,\theta} 2^{-(k-1)}}\right)^{\alpha-1} \right]\\
% &= (k-1)\log 2 - \log C_{n,\theta}.
% \end{align*}
% since $\mathbb{P}(W_k \in \mathcal{P}_{k,i(s)}) \ge C_{n,\theta} 2^{-(k-1)}$. Substituting the upper bound into Theorem~\ref{genchain} yields a finite multiscale complexity term,
% $$
% \mathbb E \bigl[\operatorname{gen}(W,S)\bigr]
% \le
% \frac{M_\theta C_2(\theta) D_\theta\bigl(2+2\|Y\|_{\psi_\theta}\bigr)}{\sqrt n}
% \sum_{k=1}^\infty 2^{-(k-1)}
% \bigl(2^{1/\theta}[(k-1)\log 2 - \log C_{n,\theta}]+4\bigr).
% $$

% \subsection{Four approaches to upper-bound $\E[X_W]$}
% We build on the example of \citet{asadi2018chaining} to illustrate that, in heavy-tailed
% regimes, chaining mutual information can yield substantially tighter bounds than
% single-scale information or classical chaining (Dudley's entropy bound).

% Let $S = \{Z_1,Z_2\}$ be i.i.d. Weibull random variables with $\mathbb{P}(Z_i\ge t)=\exp(-t^\theta),~t\ge 0,\ \theta>0.$
% Given an index set $T\subset \mathbb{R}^2$, define the canonical Weibull process indexed by $T\subset \{w\in \RR^2: \|w\|_2=1\}$: 
% \[
% X_w := \sum_{i=1}^2 w_i Z_i,\qquad w\in T~~\text{and}~~\|w\|_2 = 1. 
% \]
% Reparameterizing $w$ by its phase $\phi\in[0,2\pi)$, we have $w(\phi)=(\cos\phi,\sin\phi)$, and thus $X_\phi := X_{w(\phi)} = Z_1\cos\phi + Z_2\sin\phi$. Let 
% $$l(\phi,S):=-X_\phi.$$
% \paragraph{Verifying that $\{-X_\phi\}_{\phi\in[0,2\pi)}$ is a sub-Weibull process (Definition~\ref{subWp}).}
% For any $\phi_1,\phi_2\in[0,2\pi)$,
% \[
% X_{\phi_1}-X_{\phi_2}
% = (w(\phi_1)-w(\phi_2))^\top Z
% = (\cos\phi_1-\cos\phi_2)Z_1 + (\sin\phi_1-\sin\phi_2)Z_2.
% \]
% By the weak triangle inequality for $\psi_\theta$-Orlicz norms (see \eqref{eq:wt} in Appendix),
% \begin{align*}
% \|X_{\phi_1}-X_{\phi_2}\|_{\psi_\theta}
% &\le
% D_\theta\Bigl(|\cos\phi_1-\cos\phi_2|\,\|Z_1\|_{\psi_\theta}
% +|\sin\phi_1-\sin\phi_2|\,\|Z_2\|_{\psi_\theta}\Bigr).
% \end{align*}
% Since $\|Z_1\|_{\psi_\theta}=\|Z_2\|_{\psi_\theta}$ (i.i.d.\ coordinates), we get
% \begin{align*}
% \|X_{\phi_1}-X_{\phi_2}\|_{\psi_\theta}
% &\le
% D_\theta\|Z_1\|_{\psi_\theta}\Bigl(|\cos\phi_1-\cos\phi_2|+|\sin\phi_1-\sin\phi_2|\Bigr)\\
% &\le
% \sqrt{2}\,D_\theta\,\|Z_1\|_{\psi_\theta}\,
% \bigl\|w(\phi_1)-w(\phi_2)\bigr\|_2,
% \end{align*}
% where we used $|a|+|b|\le \sqrt{2}\|(a,b)\|_2$.
% Moreover, on the unit circle we have the chord-length bound
% $
% \|w(\phi_1)-w(\phi_2)\|_2 \le |\phi_1-\phi_2|~(\text{chord length} \le \text{arc length})$. Therefore,
% \[
% \|X_{\phi_1}-X_{\phi_2}\|_{\psi_\theta}
% \le
% \sqrt{2}\,D_\theta\,\|Z_1\|_{\psi_\theta}\,|\phi_1-\phi_2|.
% \]
% Let $d(\phi_1,\phi_2):=|\phi_1-\phi_2|$, $\{-X_\phi\}$ is a $(\theta,C)$-sub-Weibull process with
% $C:=\sqrt{2}D_\theta\|Z_1\|_{\psi_\theta}$.


% In Subsection~\ref{sec:tightness}, we consider four approaches to upper-bound $\E[X_W]$.
% We begin with the MI method. By Theorem~\ref{infob}, controlling $\E[X_W]$
% requires bounding the $\alpha$-information $I_\alpha(W;S)$.
% In the present example, the random perturbation $\xi$ contains a singular (degenerate) component,
% which implies that $I_\alpha(W;S)=\infty$. Consequently, the MI-based bound is vacuous:
% it yields $\E[X_W]\le \infty$.

% The second approach is the classical chaining method. A crude bound is
% \begin{equation}\label{eq:method2_step1}
% \E[X_W]\le \E\Big[\sup_{\phi\in[0,2\pi)} X_\phi\Big].
% \end{equation}
% For $\varepsilon\in(0,2\pi]$, the circle can be covered by at most $\lceil 2\pi/\varepsilon\rceil$ arcs of length $2\varepsilon$ with radius $\varepsilon$,
% and each such arc has chordal diameter at most its arc length. Hence,
% \begin{equation}\label{eq:circle_cover}
% N(T,d,\varepsilon)=\left\lceil \frac{2\pi}{2\varepsilon}\right\rceil=\left\lceil \frac{\pi}{\varepsilon}\right\rceil,\qquad \varepsilon\in(0,\pi].
% \end{equation}
% For $\varepsilon\ge \pi$, clearly $\log N(T,d,\varepsilon)=0$. Then Theorem~\ref{Dudley} and \eqref{eq:circle_cover} yield
% \begin{align}\label{eq:Dudley_circle}
% \E\Big[\sup_{\phi\in[0,2\pi)} X_\phi\Big]
% &\le 4CK_\theta \int_0^\infty \big[\log N(T,d,\varepsilon)\big]^{1/\theta} d\varepsilon = 4CK_\theta \int_0^{2\pi} \big[\log \left\lceil \frac{\pi}{\varepsilon}\right\rceil\big]^{1/\theta} d\varepsilon.
% \end{align}
% Let $I_\theta =\int_{0}^{\pi}\Big[\log \Big\lceil \frac{\pi}{\varepsilon}\Big\rceil\Big]^{1/\theta}d\varepsilon$. On each interval where \(\left\lceil \frac{\pi}{\varepsilon}\right\rceil\) is constant, we can compute the integral exactly. For \(k\ge 2\),
% \[
% \Big\lceil \frac{\pi}{\varepsilon}\Big\rceil = k
% \quad \Longleftrightarrow \quad
% \varepsilon \in \Big[\frac{\pi}{k},\frac{\pi}{k-1}\Big).
% \]
% Therefore,
% \[
% I_\theta
% =\sum_{k=2}^\infty \int_{\pi/k}^{\pi/(k-1)} (\log k)^{1/\theta},d\varepsilon
% =\sum_{k=2}^\infty (\log k)^{1/\theta}\Big(\frac{\pi}{k-1}-\frac{\pi}{k}\Big)=\pi \sum_{k=2}^\infty \frac{(\log k)^{1/\theta}}{k(k-1)}.
% \]
% The sum converges for every \(\theta>0\), since \((\log k)^{1/\theta}/(k(k-1)) \asymp (\log k)^{1/\theta}/k^2\). Hence,
% \begin{equation}\label{eq:method2_final}
% \E[X_W]\le \E\Big[\sup_{\phi}X_\phi\Big]
% \le 4\pi C K_\theta \sum_{k=2}^\infty \frac{(\log k)^{1/\theta}}{k(k-1)}.
% \end{equation}
% This is a purely geometric bound via covering numbers, and it does not exploit any algorithmic dependence between $W$ and $S$.



% The third approach is the R\'enyi chained mutual-information bound. Since continuous mutual information can be ill-behaved, we quantize $W$ and apply a chaining argument.

% In the present example, $e(T)=\pi$. For each integer $k\ge 1$, let 
% $$\mathcal{P}_k:=\left\{\left[0, \frac{2 \pi}{2^{k-1}}\right),\left[\frac{2 \pi}{2^{k-1}}, 2 \times \frac{2 \pi}{2^{k-1}}\right), \ldots,\left[\left(2^{k-1}-1\right) \frac{2 \pi}{2^{k-1}}, 2 \pi\right)\right\}=\cup_{m=0}^{2^{k-1}-1}\mathcal{P}_{k,m}$$
% be the equal-length partition of $\mathcal{W}=[0,2\pi)$
% into $2^{k-1}$ cells. Here:
% $\mathcal{P}_{k,m}:= \left[\frac{2\pi}{2^{k-1}}m,  \frac{2\pi}{2^{k-1}}(m+1) \right)$. 

% For any $w \in [0, 2\pi)$ and any level $k$, one can find the cell index $m$ using the floor function:
%  $m_k = \left\lfloor \frac{w}{2\pi} \cdot 2^{k-1} \right\rfloor$. The set $[w]_k$ is given by 
%  $[w]_k=\left[m_k \frac{2 \pi}{2^{k-1}},(m_k+1) \frac{2 \pi}{2^{k-1}}\right)$. Then, 
% $$[W]_k=\left[\left\lfloor \frac{W}{2\pi} \cdot 2^{k-1} \right\rfloor \frac{2 \pi}{2^{k-1}},\left(\left\lfloor \frac{W}{2\pi} \cdot 2^{k-1} \right\rfloor+1\right) \frac{2 \pi}{2^{k-1}}\right)$$
% is a discrete random variable taking values in the partition set $\mathcal{P}_k$.



% Given $S$, the minimizer $\phi^\star(S)$ is deterministic. With probability $\varepsilon$, we have $\xi=0$ and thus $W=\phi^\star(S)$,
% so $[W]_k$ falls into the ``optimal'' cell. With probability $1-\varepsilon$, $\xi$ is uniform on $(-\pi,\pi)$. Let $\mathcal{P}_{k,{m^\star}} = [\phi^\star(S)]_k$ be the cell containing the minimizer. The random variable $[W]_k \in \mathcal{P}_k$ has the following conditional probability mass function:

% 1. For the ``optimal'' cell containing $\phi^\star(S)$:
%    $$ \mathbb{P}\Big([W]_k = \mathcal{P}_{k,{m^\star}} \;\Big|\; S \Big) = \varepsilon + \frac{1 - \varepsilon}{2^{k-1}}.$$

% 2. For any of the other $2^{k-1} - 1$ cells ($m \neq m^\star$):
%    % $$ \mathbb{P}\Big([W]_k = \mathcal{P}_{k,{m}} \;\Big|\; S \Big) = \frac{1 - \varepsilon}{2^{k-1}}.$$

% Using standard trigonometry, the objective can be written as
% $$ X_\phi = Z_1 \cos\phi + Z_2 \sin\phi = \|Z\|_2 \cos(\phi - \theta_Z).$$
% where $Z = (Z_1, Z_2)$ and $\theta_Z$ is the angle of the vector $Z$ in polar coordinates.

% To minimize $X_\phi$, the cosine term must equal $-1$. This happens when the angle $\phi$ points in the exact opposite direction of $\theta_Z$. Therefore, the deterministic minimizer is:
% $$ \phi^\star(S) = \theta_Z \oplus \pi \pmod{2\pi}.$$
% Note that $Z_1, Z_2$ are i.i.d. standard strictly positive Weibull variables, and $\theta_Z \in [0, \pi/2]$. The probability density function for $\theta_Z$ is
% \begin{equation*}
%      f_{\theta_Z}(\alpha) = \frac{\theta (\cos \alpha \sin \alpha)^{\theta-1}}{(\cos^\theta \alpha + \sin^\theta \alpha)^2}, \quad \text{for } \alpha \in \left[0, \frac{\pi}{2}\right] .
% \end{equation*}
% See the corresponding Lemma \ref{LEM:theta-Z} for the proof.

%  Using the order-$2$ R\'enyi mutual information $I_2$ (chosen for algebraic simplicity), we obtain the following.
% \[
% I_2([W]_k;S)
% =\log \E_S \left[\sum_{C \in \mathcal{P}_k} \frac{\mathbb{P}([W]_k = C \mid S)^2}{\mathbb{P}([W]_k = C)} \right]=\log \left(\sum_{m=0}^{2^{k-1}-1} \frac{u^2+q_m\left(\varepsilon^2+2 \varepsilon u\right)}{\varepsilon q_m+u}\right),~u=\frac{1-\varepsilon}{2^{k-1}}.
% \]
% where $q_m=\mathbb{P}\left(\theta_Z \in([\frac{2\pi}{2^{k-1}}m-\pi,\frac{2\pi}{2^{k-1}}(m+1)-\pi) \cap[0, \pi / 2])\right)$ (the true probability that the minimizer lies in cell $m$).
% In particular, for every fixed $k$, we have 
% \begin{center}
%     $I_2([W]_k;S)\to 0$ as $\varepsilon\to 0$.
% \end{center}
% Applying Theorem~\ref{chaining} with $e(T)=\pi$ yields the following.
% \[
% \E[X_W]
% \le
% C\,\pi\sum_{k=1}^{\infty}2^{-(k-2)}
% \Big(2^{1/\theta}\big(I_2([W]_k;S)+C_{\alpha,\theta}\big)^{\I}+2\Big).
% \]
% with $C:=\sqrt{2}D_\theta\|Z_1\|_{\psi_\theta}=2^{\frac{1}{\theta}+\frac{1}{2}}\|Z_1\|_{\psi_\theta}$ for $\theta \in (0,1)$.

% The fourth approach is a direct calculation of $\mathbb{E}[X_W]$. Given $S$,
% \[
% \min_{\phi\in[0,2\pi)} l(\phi,S)=-\|Z\|_2.
% \]
% Moreover, conditional on $S$, the algorithm output $W$ is a mixture distribution:
% \begin{itemize}
%     \item With probability $\varepsilon$: $\quad \xi = 0 \implies W = \phi^\star(S)$ and \(X_W=-\|Z\|_2\);
%         \item With probability $1-\varepsilon$: $\quad \xi \sim \text{Unif}(-\pi, \pi) \implies W \sim \text{Unif}(0, 2\pi)$.
% \end{itemize}
% Hence,
%   \[
%   \mathbb E[X_W\mid S,\xi\neq 0]
%   =\frac{1}{2\pi}\int_{0}^{2\pi}\big(Z_1\cos\phi+Z_2\sin\phi\big)d\phi
%   =0. \]
%   Then, by the law of total expectations,
%   \[
% \mathbb E[X_W\mid S]=\varepsilon(-\|Z\|_2)+(1-\varepsilon)\cdot 0=-\varepsilon\|Z\|_2,
% \qquad\Rightarrow\qquad
% \mathbb E[X_W]=-\varepsilon\mathbb E\|Z\|_2.
% \]
% Each \(Z_i\) has density \(f(z)=\theta z^{\theta-1}e^{-z^\theta}\) on \([0,\infty)\). Thus
% \[
% \mathbb{E}\|Z\|_2
% =\int_{0}^{\infty}\!\!\int_{0}^{\infty}
% \sqrt{x^2+y^2}\,\theta^2 x^{\theta-1}y^{\theta-1}e^{-x^\theta-y^\theta}\,dx\,dy.
% \]
% Using polar coordinates on the first quadrant $x=r\cos t$, $y=r\sin t$,
% $t\in[0,\pi/2]$, $r\in[0,\infty)$, and $dx\,dy=r\,dr\,dt$, one obtains
% \[
% \mathbb{E}\|Z\|_2
% =\theta\,\Gamma\!\left(2+\frac{1}{\theta}\right)\int_{0}^{\pi/2}
% \frac{(\cos t\,\sin t)^{\theta-1}}{(\cos^\theta t+\sin^\theta t)^{2+1/\theta}}\,dt.
% \]
% Equivalently, after the substitutions $x=\tan^\theta t$ and then $x=t/(1-t)$,
% \[
% \mathbb E[X_W]=-\varepsilon\mathbb E\|Z\|_2
% =-\varepsilon\Gamma\!\left(2+\frac{1}{\theta}\right)\int_{0}^{1}
% \Big(t^{2/\theta}+(1-t)^{2/\theta}\Big)^{1/2}\,dt.
% \]
% If $\theta=0.5$, then
% \[
% \mathbb{E}[X_W]
% =-\varepsilon\,\Gamma(4)\int_{0}^{1}\Big(t^{4}+(1-t)^{4}\Big)^{1/2}\,dt
% =-6\varepsilon\int_{0}^{1}\sqrt{t^{4}+(1-t)^{4}}\,dt\approx -3.6\varepsilon.
% \]
% Since $\int_{0}^{1}\sqrt{t^{4}+(1-t)^{4}}\,dt \approx 0.60$.




%%%%%%%%%%%%%%%%%%%%%%%%%%%%%%%%%%%%%%%%%%%%%%%%%%
%%%%%%%%%%%%%%%%%%%Section G%%%%%%%%%%%%%%%%%%%%%
%%%%%%%%%%%%%%%%%%%%%%%%%%%%%%%%%%%%%%%%%%%%%%%%%%

\section{Proofs of Useful Lemmas and Formulas}\label{app:lemmas}

This section presents several lemmas used in the proofs of the main results. We begin by showing that the shifted-log function is convex.

\begin{lemma}\label{lemma:convex}
For $\theta>0$, the function $x \mapsto x\log^\theta(x + A),~A \geq 1$ is convex over $x>0$.
\end{lemma}

\begin{proof}
Let $L = \log(x + A)$, and define $f(x) := x\, L^{\theta}$.
The first derivative is
\$
f'(x) = L^{\theta} + x \theta L^{\theta - 1} \frac{1}{x + A}
= L^{\theta - 1} \left( L + \frac{\theta x}{x + A} \right)=L^{\theta - 1} \left( L + \theta\left[1-\frac{A}{x + A}\right] \right),
\$
while the second derivative is 
\begin{align*}
f''(x)
    &= (\theta - 1)L^{\theta - 2}\frac{1}{x + A}\left( L + \frac{\theta x}{x + A} \right)
+ L^{\theta - 1}\left(\frac{1}{x+A} + \frac{\theta A}{(x + A)^2} \right)\\
    &=\theta\frac{(\log(x+A))^{\theta-2}}{(x+A)^2}\Big[(x+2A)\log(x+A)+(\theta-1)x\Big].
\end{align*}
Since $\theta>0$, $(x+A)^2>0$, and for $x>0, A\ge 1$, we have $\log(x+A)\ge 0$, so $\log^{\theta-2}(x+A)>0$ whenever it is defined. It remains to prove that the bracketed term is positive. Define
$$B(x)=(x+2A)\log(x+A)+(\theta-1)x.$$
The worst case occurs as $\theta\to0^+$, so it suffices to show
$$(x+2A)\log(x+A)-x>0\quad\text{for all}~x>0.$$
Define $h(x)=(x+2A)\log(x+A)-x$. Then $h(0)=2A\log A\ge 0$ (since $A\ge1$), and
$$h^{\prime}(x)=\log(x+A)+\frac{x+2A}{x+A}-1=\log(x+A)+\frac A{x+A}>0\quad\mathrm{for~}x>0.$$

Thus, $h(x)$ is strictly increasing and positive for all $x>0$. Hence, $B(x)>0$ for all $x>0$, $A\ge 1$, and $\theta > 0$. Therefore, $f''(x)>0$ on $(0,\infty)$, proving convexity.
\end{proof}


\begin{lemma}\label{lemma:concave}
Let $\theta \ge 1$ and $c>0$. The function $x \mapsto \log^\theta(x+c)$ is concave on $[0,\infty)$ whenever
$c \ge \exp(\theta-1)$.
\end{lemma}

\begin{proof}
Define $f(x):=\log^\theta(x+c)$ for $x\ge 0$. Then
\[
f'(x)=\theta \bigl(\log(x+c)\bigr)^{\theta-1}\frac{1}{x+c}.
\]
Differentiating it again yields
\[
\begin{aligned}
f''(x)
&=\theta\left[
(\theta-1)\bigl(\log(x+c)\bigr)^{\theta-2}\frac{1}{(x+c)^2}
-\bigl(\log(x+c)\bigr)^{\theta-1}\frac{1}{(x+c)^2}
\right]\\
&=\frac{\theta}{(x+c)^2}\bigl(\log(x+c)\bigr)^{\theta-2}\Bigl[(\theta-1)-\log(x+c)\Bigr].
\end{aligned}
\]
Since $\theta\ge 1$ and $x+c>0$, the prefactor $\frac{\theta}{(x+c)^2}\bigl(\log(x+c)\bigr)^{\theta-2}$ is nonnegative whenever $\log(x+c)\ge 0$, and thus $f''(x)\le 0$ holds provided
\[
\log(x+c)\;\ge\;\theta-1\qquad \text{for all }x\ge 0.
\]
The left-hand side is minimized at $x=0$, so it suffices that $\log c \ge \theta-1$, i.e., $c\ge e^{\theta-1}$.
\end{proof}



To construct generalization bounds for heavy-tailed sub-Weibull, the following lemmas are needed, similar to the decorrelation lemma of \cite{chu2023unified}.

\begin{lemma}[A new Young-type inequality of non-convex function]\label{young}
For $x, y \ge 0$, we have
\[
xy \leq 2^{\frac{1}{\theta}} x [\log(x + A)]^{\frac{1}{\theta}} + \exp(y^\theta),
\]
where $A\ge (2^{\lceil \frac{2}{\theta} \rceil-2} \lceil \frac{2}{\theta} \rceil !)^2 \vee 1$ is a positive constant depending on $\theta>0$.
\end{lemma}

\begin{proof}
We discuss two cases below.\\ 

\textbf{Case 1.} If $y \leq 2^{\frac{1}{\theta}} \log(x + A)^{\frac{1}{\theta}}$, the inequality is immediate for $A\ge 1$ since $\exp(y^\theta)>0$.

\textbf{Case 2.} If $y > 2^{\frac{1}{\theta}} \log(x+A)^{\frac{1}{\theta}}$, we have
\[
xy < \left(\exp(\frac{y^\theta}{2}) - A \right) y \leq \left(\exp(\frac{y^\theta}{2}) - \sqrt{A} \right)\left(\exp(\frac{y^\theta}{2}) + \sqrt{A}\right) \leq \exp(y^\theta).
\]
The first inequality is because in this case $x<\exp(\frac{y^\theta}{2})-A$ and we can properly choose $A$ satisfying $y\leq \exp(\frac{y^\theta}{2})+\sqrt{A}$ to make the second inequality hold.

Let \(m=\lceil 2/\theta\rceil\). What does hold for all \(y\ge 0\) and \(\theta>0\) is
\[
\exp\left(\frac{y^\theta}{2}\right)\ \ge\ \frac{y^2}{2^mm!}.
\]
The claim is trivial for $y=0$, so assume $y>0$ and set $x:=y^\theta/2\ge 0$.
By the Taylor expansion of $e^x$ with nonnegative terms, for any integer $m\ge 0$,
\[
e^x=\sum_{k=0}^{\infty}\frac{x^k}{k!}\ \ge\ \frac{x^m}{m!}.
\]
Hence
\[
\exp\!\left(\frac{y^\theta}{2}\right)=e^x
\ \ge\ \frac{1}{m!}\left(\frac{y^\theta}{2}\right)^m
\ =\ \frac{y^{\theta m}}{2^m\,m!}.
\]
We now compare $y^{\theta m}$ and $y^2$.

\textbf{Case 1: $0\le y\le 1$.}
We have $\exp(y^\theta/2)\ge 1$. Also $y^2\le 1$, so
\[
\frac{y^2}{2^m m!}\le \frac{1}{2^m m!}\le 1 \le \exp\!\left(\frac{y^\theta}{2}\right).
\]

\textbf{Case 2: $y\ge 1$.}
Since $m=\left\lceil \frac{2}{\theta}\right\rceil$, we have $\theta m\ge 2$.
Therefore $y^{\theta m}\ge y^2$ for all $y\ge 1$, thus
\[
\exp\!\left(\frac{y^\theta}{2}\right)\ \ge\ \frac{y^{\theta m}}{2^m m!}
\ \ge\ \frac{y^2}{2^m m!}.
\]
So we have
\[
y-\exp\left(\frac{y^\theta}{2}\right) \le y-\frac{y^2}{2^mm!}\le 2^{m-2}m!,~y \ge 0.
\]
$A\ge 2^{m-2}m! \vee 1= (2^{\lceil \frac{2}{\theta} \rceil-2} \lceil \frac{2}{\theta} \rceil !)\vee 1$ is enough.

\end{proof}




% \begin{lemma}\label{LEM:theta-Z}
% Let $S = \{Z_1,Z_2\}$ be i.i.d. Weibull random variables with $\mathbb{P}(Z_i>t)=\exp(-t^\theta),~t\ge 0,\ \theta>0$.  Put $Z = (Z_1, Z_2)\in \mathbb{R}^+\times\mathbb{R}^+$ with its length $\|Z\|_2$. Let $\theta_Z$ be the angle of the vector $Z$ in polar coordinates. Then, the density function for $\theta_Z$ is
% \begin{equation*}
%      f_{\theta_Z}(\alpha) = \frac{\theta (\cos \alpha \sin \alpha)^{\theta-1}}{(\cos^\theta \alpha + \sin^\theta \alpha)^2}, \quad \text{for } \alpha \in \left[0, \frac{\pi}{2}\right].
% \end{equation*}
%  \end{lemma}

% \begin{proof}
% We are given \(S = \{Z_1, Z_2\}\) i.i.d. Weibull random variables with survival function
% \[
% \mathbb{P}(Z_i \ge t) = e^{-t^\theta},~t \ge 0,\ \theta > 0.
% \]
% with probability density function $f_{Z_i}(z) = \theta z^{\theta-1} e^{-z^\theta},\quad z \ge 0$. The joint density of \((Z_1, Z_2)\) is
% \[
% f_{Z_1,Z_2}(z_1,z_2) = \theta^2 (z_1 z_2)^{\theta-1} e^{-(z_1^\theta + z_2^\theta)},\quad z_1,z_2 \ge 0.
% \]

% Let \(R =\|Z\|_2= \sqrt{Z_1^2+Z_2^2}\) and \(\Theta = \theta_Z\) be the polar angle, i.e.
% \[
% Z_1 = R\cos\Theta,\quad Z_2 = R\sin\Theta,\qquad R>0,\; \Theta\in[0,\tfrac{\pi}{2}].
% \]
% The Jacobian of the transformation \((z_1,z_2) \mapsto (r,\varphi)\) is \(r\), so the joint density \((R,\Theta)\)
% \[
% f_{R,\Theta}(r,\varphi) = f_{Z_1,Z_2}(r\cos\varphi, r\sin\varphi) r
% = \theta^2 (r^2\cos\varphi\sin\varphi)^{\theta-1} e^{-r^\theta(\cos^\theta\varphi+\sin^\theta\varphi)}r.
% \]
% Simplifying,
% \[
% f_{R,\Theta}(r,\varphi) = \theta^2 r^{2\theta-1} (\cos\varphi\sin\varphi)^{\theta-1} e^{-r^\theta(\cos^\theta\varphi+\sin^\theta\varphi)},\qquad r>0,\; \varphi\in [0,\tfrac{\pi}{2}].
% \]

% Then, the marginal density of \(\Theta\) is obtained by integrating out \(r\):
% \[
% f_\Theta(\varphi) = \int_0^\infty \theta^2 r^{2\theta-1} (\cos\varphi\sin\varphi)^{\theta-1} e^{-r^\theta(\cos^\theta\varphi+\sin^\theta\varphi)}\,dr.
% \]
% Make the substitution
% \[
% u = r^\theta\bigl(\cos^\theta\varphi+\sin^\theta\varphi\bigr),\quad
% r = \left(\frac{u}{\cos^\theta\varphi+\sin^\theta\varphi}\right)^{1/\theta},\quad
% dr = \frac{1}{\theta}\,u^{1/\theta-1}\bigl(\cos^\theta\varphi+\sin^\theta\varphi\bigr)^{-1/\theta}\,du.
% \]
% Then
% \[
% r^{2\theta-1} = \left(\frac{u}{\cos^\theta\varphi+\sin^\theta\varphi}\right)^{(2\theta-1)/\theta}
% = u^{(2\theta-1)/\theta}\bigl(\cos^\theta\varphi+\sin^\theta\varphi\bigr)^{-(2\theta-1)/\theta}.
% \]
% Multiplying \(dr\) and \(r^{2\theta-1}\) gives
% \[
% dr\; r^{2\theta-1} = \frac{1}{\theta}\,u^{(1/\theta-1)+(2\theta-1)/\theta}\,
% \bigl(\cos^\theta\varphi+\sin^\theta\varphi\bigr)^{-1/\theta-(2\theta-1)/\theta}\,du= \frac{1}{\theta}\,u\,\bigl(\cos^\theta\varphi+\sin^\theta\varphi\bigr)^{-2}\,du.
% \]
% Substituting into the integral, for $\varphi\in [0,\tfrac{\pi}{2}]$.
% \begin{align*}
%     f_\Theta(\varphi) &= \theta^2 (\cos\varphi\sin\varphi)^{\theta-1}
% \int_0^\infty \frac{1}{\theta}\,u\,e^{-u}\,\bigl(\cos^\theta\varphi+\sin^\theta\varphi\bigr)^{-2}\,du\\
% &= \theta (\cos\varphi\sin\varphi)^{\theta-1} \bigl(\cos^\theta\varphi+\sin^\theta\varphi\bigr)^{-2}
% \int_0^\infty u e^{-u}du=\frac{\theta(\cos\varphi\sin\varphi)^{\theta-1}}{(\cos^\theta\varphi+\sin^\theta\varphi)^2}.
% \end{align*}
% \end{proof}







%%%%%%%%%%%%%%%%%%%%%%%%%%%%%%%%%%%%%%%%%%%%%%%%%%
%%%%%%%%%%%%%%%%%%%Section H%%%%%%%%%%%%%%%%%%%%%
%%%%%%%%%%%%%%%%%%%%%%%%%%%%%%%%%%%%%%%%%%%%%%%%%%
\section{Additional Numerical Experiments}\label{app:add_experiments}
This section provides empirical evidence for the main experimental claims. We first verify that the order-preserving power transformation used in Section~\ref{sec:rlhf-experiments} indeed produces a heavy-tailed reward. We then study a distinct heavy-tailed reward construction based on injecting centered Weibull noise and show that the qualitative behavior observed under R\'enyi regularization persists.

\subsection{RLHF experimental details}\label{app:rlhf_details}
We study the effect of R\'enyi order $\alpha$ in~\eqref{optRENYI_penalty} while fixing $\beta=1$ throughout. We sweep $\alpha\in\{1,2,\dots,10\}$, where $\alpha\to 1$ corresponds to the KL regularization. We train in the \texttt{hh-rlhf} training split (50{,}000 prompts) for up to $1500$ optimization steps and evaluate every $10$ steps on 128 prompts that are not used for training. For text generation, we cap the prompt length at 256 tokens and the completion length at 128 tokens. We sample with temperature 1.0 and top-$p$ 0.95, that is, at each decoding step we sample from the smallest set of tokens whose cumulative probability under the model is at least 0.95 (nucleus sampling). During training, we sample 4 completions for each prompt so that GRPO can compare multiple responses to the same prompt and form a group-relative reward signal for the policy update. We use a per-device batch size of 1 and accumulate gradients over 8 steps, which gives an effective batch size of 8, together with a learning rate of $10^{-5}$ and BF16 precision.\footnote{``BF16'' precision refers to the 16-bit \emph{bfloat16} floating-point format, which reduces memory usage and can improve training/inference throughput while typically maintaining numerical stability comparable to 32-bit floats.} During evaluation, we generate with batch size 16 and estimate R\'enyi divergence on a 16-prompt subset to reduce computation. For reward scoring, both reward models use a maximum input length of 512 tokens, with longer inputs truncated. We also recalibrate the proxy reward using mean- and scale statistics estimated from 256 prompts sampled from the reference policy so that reward values are on a comparable scale across runs.\\

\subsection{Tail behavior under the order-preserving power transformation}\label{sec:tail-behavior-power-transform}

As discussed in Section~\ref{sec:rlhf-experiments}, the RLHF experiments are conducted in a heavy-tailed reward regime induced by
an order-preserving power transform. In Figure~\ref{gold_proxy_trans}, we
complement that visual evidence with a quantitative diagnostic via the empirical MGF.

For the original or transformed reward samples $\{r_i\}$, we
compute
$
M(t)=\frac{1}{n}\sum_{i=1}^n e^{tr_i},~t\in\mathbb{R}.
$
Its growth for moderate values of $|t|$ provides a useful finite-sample summary
of tail amplification. $M(t)$ changes much more rapidly and explore as $t$ moves away from zero.

Figure~\ref{gold_proxy_MGF} shows that the transformed rewards exhibit
orders-of-magnitude larger empirical MGFs than the original rewards. For
$\beta=1$ and $\alpha=1$, at $t=-0.4$, the empirical MGF of the transformed
reward is approximately $5.77\times 10^{6}$, about $3.2\times 10^{6}$ times
larger than that of the original reward ($1.80$). For $\alpha=2$, the contrast
is even stronger: at $t=-0.4$, the ratio between the transformed and original
MGFs is approximately $8.34\times 10^{13}$. The logarithmic growth rate
$d\log M(t)/dt$ shows the same pattern. At $t=-0.4$, this quantity is
$-47.5$ for $\alpha=1$ and $-91.3$ for $\alpha=2$ after
transformation, compared with only $-1.47$ and $-1.57$ for the corresponding
original rewards. Taken together, these diagnostics reinforce the main-text
claim in Section~\ref{sec:rlhf-experiments}: the order-preserving power
transformation substantially amplifies tail behavior and thereby produces the
heavy-tailed reward regime used to study the robustness of
R\'enyi-regularized RLHF.


\begin{figure}[H]
\centering
\includegraphics[scale=0.225]{figure/mgf_comparison.pdf}
\caption{\small Empirical MGFs comparison for the original and transformed reward. Left: $\beta=1$ with $\alpha=1$ (KL). Right: $\beta=1$ with
$\alpha=2$ (R\'enyi). Blue curves correspond to the original reward $r(x,y)$,
and red curves correspond to the transformed reward
$\tilde r(x,y)=\operatorname{sign}(r(x,y))\,|r(x,y)|^{8}$. The much faster
growth of the empirical MGF for the transformed rewards is consistent with stronger tail amplification.}
\label{gold_proxy_MGF}
\end{figure}


%%%%%%%%%%%%%%%%%%%%%%%%%%%%%%%%%%%%%%%%%%%%%%%%%%%%%%%%%%%%%%%%%%%%%%%%%%%%%%%%%%%%%%%%%%%
%%%%%Alternative heavy-tailed reward construction by injecting centered Weibull noise%%%%%%
%%%%%%%%%%%%%%%%%%%%%%%%%%%%%%%%%%%%%%%%%%%%%%%%%%%%%%%%%%%%%%%%%%%%%%%%%%%%%%%%%%%%%%%%%%%
\subsection{Alternative heavy-tailed reward construction by injecting Weibull noise}\label{app:weibull-noise}
 

In addition to the power transform used in Section~\ref{sec:rlhf-experiments}, we consider a second heavy-tailed reward construction obtained by adding centered Weibull noise to the reward. Let
\[
r_0(x,y)=\langle \theta^{*},\psi(x,y)\rangle
\]
denote the clean reward, where $\psi(x,y)$ is the penultimate-layer representation of the same proxy reward model used in Section~\ref{sec:rlhf-experiments}, namely Qwen2.5-0.5B, and $\theta^{*}$ is the parameter vector of its trained final linear layer. We define the observed reward by
\[
r(x,y)=r_0(x,y)+\varepsilon(x,y),
\qquad
\varepsilon(x,y)=Z-\mathbb{E}[Z],
\qquad
Z\sim \mathrm{Weibull}(k,\eta(x,y)).
\]
Since $\mathbb{E}[Z]=\eta(x,y)\Gamma(1+1/k)$, the perturbation is centered and therefore preserves the conditional mean reward. We set
\(
k\in\{0.2,0.3,0.4,0.5\},
\)
with smaller $k$ corresponding to heavier-tailed noise, and take $\eta(x,y)$ to be the estimated standard error of the reward. We use the same penalized objective~\eqref{optRENYI_penalty}, fix $\beta=1$, and vary the R\'enyi order over
\(
\alpha\in\{2,4,6,8,10\}.
\)
All other implementation details are the same as those in Section~\ref{sec:rlhf-experiments}.

\begin{figure}[h]
\centering 
% \begin{minipage}[t]{0.48\textwidth}
% \centering
% \includegraphics[scale=0.26]{figure_Linear_model_LLM/divergence_vs_reward_kappa0p2.pdf}
% \end{minipage}
% \begin{minipage}[t]{0.48\textwidth}
% \centering
% \includegraphics[scale=0.26]{figure_Linear_model_LLM/divergence_vs_reward_kappa0p3.pdf}
% \end{minipage}
% \begin{minipage}[t]{0.48\textwidth}
% \centering
% \includegraphics[scale=0.26]{figure_Linear_model_LLM/divergence_vs_reward_kappa0p4.pdf}
% \end{minipage}
% \begin{minipage}[t]{0.48\textwidth}
% \centering
% \includegraphics[scale=0.26]{figure_Linear_model_LLM/divergence_vs_reward_kappa0p5.pdf}
% \end{minipage}
\includegraphics[scale=0.24]{Aligned_figures/figure9.pdf}
% \caption{Reward--divergence relationship under the centered Weibull-noise reward construction with $\beta=1$. The four subplots correspond to Weibull shape parameters 0.2, 0.3, 0.4, and 0.5, ordered from left to right and top to bottom. In each subplot, the horizontal axis reports the R\'enyi divergence and the vertical axis reports reward; the left panel shows gold reward and the right panel shows proxy reward. Colors indicate R\'enyi orders $\alpha\in\{2,4,6,8,10\}$.}
\caption{\small
Reward--divergence relationship under the centered Weibull-noise reward with shape parameters 0.2, 0.3, 0.4, and 0.5, ordered from left to right and top to bottom.
% In each subplot, the horizontal axis reports the R\'enyi divergence and the vertical axis reports reward; the left panel shows gold reward and the right panel shows proxy reward. Colors indicate R\'enyi orders $\alpha\in\{2,4,6,8,10\}$.
Across all Weibull shape parameters, the reward--divergence curves remain stable: the gold reward typically increases and then levels off, rather than exhibiting a sharp collapse, indicating that higher-order R\'enyi regularization provides a reliable trust-region with $\beta=1$.
}
\label{divergence_vs_reward_Weibull}
\end{figure}

Figures~\ref{divergence_vs_reward_Weibull} and~\ref{proxy_vs_gold_Weibull} show that the qualitative behavior observed in the main text persists under this perturbation model. Across all tested noise levels, the reward--divergence curves remain stable and the proxy--gold relationship remains monotone over the observed range. Although smaller $k$ produces heavier-tailed reward noise, it has the stabilizing effect under higher-order R\'enyi regularization. These results reinforce the main conclusion of Section~\ref{sec:rlhf-experiments}: the favorable empirical behavior of R\'enyi regularization is not an artifact of the power transformation, but persists under a distinct heavy-tailed reward construction.



\begin{figure}[H]
\centering 
% \begin{minipage}[t]{0.47\textwidth}
% \centering
% \includegraphics[scale=0.26]{figure_Linear_model_LLM/proxy_vs_gold_reward_kappa0p2.pdf}
% \end{minipage}
% \begin{minipage}[t]{0.47\textwidth}
% \centering
% \includegraphics[scale=0.26]{figure_Linear_model_LLM/proxy_vs_gold_reward_kappa0p3.pdf}
% \end{minipage}
% \begin{minipage}[t]{0.47\textwidth}
% \centering
% \includegraphics[scale=0.26]{figure_Linear_model_LLM/proxy_vs_gold_reward_kappa0p4.pdf}
% \end{minipage}
% \begin{minipage}[t]{0.47\textwidth}
% \centering
% \includegraphics[scale=0.26]{figure_Linear_model_LLM/proxy_vs_gold_reward_kappa0p5.pdf}
% \end{minipage}

\includegraphics[scale=0.45]{Aligned_figures/figure10.pdf}


% \caption{Proxy--proxy-gold reward relationship under the centered Weibull-noise reward construction with $\beta=1$. The four subplots correspond to Weibull shape parameters 0.2, 0.3, 0.4, and 0.5, ordered from left to right and top to bottom. In each subplot, the horizontal axis shows proxy reward and the vertical axis shows gold reward. Colors indicate R\'enyi orders $\alpha\in\{2,4,6,8,10\}$.}
\caption{\small
Proxy--proxy-gold reward relationship under the Weibull-noise reward construction with $\beta=1$ and shape parameters 0.2 to 0.5, ordered from left to right and top to bottom.
% In each subplot, the horizontal axis shows proxy reward and the vertical axis shows gold reward. Colors indicate R\'enyi orders $\alpha\in\{2,4,6,8,10\}$.
Across all noise levels, the proxy--gold curves remain monotone increasing over the observed range, showing that improvements in the proxy reward continue to translate consistently into improvements in the gold reward even under heavier-tailed Weibull.
}
\label{proxy_vs_gold_Weibull}
\end{figure}





\end{document}
